%% Intelligence Engineering -- the full course deck (one PDF).
%% Master file: preamble + title + framing, then \input the eleven chapter
%% fragments, ch01 to ch11.
%%
%% Build the full deck:  latexmk intelligence_engineering.tex
%%   (lualatex, two passes; see .latexmkrc)
%% Build one PDF per chapter:
%%   tools/build_slide_chapter_pdfs.sh <this dir>  ->  chapter_pdfs/chNN_*.pdf
%%
%% ALIGNED TO THE BOOK. The chapters are Huyen, Designing Machine Learning
%% Systems, one deck chapter per book chapter, in the book's order and under
%% the book's titles. The PDF sits in ../references/. The frame titles inside
%% a chapter are that chapter's own section and subsection headings.
%%
%% SPINE ONLY, WITH ONE EXCEPTION. Every frame in ch01 to ch11 is empty: the
%% structure is settled, the content is not written, and a frame carries its
%% title and nothing else. ch00_orientation.tex is the exception and is fully
%% written, because the opening lecture has to say what the course is before
%% anyone can decide to sit through it.
%%
%% ONE CHAPTER, ONE LECTURE, ONE FILE. There is no four-act grouping any
%% more; a chapter is a lecture. Each opens with its own \lecturepage divider,
%% which lives in the chapter file, and the theme's automatic section page is
%% switched off so a lecture gets one divider rather than two. Every chapter
%% also carries one spare \lecturepage mid-file as a segment-divider template.
%% A fragment carries no preamble of its own, which is what lets it build
%% alone as chapter_pdfs/chNN_*.pdf.
\documentclass[aspectratio=169]{beamer}
\usetheme{tuftedark}
% The backbone architecture diagram offsets arrow endpoints with calc syntax
% ($(node)+(dx,dy)$); the theme loads arrows.meta and positioning but not calc.
\usetikzlibrary{calc,fit,shapes.geometric}

% Drop slide images in images/; include by bare name, e.g.
%   \includegraphics[width=\linewidth]{backbone}
\graphicspath{{images/}}

% Placeholder image: \phimg[width-fraction]{label}. Draws a labeled box in place
% of an image not yet made; swap for \includegraphics once the asset exists.
\newcommand{\phimg}[2][0.5]{%
  \tikz\node[draw=tddim, thick, rounded corners, fill=tddim!12,
    minimum width={#1*\linewidth}, minimum height={0.62*#1*\linewidth},
    inner sep=4pt, align=center, text=tddim,
    font=\ttfamily\scriptsize]{\detokenize{#2}};%
}

% Placeholder side image: \phimgside[frac]{label}{caption}. The
% \sideimagecover hero at placeholder stage: a blank gray strip on the
% right, inset from the page edges so it clears the rotated watermark and
% the \slidesources marker (9mm right) and the footer band (16mm bottom).
% The label is NOT drawn; it stays in the call as the record of the
% filename to drop into images/. The caption runs up the strip's bottom
% left edge in the placeholder's own typewriter style, mirroring
% \slidesources on the right page edge. It sets the theme's strip
% bookkeeping, so a following sidebody gets the right width. Swap for
% \sideimagecover[frac]{img} + \sidecaption{caption} once the asset exists.
\makeatletter
\newcommand{\phimgside}[3][0.33]{%
  \renewcommand{\td@sidefrac}{#1}\renewcommand{\td@sideside}{east}%
  \begin{tikzpicture}[remember picture,overlay]
    \node[anchor=north east, draw=tddim, thick, fill=tddim!12, inner sep=0pt,
      minimum width={#1*\paperwidth},
      minimum height={\dimexpr\paperheight-20mm\relax}]
      (phstrip) at ([xshift=-9mm,yshift=-4mm]current page.north east){};
    \node[rotate=90, anchor=west, inner sep=0pt, text=tddim,
      font=\ttfamily\scriptsize]
      at ([xshift=-2mm]phstrip.south west){#3};
  \end{tikzpicture}%
}
\makeatother

% Placeholder QR code: \phqr{caption}. A small square placeholder box
% pinned to the page's upper right corner, inset like \phimgside so it
% clears the rotated watermark, with a typewriter caption run in on its
% left. Swap for \slideqr{caption}{url} once the target url is settled;
% the caption stays as it is.
\newcommand{\phqr}[1]{%
  \begin{tikzpicture}[remember picture,overlay]
    \node[anchor=north east, draw=tddim, thick, rounded corners,
      fill=tddim!12, minimum width=9mm, minimum height=9mm,
      inner sep=0pt, text=tddim, font=\ttfamily\tiny]
      (phqr) at ([xshift=-9mm,yshift=-5mm]current page.north east){qr};
    \node[anchor=east, inner sep=0pt, text=tddim, align=right,
      font=\ttfamily\scriptsize]
      at ([xshift=-3mm]phqr.west){#1};
  \end{tikzpicture}%
}

% Live QR code: \slideqr{caption}{url}. The compiled twin of \phqr: the
% qrcode package draws the code from the url at build time. The modules
% are white and sit straight on the dark page with no tile under them,
% which makes the page itself the quiet zone; keep the surroundings clear
% for a few module widths. An inverted code is outside ISO 18004: the iOS
% camera re-inverts and reads it, most current Android scanners do too,
% older ones do not, so a print handout wants a black on white variant.
% level=M rather than the sparser L because a 57 character url lands on
% version 4 either way: same 33x33 grid, more error correction for free.
% Longer urls can cost a version, and a version is module size, which at
% 9mm is what scanning depends on.
\usepackage{qrcode}
\newcommand{\slideqr}[2]{%
  \begin{tikzpicture}[remember picture,overlay]
    \node[anchor=north east, inner sep=0pt]
      (slideqr) at ([xshift=-9mm,yshift=-5mm]current page.north east)
      {\color{white}\qrcode[height=9mm,level=M]{#2}};
    \node[anchor=east, inner sep=0pt, text=tddim, align=right,
      font=\ttfamily\scriptsize]
      at ([xshift=-3mm]slideqr.west){#1};
  \end{tikzpicture}%
}

\title[Intelligence Engineering]{Intelligence Engineering}
\subtitle{Solving Real World Problems At Scale}
\author[Schutera]{Prof.\ Dr.-Ing.\ Mark Schutera}
\institute{\textcolor{DHBWred}{DHBW}\ Ravensburg}
\date{\today}
\setbookversion{\today}{drifting stream}

% Right-edge vertical watermark: the conditional that quietly moves under the
% system while nobody touches the code. Ch11 is where that lands.
\renewcommand{\tdwatermark}{$P(y \mid x)$}

% ---------------------------------------------------------------- citations
% Sources are cited with biblatex, numbered PER CHAPTER. Each chapter file
% wraps its frames in a \begin{refsection}, so its \cite/\slidesources markers
% run [1..N] within that chapter and its own closing References frame resolves
% them. There is no deck-level bibliography; a chapter PDF pulled out of the
% set is therefore self-contained. Reuses the notes' bib. sorting=none keeps
% the numbers in citation order, so the sources on one slide compress to a
% tidy range, e.g. [1-3].
\usepackage[style=numeric-comp,sortcites,sorting=none,backend=biber]{biblatex}
\addbibresource{../intelligenceengineering.bib}
\renewcommand*{\bibfont}{\footnotesize}
\setlength{\bibitemsep}{0.35em plus 0.15em}
\AtEveryBibitem{\clearfield{url}\clearfield{doi}\clearfield{urldate}\clearfield{note}}
\usepackage{multicol}% two-column References frame (ch00 uses it)

\setbeamercolor{bibliography entry author}{fg=tddim}
\setbeamercolor{bibliography entry title}{fg=tdfg}
\setbeamercolor{bibliography entry location}{fg=tddim}
\setbeamercolor{bibliography entry note}{fg=tddim}
\setbeamercolor{bibliography item}{fg=tddim}
\DeclareFieldFormat{labelnumberwidth}{\textcolor{tddim}{#1}}

% The side-edge source marker takes cite keys and prints their numeric labels
% ([1-3]) in the same vertical spot, in place of hand-typed author-years.
% The marker sits 3mm off the right page edge, which a \sideimage or
% \sideimagecover strip covers, so grey stops reading there; on those frames it
% goes red. \phimgside is deliberately NOT covered by this: its strip is inset
% to clear the marker, so the placeholder keeps the grey.
\makeatletter
\renewcommand{\slidesources}[1]{%
  \iftd@sideimage\def\td@srccol{DHBWred}\else\def\td@srccol{tddim}\fi
  \begin{tikzpicture}[remember picture,overlay]
    \node[rotate=90,anchor=west,inner sep=0pt,text=\td@srccol,font=\scriptsize]
      at ([xshift=-3mm,yshift=13mm]current page.south east)
      {\cite{#1}};
  \end{tikzpicture}%
}

% \sidecaption shadowed for this deck. The theme sets the caption horizontal,
% white and italic INSIDE the strip, which a photo swallows; here it reads up
% the strip's outer edge in the placeholder's own typewriter grey, mirroring
% \slidesources on the opposite page edge, so a real image and its
% \phimgside placeholder caption the same way. 13mm clears the footer band.
\renewcommand{\sidecaption}[1]{%
  \setlength{\td@stripw}{\td@sidefrac\paperwidth}%
  \def\td@westtest{west}%
  \begin{tikzpicture}[remember picture,overlay]
    \ifx\td@sideside\td@westtest
      \node[rotate=90,anchor=west,inner sep=0pt,text=tddim,
            font=\ttfamily\scriptsize]
        at ([xshift=\dimexpr\td@stripw+2mm\relax,yshift=13mm]current page.south west){#1};
    \else
      \node[rotate=90,anchor=west,inner sep=0pt,text=tddim,
            font=\ttfamily\scriptsize]
        at ([xshift=\dimexpr-\td@stripw-2mm\relax,yshift=13mm]current page.south east){#1};
    \fi
  \end{tikzpicture}%
}
\makeatother

% ------------------------------------------------------------- deck helpers

% ------------------------------------------------------------- statement
% The punchline slide, overridden deck level. The theme's own statement is
% letterspaced uppercase with a red uppercase second line; both lines shout and
% the red one competes with the first instead of supporting it.
%
% This is the Disciplined Vibe closing card's idea, editorial rather than
% centred: a short red rule, the claim flush left in large bold MIXED case with
% the red carried inline by \alertbf, then the consequence one step down in
% grey. The rule is the same motif as the subchapter divider, so dividers and
% punchlines read as one family, and flush left means a wrapping claim breaks
% cleanly instead of stranding an orphan word in the middle of the slide.
%
%   \begin{statement}
%     The algorithm is the \alertbf{small box in the middle}
%     \substatement{Everything else is this course.}
%   \end{statement}
%
% The theme's versions are deliberately shadowed, not patched: other decks share
% that file. \RenewDocumentEnvironment/Command so this fails loudly if the theme
% ever stops providing them.
\RenewDocumentEnvironment{statement}{+b}{%
  \par\vfill
  {\color{DHBWred}\rule{0.10\paperwidth}{2.5pt}}\par\vskip1.1em
  {\fontsize{28}{33}\selectfont\bfseries\raggedright #1\par}%
  \vfill\par}{}
% \substatement stays \protected so the surrounding grab cannot expand it.
\RenewDocumentCommand{\substatement}{m}{%
  \par\vskip0.9em{\large\color{tddim}\raggedright #1\par}}

% ---------------------------------------------------------- frame subtitle
% The theme's frametitle template prints the title and the hairline rule and
% drops \insertframesubtitle on the floor, so \framesubtitle{...} renders
% nothing. It is reinstated here, deck level, with the theme's own template
% copied verbatim and one branch added: grey, small, above the rule.
% Used for the source's own one line under a heading, e.g. the condition that
% names each of the five keyphrases in chapter 1.
\setbeamerfont{framesubtitle}{size=\small}
\setbeamercolor{framesubtitle}{fg=tddim}
\makeatletter
\setbeamertemplate{frametitle}{%
  \nointerlineskip%
  \vskip2.4ex%
  {\usebeamerfont{frametitle}\usebeamercolor[fg]{frametitle}\strut\insertframetitle\par}%
  \ifx\insertframesubtitle\@empty\else
    \vskip0.35ex%
    {\usebeamerfont{framesubtitle}\usebeamercolor[fg]{framesubtitle}%
      \insertframesubtitle\par}%
  \fi
  \vskip0.2ex%
  {\color{tdrule}\rule{\textwidth}{0.4pt}}%
  \vskip0.4ex%
}
\makeatother

% ------------------------------------------------------- subchapter divider
% A lecture is one chapter and gets exactly ONE \lecturepage. Inside it, the
% book's own sections are subchapters, and each opens on this divider: the same
% mini title page the theme draws for a section, but driven by \subsection.
% Two reasons for \subsection rather than \section: the number restarts inside
% every lecture, and the footer keeps naming the lecture instead of switching to
% the subchapter.
% Set \subchaptergloss{...} just before \subsection{...} for the grey line under
% the title; it clears itself, so the next divider has none unless set again.
\newcommand{\subchapterglosstext}{}
\newcommand{\subchaptergloss}[1]{\gdef\subchapterglosstext{#1}}
\setbeamertemplate{subsection page}{%
  \vfill
  {\color{DHBWred}\rule{0.13\paperwidth}{2pt}}\par\vskip0.7em%
  {\usebeamercolor[fg]{section number}\large\bfseries
    \insertsubsectionnumber}\par\vskip0.25em%
  {\usebeamerfont{section title}\usebeamercolor[fg]{section title}%
    \insertsubsectionhead\par}%
  \ifdefempty{\subchapterglosstext}{}{%
    \vskip0.35em
    {\usebeamerfont{section title}\color{tddim}\subchapterglosstext\par}}%
  \gdef\subchapterglosstext{}%
  \vfill
}
\AtBeginSubsection[]{\begin{frame}\subsectionpage\end{frame}}

% \lecturepage{IV}{Training Data}{one line on what the lecture is for}
% One chapter is one lecture, and this plain frame is what a lecture opens on:
% numeral, title, red rule, then the grey gloss. Exactly one per chapter, at
% the top of the chapter file. Inside a lecture the divider is \subsection,
% which draws the subchapter page above; do not use \lecturepage for that.
% Argument 3 may be left empty; the gloss line then simply does not appear.
\newcommand{\lecturepage}[3]{%
  \begin{frame}[plain]
    \vfill
    \centering
    {\color{tddim}\footnotesize LECTURE\ \ #1}\par\vskip0.4em
    {\LARGE\bfseries #2}\par\vskip0.7em
    \textcolor{DHBWred}{\rule{0.18\textwidth}{1.2pt}}\par\vskip0.7em
    {\color{tddim}\small #3}
    \vfill
  \end{frame}}

% \increment{...} -- the one sentence of backbone work a block produces, lifted
% verbatim from the chapter header's "Increment:" line. Every block has exactly
% one. Renders as a bordered card so it reads as the deliverable, not as prose.
\newcommand{\increment}[1]{%
  \begin{tdblock}
    \alertbf{Increment.}\; #1
  \end{tdblock}}

% \tofill{<the author's gloss>} -- scaffold marker. The notes carry a one-line
% gloss under every heading and no prose yet; the gloss goes here so the frame
% says something true, and the short red rule above it says the frame is
% unfinished. Deleting the marker is what "this frame is written" means, and it
% also hands the slide's one red element back to whatever replaces it, so an
% unfilled frame is the only kind whose red is the marker.
\newcommand{\tofill}[1]{%
  \par\vskip0.3em
  {\color{DHBWred}\rule{2.2em}{0.6pt}}\par\vskip0.4em
  {\color{tddim}\itshape #1}}

% Add-on group frame: one booktabs table, "add-on" beside "the core unit it
% attaches to". The notes' default is to weave an add-on into its core unit
% rather than give it a heading, so the whole add-on set of a chapter lands on
% one frame instead of one frame each.
% Sized at \footnotesize with a tight arraystretch on purpose: every row is two
% lines (the add-on, then its source in \scriptsize), and the chapters that run
% to seven add-ons overflow the frame at \small.
\newenvironment{addontable}{%
  \vskip0.2em\footnotesize
  \renewcommand{\arraystretch}{0.95}%
  \begin{tabular}{@{}>{\raggedright\arraybackslash}p{0.40\textwidth}%
                    @{\hspace{5mm}}%
                    >{\raggedright\arraybackslash}p{0.50\textwidth}@{}}%
    \toprule
    \textbf{Add-on} & \textbf{Attaches to}\\
    \midrule}{%
    \bottomrule
  \end{tabular}}

% ------------------------------------------------------------- footer nav
% The theme's footer is beamer's mini-frame navigation: one hollow dot per
% frame, grouped under its section name. It is sized for a deck the shape of
% Disciplined Vibe, eight sections over fifty frames. This deck mirrors a
% seventeen block course: sixteen sections over about a hundred and eighty
% frames, which overruns the footer band by some 340pt on every single page.
% Neither knob fixes it, the dots alone are already three times the paper
% width. So the nav is replaced here, at deck level, with the same information
% compressed: the section you are in on the left, a progress track in the
% middle, the counter on the right. The theme is untouched and the other decks
% keep their dots. The track is drawn in tdfg on a tdrule ground rather than in
% red, so the footer never spends the slide's one red element.
% Drawn with \rule rather than tikz on purpose: a tikz coordinate resolves a
% bare length register through pgfmath, which hands back a unitless number that
% tikz then reads as centimetres, and the bar blows past the maximum dimension.
\newlength{\ieprogwd}
\newlength{\ieprogfill}
\setlength{\ieprogwd}{0.30\paperwidth}
% The fraction goes through pgfmath, which works in floating point. Doing it in
% TeX dimen arithmetic does not survive the first pass: before the aux file is
% written \inserttotalframenumber is stale, the \divide leaves \ieprogwd whole,
% and \ieprogwd times a frame number past about 120 exceeds the maximum TeX
% dimension. The clamp covers the same stale-aux case, where the current frame
% number can exceed the recorded total and the fill would run past its track.
\newcommand{\ieprogressbar}{%
  \ifnum\inserttotalframenumber>0\relax
    \ifnum\insertframenumber<\inserttotalframenumber\relax
      % The ratio is parenthesised so it is evaluated first. pgfmath keeps its
      % intermediates in a dimen register, so multiplying the track width by the
      % frame number before dividing overflows once past frame 120.
      \pgfmathsetlength{\ieprogfill}{\ieprogwd*(\insertframenumber/\inserttotalframenumber)}%
    \else
      \setlength{\ieprogfill}{\ieprogwd}%
    \fi
  \else
    \setlength{\ieprogfill}{0pt}%
  \fi
  \raisebox{0.6ex}{\makebox[\ieprogwd][l]{%
    \rlap{\textcolor{tdrule}{\rule{\ieprogwd}{1.1pt}}}%
    \textcolor{tdfg}{\rule{\ieprogfill}{1.1pt}}}}%
}
\setbeamertemplate{footline}{%
  \begin{beamercolorbox}[wd=\paperwidth,ht=6mm,dp=5mm,leftskip=6mm,rightskip=6mm]{section in head/foot}%
    {\usebeamerfont{section in head/foot}\color{tddim}\insertsectionhead}%
    \hfill
    \ieprogressbar
    \hfill
    {\usebeamerfont{section in head/foot}\color{tddim}%
      \insertframenumber\,/\,\inserttotalframenumber}%
  \end{beamercolorbox}%
}


\begin{document}

\begin{frame}[plain]
  \titlepage
\end{frame}

% ------------------------------------------------------------- front matter
% The framing frames that used to sit here, commented out behind \iffalse,
% are gone. Their job is now done by ch00_orientation.tex, which is written
% from the book's preface and the opening of its chapter 1 rather than from
% the earlier hand-written framing.

% ------------------------------------------------------------------ chapters
% One \input per book chapter, in the book's order. Each fragment carries its
% own \lecturepage divider and its \section. ch00 has no book chapter: it is
% the orientation lecture, built from the preface and the run-up to chapter 1's
% first heading, and it opens on a blank numeral (LECTURE _) for that reason.
%%% ch00_orientation.tex -- Intelligence Engineering, chapter 00 "Orientation".
%%
%% The orientation lecture. It has no book chapter of its own: segments 1 to 3
%% are Huyen's preface and segment 4 is the opening of chapter 1, everything
%% before the heading "When to Use Machine Learning". ch01 opens on exactly that
%% heading, so the two files meet without overlap and without a gap.
%%
%% STATE: titles, subtitles and TEMPLATE only. No body text. Every frame carries
%% the layout it is meant to be filled into, and a % TEMPLATE comment saying why
%% that layout. The subtitles are the book's own lines, not paraphrases.
%% The earlier prose draft is kept at _archive/ch00_orientation_prose.tex.bak.
%%
%% NUMERAL. Every other chapter opens on \lecturepage{I} .. \lecturepage{XI} and
%% the numeral equals the book chapter. This lecture has no book chapter, so its
%% numeral is a blank: LECTURE _.
%%
%% IMAGES. \phimg draws a labelled placeholder box; the label is the filename to
%% drop into images/. Swapping in the asset is
%% \includegraphics[width=\linewidth]
%% in place of the \phimg call, nothing else moves.
%%
%% Fragment of intelligence_engineering.tex. One chapter is one lecture and gets
%% exactly ONE \lecturepage, at the top. Inside it every \subsection opens on
%% the subchapter divider, emitted by the master via \AtBeginSubsection.

\begin{refsection}

\lecturepage{\_}{Orientation}{what this course is, what it assumes,
  what yor are going to know after.}
\section[Orientation]{Orientation}

% ----------------------------------------------------------------------



\begin{frame}{}
  \begin{statement}
    \substatement{You have been given a lot of raw \alert{data}, }  
    and you want to engineer the shit out of it.
  \end{statement}
\end{frame}

\begin{frame}{}
  \begin{statement}
    \substatement{You have been given a business problem, }  
    and you want to choose the right \alert{metrics and objectives} to solve it.
  \end{statement}
\end{frame}

\begin{frame}{}
  \begin{statement}
    \substatement{Your initial models perform well in offline experiments, }  
    and you
    want to \alert{deploy} them at scale.
    \end{statement}
\end{frame}

\begin{frame}{}
  \begin{statement}
    \substatement{You want to harden your deployments  }  
    to quickly detect, debug, 
and \alert{address issues} that arise in production.
  \end{statement}
\end{frame}

\begin{frame}{}
  \begin{statement}
    \substatement{There is a lot of manual work and handling involved so far,  }  
    and you want to \alert{operation vacation} your machine learning product or service.
  \end{statement}
\end{frame}

\begin{frame}{}
  \begin{statement}
    \substatement{You want to learn about a framework and foundation for building machine learning systems}  
    that can be
    \alert{shared and reused} across machine learning use cases and data products.  
  \end{statement}
\end{frame}




\subchaptergloss{Machine Learning and Data Engineering}
\subsection{Overview}

% TEMPLATE: overview table. One row per lecture, the numeral equals the
% \lecturepage numeral and the title is the book's own chapter title. The
% orientation row sits above the midrule because it has no book chapter.
\begin{frame}{The lecture in one slide}
  \sideimagecover[0.33]{placeholder_gray.png}
  \sidecaption{Machine Learning Systems Overview}
  \begin{sidebody}
    \footnotesize
    \renewcommand{\arraystretch}{1.1}%
    \begin{tabular}{@{}ll@{}}
        \toprule
        Lecture & \alert{Contents}\\
        \midrule
        \alertbf{I} & Overview of Machine Learning Systems\\
        \alertbf{II} & Introduction to Machine Learning Systems Design\\
        \alertbf{III} & Data Engineering Fundamentals\\
        \alertbf{IV} & Training Data\\
        \alertbf{V} & Feature Engineering\\
        \alertbf{VI} & Model Development and Offline Evaluation\\
        \alertbf{VII} & Model Deployment and Prediction Service\\
        \alertbf{VIII} & Data Distribution Shifts and Monitoring\\
        \alertbf{IX} & Continual Learning and Test in Production\\
        \alertbf{X} & Infrastructure and Tooling for MLOps\\
        \alertbf{XI} & The Human Side of Machine Learning\\
        \bottomrule
    \end{tabular}
  \end{sidebody}
  \slidesources{Huyen2022}
\end{frame}


% ----------------------------------------------------------------------
% TEMPLATE: the one slide, unfolded. Three lectures per frame, one column
% per lecture: the chapter header on top, its subchapters below, titles
% only. The frame title stays the overview frame's. The size command sits
% inside each column on purpose: a bare group wrapping the columns block
% silently drops it.
\begin{frame}{The lecture in one slide}
  \begin{columns}[t]
    \begin{column}{0.31\textwidth}
      \footnotesize
      \renewcommand{\arraystretch}{1.1}%
      \begin{tabular}[t]{@{}>{\raggedright\arraybackslash}p{\linewidth}@{}}
        \toprule
        \parbox[t][2\baselineskip][t]{\linewidth}{\raggedright\alertbf{I}\quad \textbf{Overview of Machine Learning Systems}}\\
        \midrule
        When to Use Machine Learning\\[0.2em]
        Understanding Machine Learning Systems\\[0.2em]
        \bottomrule
      \end{tabular}
    \end{column}
    \begin{column}{0.31\textwidth}
      \footnotesize
      \renewcommand{\arraystretch}{1.1}%
      \begin{tabular}[t]{@{}>{\raggedright\arraybackslash}p{\linewidth}@{}}
        \toprule
        \parbox[t][2\baselineskip][t]{\linewidth}{\raggedright\alertbf{II}\quad \textbf{Introduction to Machine Learning Systems Design}}\\
        \midrule
        Business and ML Objectives\\[0.2em]
        Requirements for ML Systems\\[0.2em]
        Iterative Process\\[0.2em]
        Framing ML Problems\\[0.2em]
        Mind Versus Data\\[0.2em]
        \bottomrule
      \end{tabular}
    \end{column}
    \begin{column}{0.31\textwidth}
      \footnotesize
      \renewcommand{\arraystretch}{1.1}%
      \begin{tabular}[t]{@{}>{\raggedright\arraybackslash}p{\linewidth}@{}}
        \toprule
        \parbox[t][2\baselineskip][t]{\linewidth}{\raggedright\alertbf{III}\quad \textbf{Data Engineering Fundamentals}}\\
        \midrule
        Data Sources\\[0.2em]
        Data Formats\\[0.2em]
        Data Models\\[0.2em]
        Data Storage Engines and Processing\\[0.2em]
        Modes of Dataflow\\[0.2em]
        Batch Processing Versus Stream Processing\\[0.2em]
        \bottomrule
      \end{tabular}
    \end{column}
  \end{columns}
  \slidesources{Huyen2022}
\end{frame}

\begin{frame}{The lecture in one slide}
  \begin{columns}[t]
    \begin{column}{0.31\textwidth}
      \footnotesize
      \renewcommand{\arraystretch}{1.1}%
      \begin{tabular}[t]{@{}>{\raggedright\arraybackslash}p{\linewidth}@{}}
        \toprule
        \parbox[t][2\baselineskip][t]{\linewidth}{\raggedright\alertbf{IV}\quad \textbf{Training Data}}\\
        \midrule
        Sampling\\[0.2em]
        Labeling\\[0.2em]
        Class Imbalance\\[0.2em]
        Data Augmentation\\[0.2em]
        \bottomrule
      \end{tabular}
    \end{column}
    \begin{column}{0.31\textwidth}
      \footnotesize
      \renewcommand{\arraystretch}{1.1}%
      \begin{tabular}[t]{@{}>{\raggedright\arraybackslash}p{\linewidth}@{}}
        \toprule
        \parbox[t][2\baselineskip][t]{\linewidth}{\raggedright\alertbf{V}\quad \textbf{Feature Engineering}}\\
        \midrule
        Learned Features Versus Engineered Features\\[0.2em]
        Common Operations\\[0.2em]
        Data Leakage\\[0.2em]
        Engineering Good Features\\[0.2em]
        \bottomrule
      \end{tabular}
    \end{column}
    \begin{column}{0.31\textwidth}
      \footnotesize
      \renewcommand{\arraystretch}{1.1}%
      \begin{tabular}[t]{@{}>{\raggedright\arraybackslash}p{\linewidth}@{}}
        \toprule
        \parbox[t][2\baselineskip][t]{\linewidth}{\raggedright\alertbf{VI}\quad \textbf{Model Development and Offline Evaluation}}\\
        \midrule
        Model Development and Training\\[0.2em]
        Model Offline Evaluation\\[0.2em]
        \bottomrule
      \end{tabular}
    \end{column}
  \end{columns}
  \slidesources{Huyen2022}
\end{frame}

\begin{frame}{The lecture in one slide}
  \begin{columns}[t]
    \begin{column}{0.31\textwidth}
      \footnotesize
      \renewcommand{\arraystretch}{1.1}%
      \begin{tabular}[t]{@{}>{\raggedright\arraybackslash}p{\linewidth}@{}}
        \toprule
        \parbox[t][2\baselineskip][t]{\linewidth}{\raggedright\alertbf{VII}\quad \textbf{Model Deployment and Prediction Service}}\\
        \midrule
        ML Deployment Myths\\[0.2em]
        Batch Prediction Versus Online Prediction\\[0.2em]
        Model Compression\\[0.2em]
        ML on the Cloud and on the Edge\\[0.2em]
        \bottomrule
      \end{tabular}
    \end{column}
    \begin{column}{0.31\textwidth}
      \footnotesize
      \renewcommand{\arraystretch}{1.1}%
      \begin{tabular}[t]{@{}>{\raggedright\arraybackslash}p{\linewidth}@{}}
        \toprule
        \parbox[t][2\baselineskip][t]{\linewidth}{\raggedright\alertbf{VIII}\quad \textbf{Data Distribution Shifts and Monitoring}}\\
        \midrule
        Causes of ML System Failures\\[0.2em]
        Data Distribution Shifts\\[0.2em]
        Monitoring and Observability\\[0.2em]
        \bottomrule
      \end{tabular}
    \end{column}
    \begin{column}{0.31\textwidth}
      \footnotesize
      \renewcommand{\arraystretch}{1.1}%
      \begin{tabular}[t]{@{}>{\raggedright\arraybackslash}p{\linewidth}@{}}
        \toprule
        \parbox[t][2\baselineskip][t]{\linewidth}{\raggedright\alertbf{IX}\quad \textbf{Continual Learning and Test in Production}}\\
        \midrule
        Continual Learning\\[0.2em]
        Test in Production\\[0.2em]
        \bottomrule
      \end{tabular}
    \end{column}
  \end{columns}
  \slidesources{Huyen2022}
\end{frame}

\begin{frame}{The lecture in one slide}
  \begin{columns}[t]
    \begin{column}{0.31\textwidth}
      \footnotesize
      \renewcommand{\arraystretch}{1.1}%
      \begin{tabular}[t]{@{}>{\raggedright\arraybackslash}p{\linewidth}@{}}
        \toprule
        \parbox[t][2\baselineskip][t]{\linewidth}{\raggedright\alertbf{X}\quad \textbf{Infrastructure and Tooling for MLOps}}\\
        \midrule
        Storage and Compute\\[0.2em]
        Development Environment\\[0.2em]
        Resource Management\\[0.2em]
        ML Platform\\[0.2em]
        Build Versus Buy\\[0.2em]
        \bottomrule
      \end{tabular}
    \end{column}
    \begin{column}{0.31\textwidth}
      \footnotesize
      \renewcommand{\arraystretch}{1.1}%
      \begin{tabular}[t]{@{}>{\raggedright\arraybackslash}p{\linewidth}@{}}
        \toprule
        \parbox[t][2\baselineskip][t]{\linewidth}{\raggedright\alertbf{XI}\quad \textbf{The Human Side of Machine Learning}}\\
        \midrule
        User Experience\\[0.2em]
        Team Structure\\[0.2em]
        Responsible AI\\[0.2em]
        \bottomrule
      \end{tabular}
    \end{column}
    \begin{column}{0.31\textwidth}
    \end{column}
  \end{columns}
  \slidesources{Huyen2022}
\end{frame}



% ----------------------------------------------------------------------
% Course work, two frames under one title like the repeated "lecture in
% one slide" frames. First the contract: format in the block title, the
% project brief in its body, logistics in grey at the bottom. The judging
% lens gets the second frame. Side strip on both: the bee hive.
\begin{frame}{Course Work}
  \sideimagecover[0.3]{bees}
  \sidecaption{The bee}
  \begin{sidebody}
    \begin{block}{Presentation (15\,min) + Q\&A (open)}
      You get recordings of a bee hive entrance, and you are to
      contribute by data engineering, an annotation policy, a prediction
      service or specialized tooling.
    \end{block}
    \vskip0.8em
    {\color{tddim}Open whiteboard, three slides or one page only.
    Be ready when the class starts.}
  \end{sidebody}
\end{frame}

% The judging lens: bold term + grey rest, one row per talk section, in
% the order a talk runs.
\begin{frame}{Course Work}
  \sideimagecover[0.3]{bee_mloverlay}
  \sidecaption{The bee}
  \begin{sidebody}
    What I am looking for ..\\
    \medskip

    \textbf{\alertbf{I}ntroduction}\enskip{\color{tddim}by Motivation and Literature}\\[0.5em]
    \textbf{\alertbf{M}ethod}\enskip{\color{tddim}by Explanation and or Demonstration}\\[0.5em]
    \textbf{\alertbf{R}esults}\enskip{\color{tddim}show what was achieved and how well}\\[0.5em]
    \textbf{\alertbf{D}iscussion}\enskip{\color{tddim}on findings, limitations and what's next}\\[0.5em]
    \textbf{Smart and Aesthetic}\enskip{\color{tddim}Thoughts, Designs, Ideas and \alert{Questions}.}
  \end{sidebody}
\end{frame}

% ----------------------------------------------------------------------
% The extra mile. The offer in the one block, the terms in grey under it,
% and the trade as the closing line carrying the slide's single red peak.
% Side strip: the large banana picture.
\begin{frame}{Extra Mile}
  \sideimagecover[0.33]{banana_1}
  \sidecaption{A large banana}
  \begin{sidebody}
    \begin{block}{Special projects}
      If you are in for an extra mile, let me know. I have special
      projects that I can assign to you.
    \end{block}
    \vskip0.8em
    {\color{tddim}Wrapped in the same presentation and Q\&A format, yes it will be about Bananas, and yes it will be
    more challenging, more time and effort.}

    \bigskip
    Classic loose loose short term, yet \alertbf{win long term}.
  \end{sidebody}
\end{frame}


\begin{frame}{}
  \slideqr{The Hive is Real. Dataset on Hugging Face.}{https://huggingface.co/datasets/mrkschtr/thehiveisreal}
  \begin{statement}
    \substatement{By the end of this lecture, }
    Get yourself a \alert{presentation slot} and brainstorm a project \alert{idea to work} on.
    \end{statement}
\end{frame}

% ----------------------------------------------------------------------
% This chapter's own references. The refsection opened above scopes the
% citation numbering to this chapter, so chapter_pdfs/<this>.pdf resolves
% its own [n] markers. Two columns and a step down in size: this lecture
% cites the whole bib and it does not fit one column at \footnotesize.
\begin{frame}{References}
  \renewcommand*{\bibfont}{\scriptsize}%
  \setlength{\bibitemsep}{0.25em}%
  \begin{multicols}{2}
    \printbibliography[heading=none]
  \end{multicols}
\end{frame}
\end{refsection}

%%% ch01_overview.tex -- Intelligence Engineering, chapter 01 "Overview of
%% Machine Learning Systems".
%% Spine transferred from Huyen, Designing Machine Learning Systems, chapter 1.
%% The subchapters, the frame titles and the frame subtitles are the book's own
%% headings and its own one-line conditions, in the book's order.
%%
%% STATE: titles, subtitles and TEMPLATE only. No body text. Every frame carries
%% the layout it is meant to be filled into, and a % TEMPLATE comment saying why
%% that layout and what belongs in each slot. Filling them is the next pass.
%%
%% Fragment of intelligence_engineering.tex. One chapter is one lecture and
%% gets exactly ONE \lecturepage, at the top. Inside it every \subsection
%% opens on the subchapter divider, which the master emits automatically via
%% \AtBeginSubsection; do not write those divider frames by hand.

\begin{refsection}

\lecturepage{I}{Overview of a Machine Learning System}{}
\section[Overview]{Overview of a Machine Learning System}

% ======================================================================
\subchaptergloss{Don't confuse the model with the machine learning system.}
\subsection{The small part and the big picture}

% TEMPLATE: withmargin + tddef. The definition anchors the whole subchapter, so
% it takes the slide's one titled box; the five keyphrases are numbered inside
% it and each gets a frame of its own below. Margin holds the caveat that
% qualifies the definition (not a magic tool, not always optimal).
\begin{frame}{Enabling the delivery of the logic}
  \begin{withmargin}
    \begin{tdmain}
      The \textbf{model} is the small part of a machine learning system. 
      With foundation models and AI engineering, that part got smaller still. 
      The essential work of a machine learning system remains and stacks around business understanding, infrastructure, design pattern and data engineering.
    \end{tdmain}
    \begin{tdmargin}
      Before learning about solving a problem with machine learning, let's first assess when it is actually a good idea to do so.  
    \end{tdmargin}
  \end{withmargin}
\end{frame}

% TEMPLATE: withmargin. Main column takes the condition and what follows from
% it; margin takes the book's worked example (the Airbnb listing and its price).
\begin{frame}{(1) Learn}
  \begin{withmargin}
    \begin{tdmain}
        The problem you are solving, requires your system to have the capability to learn.
        Learning can look different depending on the problem and its context and boundary conditions.
        You already know about supervised learning, unsupervised learning and reinforcement learning paradigms.
    \end{tdmain}
    \begin{tdmargin}
        Machine learning is an approach to (1) learn (2) complex patterns
        from (3) existing data and use these patterns to make (4)
        predictions on (5) unseen data.
    \end{tdmargin}
  \end{withmargin}
\end{frame}

\begin{frame}{(2) Pattern}
    \begin{withmargin}
      \begin{tdmain}
          Learning requires the existence of a pattern, as opposed to noise or randomness. 
          It is not trivial nor obvious whether there is a pattern in the data, and whether it is learnable.
          Complex pattern is more difficult to embed or state explicitly, this is where machine learning has an edge over classical approaches.
      \end{tdmain}
      \begin{tdmargin}
          Machine learning is an approach to (1) learn (2) complex patterns
          from (3) existing data and use these patterns to make (4)
          predictions on (5) unseen data.
      \end{tdmargin}
    \end{withmargin}
  \end{frame}


  \begin{frame}{(3) Data}
    \begin{withmargin}
      \begin{tdmain}
          Data is available, accessible or acquirable. 
          The data is of sufficient quality and quantity to support learning, (and to validate for that matter).
          If your approach does not require data, it is because you build on embeddings or foundation models, 
          which themselves are built on data in the first place, or you need to rely on data acquisition and curation as part of your system.
      \end{tdmain}
      \begin{tdmargin}
          Machine learning is an approach to (1) learn (2) complex patterns
          from (3) existing data and use these patterns to make (4)
          predictions on (5) unseen data.
      \end{tdmargin}
    \end{withmargin}
  \end{frame}

  \begin{frame}{(4) Predictions}
    \begin{withmargin}
      \begin{tdmain}
          Machine learning is about making predictions, always. Predictions can be of different types, such as 
          classification, regression, ranking, clustering, generation and on and on.
          By reframing problems in terms of predictions, we can leverage machine learning to solve them.
          Predictions are approximations by nature and therefore come with uncertainty, 
          which needs to be accounted for.

      \end{tdmain}
      \begin{tdmargin}
          Machine learning is an approach to (1) learn (2) complex patterns
          from (3) existing data and use these patterns to make (4)
          predictions on (5) unseen data.
      \end{tdmargin}
    \end{withmargin}
  \end{frame}

  \begin{frame}{(5) Unseen Data}
    \begin{withmargin}
      \begin{tdmain}
          In a closed system or domain, the data is finite and the system can be fully specified by optimization.
          However, often the domain is open and infinite, and the system needs to be able to generalize to unseen data.
          Of course, the unseen data is not completely unknown, it is drawn from the same distribution (i.i.d.) as the 
          training data and is thus able to generalize to the unseen data inside a specific domain or distribution.
      \end{tdmain}
      \begin{tdmargin}
          Machine learning is an approach to (1) learn (2) complex patterns
          from (3) existing data and use these patterns to make (4)
          predictions on (5) unseen data.
      \end{tdmargin}
    \end{withmargin}
  \end{frame}


% TEMPLATE: two by two block grid. Four peer characteristics, none subordinate
% to another, so equal boxes rather than a bullet list. The equal height group
% keeps each row level.
\begin{frame}{Catalysts}
    A catalyst is a substance that increases the rate of reaction without 
    undergoing change itself.
  \begin{columns}[T,onlytextwidth]
    \begin{column}{0.47\textwidth}
      \begin{tdblockt}[equal height group=cat]{The task is repetitive.}
        In a sense that we have a lot of data to learn a pattern from.
      \end{tdblockt}
      \vskip0.5em
      \begin{tdblockt}[equal height group=cat]{The context changes.}
        Meaning that hard coded or other solutions will be brittle and break.
      \end{tdblockt}
    \end{column}
    \begin{column}{0.47\textwidth}
      \begin{tdblockt}[equal height group=cat]{It scales.}
        The need to scale a specific prediction capability, not only a few times.
      \end{tdblockt}
      \vskip0.5em
      \begin{tdblockt}[equal height group=cat]{Being wrong is not too bad.}
        Approximations on unseen data. Failures need to be manageable.  
      \end{tdblockt}
    \end{column}
  \end{columns}
\end{frame}

% TEMPLATE: three block row plus a closing line. The book gives exactly three
% disqualifying conditions, so one box each rather than three bullets; the line
% under them is the escape hatch, decompose and use ML on a component.
\begin{frame}{When Not to Use Machine Learning}
  \begin{columns}[T,onlytextwidth]
    \begin{column}{0.31\textwidth}
      \begin{tdblockt}[equal height group=not]{It is unethical.}
        Well, that might also be a general rule. But doing stuff at scale, without human oversight, has its implications.$^*$
      \end{tdblockt}
    \end{column}
    \begin{column}{0.31\textwidth}
      \begin{tdblockt}[equal height group=not]{There's a simple solution.}
        A simple solution, heuristic, or baseline should always be the first step regardless.
      \end{tdblockt}
    \end{column}
    \begin{column}{0.31\textwidth}
      \begin{tdblockt}[equal height group=not]{It's not economically viable.}
        The value of the prediction, needs to be greater than the cost of the system, including development, deployment and maintenance.
      \end{tdblockt}
    \end{column}
  \end{columns}
  \vskip0.8em
  $^*$ More? Enroll to my Philosophy of AI course.
  \vskip0.8em
\end{frame}







\begin{frame}{}
    \begin{statement}
      \substatement{Describe two use-cases and the rationale for whether or not to use machine learning.}  
      \alert{One} Machine Learning Use-Case and \alert{One} which is not.
      \end{statement}
\end{frame}











% The next four frames map value over technical feasibility and how that map
% moved between 2018 and 2026. Order: the 2018 baseline, the Eisenhower matrix
% in the middle, the 2026 survey numbers, the labs' own measurements. Four
% rather than the asked two to three, so there is something to cut.

% TEMPLATE: two block contrast plus a dimmed line. The 2018/2019 baseline is a
% tension between two documents, the trillion-dollar map and the months-long
% deployment reality, so each takes one levelled box and the dim line reads
% them against each other.
\begin{frame}{Mapping the value, 2018}
  \begin{columns}[T,onlytextwidth]
    \begin{column}{0.47\textwidth}
      \begin{tdblockt}[equal height group=map18]{The map: value in the trillions}
        \footnotesize
        McKinsey sized \alertbf{400+} use cases across \alertbf{19} industries and nine business functions: deep learning promised \alertbf{\$3.5 to \$5.8 trillion} in value per year, the largest shares in marketing and sales plus supply chain and manufacturing.
      \end{tdblockt}
    \end{column}
    \begin{column}{0.47\textwidth}
      \begin{tdblockt}[equal height group=map18]{The ground: months per model}
        \footnotesize
        Algorithmia surveyed \alertbf{745} practitioners in October 2019: most companies needed \alertbf{8 to 90 days} to deploy a single model, 18\% needed longer than 90 days, and deployment time grew year over year.
      \end{tdblockt}
    \end{column}
  \end{columns}
  \vskip0.8em
  {\color{tddim}The value was mapped in trillions while shipping one model took a quarter: feasibility, not value, was the bottleneck.}
  \slidesources{Chui2018,Algorithmia2020}
\end{frame}

% TEMPLATE: withmargin with the Eisenhower matrix in the main column. Redrawn
% as a quadrant matrix rather than a scatter: nobody measured the feasibility
% of customer service, so plotted dots claimed a precision the data does not
% have, and eleven of them collided with each other and with the quadrant
% verbs. Now each quadrant is its verb plus the domains that sit in it. The
% frontier between hard and works is the one red element and it carries the
% claim: it moves right. Margin holds the numbers that moved it, one per
% source, so the diagram states the argument and the margin evidences it.
\begin{frame}{Value over technical feasibility}
  \begin{withmargin}
    \begin{tdmain}
      \centering
      \begin{tikzpicture}[
          ax/.style={tddim, semithick, -{Latex[length=4pt,width=3pt]}},
          split/.style={tddim, thin, densely dashed},
          quad/.style={text=tddim, font=\ttfamily\scriptsize, anchor=north west},
          dom/.style={text=tdfg, font=\tiny, anchor=north west, align=left,
                      inner sep=0pt}]
        % axes: left and bottom only, no enclosing box
        \draw[ax] (0,0) -- (0,4.40);
        \draw[ax] (0,0) -- (7.75,0);
        % the value split is a plain guide
        \draw[split] (0,2.10) -- (7.6,2.10);
        % the feasibility split is the frontier, and it moves
        \draw[DHBWred, semithick, densely dashed] (3.8,0) -- (3.8,4.25);
        \draw[DHBWred, semithick, -{Latex[length=4pt,width=3pt]}]
          (3.8,4.45) -- (4.6,4.45);
        \node[text=DHBWred, font=\tiny, anchor=west] at (4.68,4.45)
          {the frontier moves};
        % x axis: technical feasibility
        \node[tddim, font=\tiny, anchor=north west] at (0,-0.08) {hard today};
        \node[tddim, font=\tiny, anchor=north east] at (7.6,-0.08) {works today};
        \node[tddim, font=\scriptsize, anchor=north] at (3.8,-0.45)
          {technical feasibility};
        % y axis: business value
        \node[tddim, font=\tiny, rotate=90, anchor=south west] at (-0.08,0) {modest};
        \node[tddim, font=\tiny, rotate=90, anchor=south east] at (-0.08,4.25)
          {transformative};
        \node[tddim, font=\scriptsize, rotate=90, anchor=south] at (-0.52,2.10)
          {business value};
        % invest: high value, still hard
        \node[quad] at (0.20,4.20) {invest};
        \node[dom] at (0.20,3.80) {autonomous driving\\medical diagnosis\\agentic workflows end to end};
        % deploy now: high value, works today
        \node[quad] at (4.05,4.20) {deploy now};
        \node[dom] at (4.05,3.80) {coding agents\\recommender systems\\fraud detection\\machine translation};
        % skip: hard and modest
        \node[quad] at (0.20,1.95) {skip};
        \node[dom] at (0.20,1.55) {bespoke one off models};
        % automate cheaply: works today, modest value
        \node[quad] at (4.05,1.95) {automate cheaply};
        \node[dom] at (4.05,1.55) {speech transcription\\document OCR};
      \end{tikzpicture}
    \end{tdmain}
    \begin{tdmargin}
      SWE-bench Verified rose from 60\% to near 100\% within a year.
      \vskip0.5em
      Claude Opus 4.1 wins or ties against human experts on 47.6\% of GDPval tasks.
      \vskip0.5em
      Among organizations above \$1 billion revenue, 40\% now scale AI agents, up from 27\% in one year.
    \end{tdmargin}
  \end{withmargin}
  \slidesources{AIIndex2026,OpenAI2025,McKinsey2026}
\end{frame}

% TEMPLATE: side image hero plus two label led columns. Mark's call, and the
% reason holds: a pair of cards frames the two shares as peers of equal weight
% when the point is that they disagree, so the boxes come off and the column
% split alone carries the contrast. The strip takes the image; the headings
% survive as bold lead lines rather than block titles.
\begin{frame}{The field is moving, 2026}
  \sideimagecover[0.33]{placeholder_gray.png}
  % Asset brief: the adoption curve against the value capture curve on one set
  % of axes, 2023 to 2026, the two diverging. Caption stays under ~50
  % characters: it is a rotated node with no text width and the page clips
  % whatever runs past the top edge.
  \sidecaption{Adoption climbs, value capture flat}
  \begin{sidebody}
    \scriptsize
    \begin{columns}[T,onlytextwidth]
      \begin{column}{0.47\linewidth}
        \textbf{Adoption is near universal}
        \begin{itemize}
          \item \alertbf{88\%} of organizations use AI in at least one business function, \alertbf{72\%} use gen AI.
          \item Gen AI reached 53\% population adoption within three years, faster than the PC or the internet.
          \item US private AI investment reached \$285.9 billion in 2025.
        \end{itemize}
      \end{column}
      \hfill
      \begin{column}{0.47\linewidth}
        \textbf{Value capture is not}
        \begin{itemize}
          \item \alertbf{80\%} of AI users report individual productivity gains.
          \item \alertbf{37\%} of organizations attribute any EBIT impact to AI, flat since 2025.
          \item \alertbf{6\%} qualify as AI high performers.
        \end{itemize}
      \end{column}
    \end{columns}
  \end{sidebody}
  \slidesources{McKinsey2026,AIIndex2026}
\end{frame}

% TEMPLATE: two block contrast plus a dimmed line. The two labs publish the
% same map from their own vantage points, a benchmark against experts and a
% telescope on live traffic, so one box per lab and the dim line reads the
% two together.
\begin{frame}{What the labs measure}
  \begin{columns}[T,onlytextwidth]
    \begin{column}{0.47\textwidth}
      \begin{tdblockt}[equal height group=labs]{OpenAI: GDPval}
        \footnotesize
        \alertbf{1{,}320} tasks from \alertbf{44} occupations across nine sectors, graded head to head by experts averaging 14 years in the field. Frontier models finish them about \alertbf{100$\times$} faster and \alertbf{100$\times$} cheaper than the experts.
      \end{tdblockt}
    \end{column}
    \begin{column}{0.47\textwidth}
      \begin{tdblockt}[equal height group=labs]{Anthropic: Economic Index}
        \footnotesize
        \alertbf{49\%} of jobs show Claude on at least a quarter of their tasks. Augmentation leads automation \alertbf{52\% to 45\%}. The single most common task: modifying software to correct errors, 6\% of all usage.
      \end{tdblockt}
    \end{column}
  \end{columns}
  \vskip0.8em
  {\color{tddim}Broad across occupations, thin within them, and the densest use is code: the labs' own traffic redraws the feasibility axis every quarter.}
  \slidesources{OpenAI2025,Anthropic2026}
\end{frame}











% ======================================================================
\subchaptergloss{Research and Production are not the same.}
\subsection{Understanding Machine Learning Systems}

% TEMPLATE: comparison table. Table 1-1 of the book is already a table and
% should stay one: five rows, research against production. booktabs only, no
% vertical rules. The five row labels are the book's.
\begin{frame}{Horizontal versus Vertical}
  \footnotesize
  Nothing motivates \alert{system} thinking more than the need to go \alert{horizontal}.
  Academia is vertical, with a narrow focus on a single problem, where drilling deep, fast is rewarded and
  focus allows to teach a wide range of concepts.
  \vspace{1.5em}

  \renewcommand{\arraystretch}{1.25}
  
  \begin{tabular}{@{}%
    >{\raggedright\arraybackslash}p{0.22\textwidth}%
    >{\raggedright\arraybackslash}p{0.35\textwidth}%
    >{\raggedright\arraybackslash}p{0.35\textwidth}@{}}
    \toprule
     & \textbf{Research} & \textbf{Production}\\
    \midrule
    \textbf{Objectives} & State of the Art highly engineered models &  Stakeholder outcome driven \\
    \textbf{Compute} & Training and high throughput &  Fast inference with low latency \\
    \textbf{Data} & Often static benchmark oriented & Shifting domains and structures  \\
    \textbf{Usability metrics} & Usually not considered & UI, Fairness, Robust, Interpretability \\
    \bottomrule
  \end{tabular}
\end{frame}









% TEMPLATE: side image hero plus a two row card grid. The five stakeholders are
% peers, none subordinate to another, so a grid rather than a stack, and two
% rows keep the sidebody inside its height budget. Each card carries an
% objective and the solution space its owner already has in mind, which are the
% two axes of the Eisenhower map in the strip. No red on the cards: the one red
% element on this frame is the term the whole slide is about.
\begin{frame}{Stakeholders and Objectives}
  \sideimagecover[0.33]{placeholder_gray.png}
  % Asset brief: an Eisenhower map, concern or objective over solution space,
  % both axes running from business to technical. The visible caption has to
  % stay under ~50 characters: it is a rotated node with no text width and
  % the page clips whatever runs past the top edge.
  \sidecaption{Eisenhower: objective over solution space}
  \begin{sidebody}
    \scriptsize
    Tinder for horses. In research the objective is one number on one benchmark.
    Ship the same model as a product and five people want five different things
    from it, each arriving with a solution space already in mind:
    \alertbf{conflicting requirements} are the norm, and not all of them are must haves.
    \vskip0.7em
    \tiny
    \begin{columns}[T,onlytextwidth]
      \begin{column}{0.315\linewidth}
        \begin{tdblockt}[coltitle=tdfg,before upper=\raggedright,equal height group=sh,left=4pt,right=4pt,toptitle=1pt,top=1pt,bottom=2pt]{ML engineers}
          {\color{tddim}Objective} better matches.\par
          {\color{tddim}Solution space} more data, a bigger model.
        \end{tdblockt}
      \end{column}
      \hfill
      \begin{column}{0.315\linewidth}
        \begin{tdblockt}[coltitle=tdfg,before upper=\raggedright,equal height group=sh,left=4pt,right=4pt,toptitle=1pt,top=1pt,bottom=2pt]{Platform team}
          {\color{tddim}Objective} nothing pages at 3am.\par
          {\color{tddim}Solution space} freeze the model, build the platform.
        \end{tdblockt}
      \end{column}
      \hfill
      \begin{column}{0.315\linewidth}
        \begin{tdblockt}[coltitle=tdfg,before upper=\raggedright,equal height group=sh,left=4pt,right=4pt,toptitle=1pt,top=1pt,bottom=2pt]{Product team}
          {\color{tddim}Objective} the swipe feels instant.\par
          {\color{tddim}Solution space} a latency budget of 100 ms.
        \end{tdblockt}
      \end{column}
    \end{columns}
    \vskip0.35em
    \begin{columns}[T,onlytextwidth]
      \begin{column}{0.315\linewidth}
        \begin{tdblockt}[coltitle=tdfg,before upper=\raggedright,equal height group=sh2,left=4pt,right=4pt,toptitle=1pt,top=1pt,bottom=2pt]{Sales team}
          {\color{tddim}Objective} revenue per match.\par
          {\color{tddim}Solution space} rank the premium stallions first.
        \end{tdblockt}
      \end{column}
      \hfill
      \begin{column}{0.315\linewidth}
        \begin{tdblockt}[coltitle=tdfg,before upper=\raggedright,equal height group=sh2,left=4pt,right=4pt,toptitle=1pt,top=1pt,bottom=2pt]{Manager}
          {\color{tddim}Objective} margin.\par
          {\color{tddim}Solution space} fewer models, and maybe fewer engineers.
        \end{tdblockt}
      \end{column}
      \hfill
      \begin{column}{0.315\linewidth}
      \end{column}
    \end{columns}
  \end{sidebody}
  \slidesources{Huyen2022}
\end{frame}


% TEMPLATE: side image hero plus a bullet list. Three separate charges against
% the same practice, so a bullet each rather than boxes; none of them is a
% comparison, so no paired blocks. The strip takes the joke the criticism is
% named after. Picard's seed scan is the whole argument of the second bullet,
% so its numbers sit in the bullet and not in a caption.
\begin{frame}{Leaderboards and Benchmarks}
  \sideimagecover[0.33]{placeholder_gray.png}
  % Asset brief: a trust me bro benchmark chart, bars with no axis, no
  % baseline and no error bars, posted to announce a model. Caption kept
  % under ~50 characters, see the note on the frame above.
  \sidecaption{Trust me bro benchmark: no axis, no error bars}
  \begin{sidebody}
    \scriptsize
    A leaderboard scores the \alertbf{model}. A system is not a model.
    \vskip0.5em
    \begin{itemize}
      \item What a benchmark measures is one metric on one frozen test set. What it
        never measures is the pipeline around it: acquisition, cleaning, labelling,
        feature engineering, deployment, monitoring.
      \item Thousands of entries against the same test set is multiple hypothesis
        testing, so part of the podium is pure chance. Scanning 10,000 seeds on
        CIFAR 10 spreads accuracy from 89.01 to 90.83 percent, 1.82 points, with
        nothing changed but the seed.
      \item A benchmark is a proxy for a single metric. It was never forged in
        production, where objectives conflict, data shifts, and the users are real.
    \end{itemize}
  \end{sidebody}
  \slidesources{Picard2021}
\end{frame}

\begin{frame}{Computational priorities}
  \begin{columns}[T,onlytextwidth]
    \begin{column}{0.47\textwidth}
      \begin{tdblockt}[equal height group=cp]{Research: fast training}
        \scriptsize
        \begin{itemize}
          \item Training is the bottleneck: many models, each doing multiple passes over the data.
          \item Prioritizes \alertbf{high throughput}, how many queries are processed in a period of time.
        \end{itemize}
      \end{tdblockt}
    \end{column}
    \begin{column}{0.47\textwidth}
      \begin{tdblockt}[equal height group=cp]{Production: fast inference}
        \scriptsize
        \begin{itemize}
          \item Once deployed, inference is the bottleneck: the model's job is to generate predictions.
          \item Prioritizes \alertbf{low latency}, the time from receiving a query to returning the result.
        \end{itemize}
      \end{tdblockt}
    \end{column}
  \end{columns}
  \vskip0.3em
  \scriptsize
  {\centering
    \begin{tabular}{@{}lll@{}}
      \toprule
      \textbf{Processing} & \textbf{Average latency} & \textbf{Throughput}\\
      \midrule
      One query at a time & 10 ms & 100 queries/second\\
      One query at a time & 100 ms & 10 queries/second\\
      Batches of 10 queries & 10 ms & 1,000 queries/second\\
      Batches of 50 queries & 20 ms & 2,500 queries/second\\
      \bottomrule
    \end{tabular}\par}
  \vskip0.3em
  {\color{tddim}Batched, higher latency might also mean higher throughput.}
  \slidesources{Huyen2022}
\end{frame}

% TEMPLATE: withmargin plus a takeaway. The main column runs the book's ten
% measured latencies and its percentile argument in order, because the point is
% arithmetic. Margin holds the three dated production studies as the evidence
% that latency matters. The takeaway is the Amazon observation, the reason the
% tail of the distribution is not a rounding error.
\begin{frame}{Latency is a distribution}
  \begin{withmargin}
    \begin{tdmain}
      \small
      Ten requests, latency in ms: 100, 102, 100, 100, 99, 104, 110, 90, \alertbf{3,000}, 95.

      The average latency is 390 ms, which makes your system seem slower than it actually is.

      Better to think in percentiles: p50 is the median, p90 and p99 expose the outliers.

      The p90 here is 3,000 ms, the one request to investigate.
    \end{tdmain}
    \begin{tdmargin}
      Latency matters a lot.
      \vskip0.5em
      2017, Akamai: a 100 ms delay can hurt conversion rates by 7\%.
      \vskip0.5em
      2019, Booking.com: about 30\% more latency cost about 0.5\% in conversion rates, ``a relevant cost for our business''.
      \vskip0.5em
      2016, Google: more than half of mobile users leave a page that takes more than three seconds to load.
    \end{tdmargin}
  \end{withmargin}
  \begin{takeaway}
    On Amazon the slowest requests come from the most valuable customers.
  \end{takeaway}
  \slidesources{Huyen2022}
\end{frame}

% TEMPLATE: two block contrast. The book's Data subheading is a straight
% research against production split, so each world takes one box, headed with
% its own entry from Table 1-1 and filled with its data properties in the
% book's words. The dimmed line under them is the book's Karpathy anchor for
% the same contrast.
\begin{frame}{Data}
  \begin{columns}[T,onlytextwidth]
    \begin{column}{0.47\textwidth}
      \begin{tdblockt}[equal height group=data]{Research: static}
        \footnotesize
        \begin{itemize}
          \item Clean and well formatted
          \item Static, so the community can benchmark
          \item Quirks of the dataset are known
          \item Open source scripts to feed your models
          \item Historical data, already stored somewhere
        \end{itemize}
      \end{tdblockt}
    \end{column}
    \begin{column}{0.47\textwidth}
      \begin{tdblockt}[equal height group=data]{Production: constantly shifting}
        \footnotesize
        \begin{itemize}
          \item Messy, noisy, possibly unstructured
          \item \alert{Likely biased, and you do not know how}
          \item Labels sparse, imbalanced, or incorrect
          \item Changing requirements can mean relabeling
          \item Privacy and regulatory concerns
          \item Generated by users, systems, third parties
        \end{itemize}
      \end{tdblockt}
    \end{column}
  \end{columns}
  \vskip0.5em
  {\color{tddim}\footnotesize Andrej Karpathy, director of AI at Tesla, drew the data problems of his PhD against those of his time at Tesla.}
  \slidesources{Huyen2022}
\end{frame}

% TEMPLATE: one quote block between two lines. The book's fairness point is a
% quotation and its rebuttal rather than a comparison, so the deferral gets the
% slide's one box and the verdict lands inside it. The line above is the
% research phase condition, the line below the asymmetry, and the takeaway the
% leaderboard that does not exist.
\begin{frame}{Fairness}
  During the research phase, a model is not yet used on people, so it is easy for researchers to put off fairness as an afterthought.
  \vskip0.4em
  \begin{tdblock}
    {\color{tddim}``Let's try to get state of the art first and worry about fairness when we get to production.''}
    \vskip0.5em
    \alertbf{When it gets to production, it is too late.}
  \end{tdblock}
  \vskip0.4em
  Optimize for better accuracy or lower latency and you can show that you beat state of the art.
  \begin{takeaway}
    As of writing the book, there is no equivalent state of the art for fairness metrics.
  \end{takeaway}
  \slidesources{Huyen2022}
\end{frame}

% TEMPLATE: withmargin with a bullet list. The book argues this one by naming
% victims, so the main column is one book example per bullet, opened by the
% line that explains them all and closed by what deployment at scale does to
% them. Margin takes the sourced evidence, the Berkeley lending finding and the
% McKinsey number that says how few companies act on it.
\begin{frame}{Fairness at scale}
  \begin{withmargin}
    \begin{tdmain}
      \small
      Machine learning algorithms do not predict the future, they encode the past, thus perpetuating the biases in the data.
      \begin{itemize}
        \item Your loan application might be rejected because the algorithm picks on your \alert{zip code}, which embodies biases about one's socioeconomic background.
        \item Your resume might be ranked lower because the ranking system employers use picks on the \alert{spelling of your name}.
        \item Your mortgage might get a higher interest rate because it relies partially on \alert{credit scores}, which favor the rich and punish the poor.
      \end{itemize}
      \alertbf{Deployed at scale, they discriminate at scale}, judging millions in split seconds.
    \end{tdmain}
    \begin{tdmargin}
      In 2019, Berkeley researchers found that both face-to-face and online lenders rejected \textbf{1.3 million} creditworthy Black and Latino applicants between 2008 and 2015. With the income and credit scores kept and the race identifiers deleted, the mortgage application was accepted.
      \vskip0.5em
      McKinsey \& Company, 2019: only \textbf{13\%} of the large companies surveyed said they are taking steps to mitigate risks to equity and fairness.
    \end{tdmargin}
  \end{withmargin}
  \slidesources{Huyen2022}
\end{frame}

% TEMPLATE: two block contrast plus a dimmed closing line. Hinton's question is
% a forced choice between two options, so each option takes one box and the two
% cure rates read off against one another. The lines above set the scenario, the
% quoted question sits under the boxes, and the dimmed line is what a room of
% executives answered when the book's author put the question to them.
\begin{frame}{Interpretability}
  In early 2020, Turing Award winner Geoffrey Hinton posed a heatedly debated question.
  Suppose you have cancer and you have to choose.
  \vskip0.9em
  \begin{columns}[T,onlytextwidth]
    \begin{column}{0.47\textwidth}
      \begin{tdblockt}[equal height group=surgeon]{Black box AI surgeon}
        \alertbf{90\%} cure rate, and it cannot explain how it works.
      \end{tdblockt}
    \end{column}
    \begin{column}{0.47\textwidth}
      \begin{tdblockt}[equal height group=surgeon]{Human surgeon}
        \alertbf{80\%} cure rate.
      \end{tdblockt}
    \end{column}
  \end{columns}
  \vskip0.9em
  \alertbf{``Do you want the AI surgeon to be illegal?''}
  \vskip0.6em
  {\color{tddim}Put to a group of 30 technology executives at public nontech companies, half of them wanted the AI surgeon to operate on them. The other half wanted the human surgeon.}
  \slidesources{Huyen2022}
\end{frame}

% TEMPLATE: withmargin. The main column carries the asymmetry this subchapter
% is built on, research optimizes one objective and industry owes an
% explanation, with the book's two reasons as bullets. Margin holds why the
% discomfort is specific to AI, and the number showing that a requirement is
% not yet a practice.
\begin{frame}{Interpretability in production}
  \begin{withmargin}
    \begin{tdmain}
      \small
      Most machine learning research is still evaluated on a single objective, model performance, so researchers are not incentivized to work on model interpretability.
      In industry, interpretability is not just optional for most use cases but a requirement.
      \begin{itemize}
        \item Important for \alertbf{users}, both business leaders and end users, to understand why a decision is made, so that they can trust a model and detect potential biases.
        \item Important for \alertbf{developers}, to be able to debug and improve a model.
      \end{itemize}
    \end{tdmain}
    \begin{tdmargin}
      Most of us are comfortable using a microwave without understanding how it works. Many do not feel the same way about AI, especially if that AI makes important decisions about their lives.
      \vskip0.5em
      As of 2019, only \textbf{19\%} of large companies are working to improve the explainability of their algorithms.
    \end{tdmargin}
  \end{withmargin}
  \slidesources{Huyen2022}
\end{frame}


% TEMPLATE: single block plus a takeaway. The book closes the research versus
% production subchapter with an argument and its rebuttal rather than a
% comparison, so the claim takes the slide's one box and the takeaway carries
% what the book answers. Neither a margin nor a pair of boxes, because both
% neighbours already use those.
\begin{frame}{Discussion}
  \begin{tdblock}
  \end{tdblock}
  \begin{takeaway}
  \end{takeaway}
  \slidesources{Huyen2022}
\end{frame}

% TEMPLATE: two block contrast plus a takeaway. Departs from the stub's
% withmargin on purpose: the frame before it is a withmargin and a second in a
% row stalls, and the book's point here is symmetric, what software engineering
% assumes against what a machine learning system is actually made of, which
% reads as two levelled boxes rather than a claim with an aside. Each box holds
% one world in the book's own properties; the takeaway is the book's own answer
% to the question the subheading opens with.
\begin{frame}{Machine Learning Systems Versus Traditional Software}
  \begin{columns}[T,onlytextwidth]
    \begin{column}{0.47\textwidth}
      \begin{tdblockt}[equal height group=mlsw]{Traditional software engineering}
        \scriptsize
        \begin{itemize}
          \item An underlying assumption that \alertbf{code and data are separated}.
          \item Keep things as modular and separate as possible: separation of concerns.
          \item Many of its tools do work for machine learning.
        \end{itemize}
      \end{tdblockt}
    \end{column}
    \begin{column}{0.47\textwidth}
      \begin{tdblockt}[equal height group=mlsw]{Machine learning systems}
        \scriptsize
        \begin{itemize}
          \item \alertbf{Part code, part data}, and part artifacts created from the two.
          \item Applications with the most and best data win: companies improve data, not algorithms.
          \item Data can change quickly, so applications need to be adaptive.
        \end{itemize}
      \end{tdblockt}
    \end{column}
  \end{columns}
  \begin{takeaway}
    ML production would be a much better place if ML experts were better software engineers.
  \end{takeaway}
  \slidesources{Huyen2022}
\end{frame}

% TEMPLATE: withmargin, the layout this chapter uses when one claim needs one
% worked example to land, and the stub's intent honoured one frame later. Main
% column runs the book's argument in its order: code only, then code and data,
% then the two open questions, then the inequality of samples, closing on the
% risk that follows. Margin holds the book's lung scan counts, which are the
% whole proof that samples are not equal.
\begin{frame}{Test and version the data too}
  \begin{withmargin}
    \begin{tdmain}
      \small
      In traditional software engineering, you only need to focus on testing and versioning your code.
      With machine learning, we have to test and version our data too, and \alertbf{that is the hard part}.
      \begin{itemize}
        \item How to version large datasets?
        \item How to know if a data sample is good or bad for your system?
      \end{itemize}
      Not all data samples are equal. Some are more valuable to your model than others.
      \vskip0.3em
      \alertbf{Indiscriminately accepting all available data might hurt performance and make your model susceptible to data poisoning attacks.}
    \end{tdmain}
    \begin{tdmargin}
      Your model has already trained on \textbf{one million} scans of normal lungs and only \textbf{one thousand} scans of cancerous lungs.
      \vskip0.5em
      A scan of a cancerous lung is much more valuable than a scan of a normal lung.
    \end{tdmargin}
  \end{withmargin}
  \slidesources{Huyen2022}
\end{frame}

% TEMPLATE: three block row, two dimmed asides, one closing band. The book
% names three engineering challenges of the same kind, size, speed and
% visibility, so one box each rather than a bullet list, levelled by the equal
% height group. Under the row sit the book's own two one-liners, the joke about
% how fast this ages and the autocompletion example that says what fast enough
% means. The band breaks the row on purpose: it is the last thing the chapter
% says, so the good news gets its own box and keeps the BERT numbers that were
% once called impractical.
\begin{frame}{Size, speed and visibility}
  \footnotesize
  \begin{columns}[T,onlytextwidth]
    \begin{column}{0.31\textwidth}
      \begin{tdblockt}[equal height group=ssv]{Size}
        \scriptsize
        As of 2022 it is common for models to have \alertbf{hundreds of millions, if not billions, of parameters}, needing gigabytes of RAM.
      \end{tdblockt}
    \end{column}
    \begin{column}{0.31\textwidth}
      \begin{tdblockt}[equal height group=ssv]{Speed}
        \scriptsize
        Getting large models into production, \alertbf{especially on edge devices}, is a massive challenge, as is running them fast enough.
      \end{tdblockt}
    \end{column}
    \begin{column}{0.31\textwidth}
      \begin{tdblockt}[equal height group=ssv]{Visibility}
        \scriptsize
        Monitoring and debugging in production is also nontrivial: models get complex, with a \alertbf{lack of visibility} into their work.
      \end{tdblockt}
    \end{column}
  \end{columns}
  \vskip0.3em
  {\color{tddim}\scriptsize
    An autocompletion model is useless if the time it takes to suggest the next character is longer than the time it takes for you to type.}
  \vskip0.3em
  \begin{tdblockt}{The good news: these challenges are being tackled at a breakneck pace}
    \scriptsize
    In 2018 people said BERT was too big, too complex and too slow to be practical. The pretrained large BERT model has \alertbf{340 million} parameters and is \alertbf{1.35 GB}. Two years later, BERT and its variants were used in almost every English search on Google.
  \end{tdblockt}
  \slidesources{Huyen2022}
\end{frame}


% ----------------------------------------------------------------------
% This chapter's own references. The refsection opened above scopes the
% citation numbering to this chapter, so chapter_pdfs/<this>.pdf resolves
% its own [n] markers instead of pointing at a list it does not contain.
\begin{frame}{References}
  \printbibliography[heading=none]
\end{frame}
\end{refsection}

%% ch02_systems_design.tex -- Intelligence Engineering, chapter 02
%% "Introduction to Machine Learning Systems Design".
%% Spine transferred from Huyen, Designing Machine Learning Systems, chapter
%% 2. The subchapters and the frame titles under them are the book's own
%% section and subsection headings, in the book's order.
%%
%% STATE: titles and TEMPLATE layout only. No body text. Every frame carries
%% the layout it is meant to be filled into, and a % TEMPLATE comment saying
%% why that layout and what belongs in each slot. Filling them is the next pass.
%%
%% Fragment of intelligence_engineering.tex. One chapter is one lecture and
%% gets exactly ONE \lecturepage, at the top. Inside it every \subsection
%% opens on the subchapter divider, which the master emits automatically via
%% \AtBeginSubsection; do not write those divider frames by hand.

\begin{refsection}

\lecturepage{II}{Introduction to Machine Learning Systems Design}{}
\section[Systems Design]{Introduction to Machine Learning Systems Design}

% ----------------------------------------------------------------------
\subsection{Business and ML Objectives}

% TEMPLATE: two block contrast. The heading names both sides, so each side
% takes one box rather than one shared list, levelled by the equal height
% group. Left box takes the metrics the business actually reports, right box
% the metrics a model is trained to move. The filling pass has to make the
% mapping from the right box to the left one explicit.
\begin{frame}{Business and ML Objectives}
  \begin{columns}[T,onlytextwidth]
    \begin{column}{0.47\textwidth}
      \begin{tdblockt}[equal height group=obj]{Business objectives}
      \end{tdblockt}
    \end{column}
    \begin{column}{0.47\textwidth}
      \begin{tdblockt}[equal height group=obj]{ML objectives}
      \end{tdblockt}
    \end{column}
  \end{columns}
  \slidesources{Huyen2022}
\end{frame}

% ----------------------------------------------------------------------
\subsection{Requirements for ML Systems}

% TEMPLATE: withmargin plus tddef. This is the first of the book's four
% requirements and the book states it as a definition, so it takes the
% slide's one titled box. Margin holds the aside that makes the requirement
% hard for a machine learning system: it can fail silently and still return
% a prediction that looks plausible.
\begin{frame}{Reliability}
  \begin{withmargin}
    \begin{tdmain}
      \begin{tddef}{Reliability}
      \end{tddef}
    \end{tdmain}
    \begin{tdmargin}
    \end{tdmargin}
  \end{withmargin}
  \slidesources{Huyen2022}
\end{frame}

% TEMPLATE: three plain boxes in a row. Scalability is not one thing in the
% book, it is growth along several axes at once, so each axis takes one box
% and the row makes the count visible instead of burying it in a list. The
% boxes stay unheaded until the filling pass names the axes; the equal
% height group keeps the row level.
\begin{frame}{Scalability}
  \begin{columns}[T,onlytextwidth]
    \begin{column}{0.31\textwidth}
      \begin{tdblock}[equal height group=scale]
      \end{tdblock}
    \end{column}
    \begin{column}{0.31\textwidth}
      \begin{tdblock}[equal height group=scale]
      \end{tdblock}
    \end{column}
    \begin{column}{0.31\textwidth}
      \begin{tdblock}[equal height group=scale]
      \end{tdblock}
    \end{column}
  \end{columns}
  \slidesources{Huyen2022}
\end{frame}

% TEMPLATE: withmargin, and a deliberate break from the row above it. Main
% column takes what maintainability asks of the artifacts themselves, code,
% data and documentation. Margin takes who those artifacts are handed to,
% the book's point that an ML system is built by people who do not share a
% vocabulary.
\begin{frame}{Maintainability}
  \begin{withmargin}
    \begin{tdmain}
    \end{tdmain}
    \begin{tdmargin}
    \end{tdmargin}
  \end{withmargin}
  \slidesources{Huyen2022}
\end{frame}

% TEMPLATE: tddef alone. The book states this requirement as a definition
% too, and the frame reads stronger bare than boxed inside a margin layout,
% so the term carries the whole slide. The box takes the capacities the book
% asks for and nothing else; what they imply for the workflow is the next
% subchapter's job.
\begin{frame}{Adaptability}
  \begin{tddef}{Adaptability}
  \end{tddef}
  \slidesources{Huyen2022}
\end{frame}

% ----------------------------------------------------------------------
\subsection{Iterative Process}

% TEMPLATE: side image hero, the first of this chapter's two figures. The
% heading is a process, so the book's own cycle owns the frame and the grey
% placeholder holds its place until the diagram is drawn. sidebody takes the
% steps in the book's order, the phimgside caption the one line that says why the
% arrows run backwards as often as forwards.
\begin{frame}{Iterative Process}
  \phimgside[0.33]{iterative_process.png}{}
  \begin{sidebody}
  \end{sidebody}
  \slidesources{Huyen2022}
\end{frame}

% ----------------------------------------------------------------------
\subsection{Framing ML Problems}

% TEMPLATE: two by two grid of plain boxes, the opener for the four frames
% below it. The heading introduces a countable set, so the grid is sized to
% that count and each box previews one of the four headings that follow
% rather than arguing anything itself. Unheaded on purpose: the detail
% frames own the names.
\begin{frame}{Types of ML Tasks}
  \begin{columns}[T,onlytextwidth]
    \begin{column}{0.47\textwidth}
      \begin{tdblock}[equal height group=tasks]
      \end{tdblock}
      \vskip0.5em
      \begin{tdblock}[equal height group=tasks]
      \end{tdblock}
    \end{column}
    \begin{column}{0.47\textwidth}
      \begin{tdblock}[equal height group=tasks]
      \end{tdblock}
      \vskip0.5em
      \begin{tdblock}[equal height group=tasks]
      \end{tdblock}
    \end{column}
  \end{columns}
  \slidesources{Huyen2022}
\end{frame}

% TEMPLATE: two block contrast. The heading is a versus, so each task takes
% one box with its header read straight off the heading, levelled by the
% equal height group. Each box takes what that task predicts in the book's
% words, and the filling pass owes the book's reminder that either one can
% be turned into the other.
\begin{frame}{Classification versus regression}
  \begin{columns}[T,onlytextwidth]
    \begin{column}{0.47\textwidth}
      \begin{tdblockt}[equal height group=clsreg]{Classification}
      \end{tdblockt}
    \end{column}
    \begin{column}{0.47\textwidth}
      \begin{tdblockt}[equal height group=clsreg]{Regression}
      \end{tdblockt}
    \end{column}
  \end{columns}
  \slidesources{Huyen2022}
\end{frame}

% TEMPLATE: comparison table. A second pair of boxes in a row would stall,
% and the book separates these two along named dimensions rather than in
% prose, so the columns are the two sides from the heading and the empty
% rows are for those dimensions, one per row. booktabs only, no vertical
% rules.
\begin{frame}{Binary versus multiclass classification}
  \footnotesize
  \renewcommand{\arraystretch}{1.35}
  \begin{tabular}{@{}>{\raggedright\arraybackslash}p{0.24\textwidth} >{\raggedright\arraybackslash}p{0.32\textwidth} >{\raggedright\arraybackslash}p{0.32\textwidth}@{}}
    \toprule
     & \textbf{Binary} & \textbf{Multiclass}\\
    \midrule
     & & \\
     & & \\
     & & \\
     & & \\
    \bottomrule
  \end{tabular}
  \slidesources{Huyen2022}
\end{frame}

% TEMPLATE: two code panels. The book settles this distinction on how a
% label is encoded, so the contrast belongs in the mono panels where a label
% vector can be shown literally instead of described. The captions are the
% two sides from the heading; each panel takes the encoding the book prints
% for that case, and the ambiguity it creates.
\begin{frame}{Multiclass versus multilabel classification}
  \begin{columns}[T,onlytextwidth]
    \begin{column}{0.47\textwidth}
      \begin{tdcodet}{Multiclass}
      \end{tdcodet}
    \end{column}
    \begin{column}{0.47\textwidth}
      \begin{tdcodet}{Multilabel}
      \end{tdcodet}
    \end{column}
  \end{columns}
  \slidesources{Huyen2022}
\end{frame}

% TEMPLATE: two plain boxes. The heading promises more than one framing of
% the same problem, so the framings sit side by side where they can be read
% off against each other. Unheaded, because the heading names no framing;
% each box takes one framing with its inputs, its outputs and what happens
% to it when the set of classes grows.
\begin{frame}{Multiple ways to frame a problem}
  \begin{columns}[T,onlytextwidth]
    \begin{column}{0.47\textwidth}
      \begin{tdblock}[equal height group=ways]
      \end{tdblock}
    \end{column}
    \begin{column}{0.47\textwidth}
      \begin{tdblock}[equal height group=ways]
      \end{tdblock}
    \end{column}
  \end{columns}
  \slidesources{Huyen2022}
\end{frame}

% TEMPLATE: withmargin, the opener for the frame below it. Main column takes
% what an objective function is for and the losses the book names, margin
% the aside that the objective is a choice and not a given. Kept general on
% purpose: the detail, two objectives inside one loss, is the next frame.
\begin{frame}{Objective Functions}
  \begin{withmargin}
    \begin{tdmain}
    \end{tdmain}
    \begin{tdmargin}
    \end{tdmargin}
  \end{withmargin}
  \slidesources{Huyen2022}
\end{frame}

% TEMPLATE: side image hero, the chapter's second and last figure. The
% heading is an architecture decision, one model per objective with the
% predictions combined afterwards, which reads as a diagram and not as
% prose. sidebody takes the argument for decoupling in the book's order,
% the phimgside caption what the decoupling costs.
\begin{frame}{Decoupling objectives}
  \phimgside[0.33]{decoupling_objectives.png}{}
  \begin{sidebody}
  \end{sidebody}
  \slidesources{Huyen2022}
\end{frame}

% ----------------------------------------------------------------------
\subsection{Mind Versus Data}

% TEMPLATE: two block contrast. This is the chapter's loudest versus and the
% book gives it as a standing argument with named people on both sides, so
% each camp takes one box, headed from the heading and levelled by the equal
% height group. Each box takes that camp's position and the evidence it
% leans on.
\begin{frame}{Mind Versus Data}
  \begin{columns}[T,onlytextwidth]
    \begin{column}{0.47\textwidth}
      \begin{tdblockt}[equal height group=mind]{Mind}
      \end{tdblockt}
    \end{column}
    \begin{column}{0.47\textwidth}
      \begin{tdblockt}[equal height group=mind]{Data}
      \end{tdblockt}
    \end{column}
  \end{columns}
  \slidesources{Huyen2022}
\end{frame}

% ----------------------------------------------------------------------
\subsection{Summary}

% TEMPLATE: takeaway alone. A summary earns no scaffold of its own, so the
% chapter closes on one centred line with nothing competing against it. The
% line is the single sentence the chapter should leave behind, lifted from
% the book's own summary rather than written fresh.
\begin{frame}{Summary}
  \begin{takeaway}
  \end{takeaway}
  \slidesources{Huyen2022}
\end{frame}

% ----------------------------------------------------------------------
% This chapter's own references. The refsection opened above scopes the
% citation numbering to this chapter, so chapter_pdfs/<this>.pdf resolves
% its own [n] markers instead of pointing at a list it does not contain.
\begin{frame}{References}
  \printbibliography[heading=none]
\end{frame}
\end{refsection}

%%% ch03_data_engineering.tex -- Intelligence Engineering, chapter 03 "Data
%% Engineering Fundamentals".
%% Spine transferred from Huyen, Designing Machine Learning Systems, chapter
%% 3. The subchapters and the frame titles under them are the book's own
%% section and subsection headings, in the book's order.
%%
%% STATE: titles and TEMPLATE layout only. No body text. Every frame carries
%% the layout it is meant to be filled into, and a % TEMPLATE comment saying
%% why that layout and what belongs in each slot. Filling them is the next pass.
%%
%% Fragment of intelligence_engineering.tex. One chapter is one lecture and
%% gets exactly ONE \lecturepage, at the top. Inside it every \subsection
%% opens on the subchapter divider, which the master emits automatically via
%% \AtBeginSubsection; do not write those divider frames by hand.

\begin{refsection}

\lecturepage{III}{Data Engineering Fundamentals}{}
\section[Data Engineering]{Data Engineering Fundamentals}

% ----------------------------------------------------------------------
\subsection{Data Sources}

% TEMPLATE: two by two block grid. The heading is plural and the book answers
% it with four peer sources, user input data, system generated data, internal
% databases and third party data, so one box each rather than a bullet list,
% levelled by the equal height groups. The boxes stay untitled because the
% heading names no source; the later pass opens each box with its source name
% in bold, and the book's warning about malformed user input goes in the first.
\begin{frame}{Data Sources}
  \begin{columns}[T,onlytextwidth]
    \begin{column}{0.47\textwidth}
      \begin{tdblock}[equal height group=src]
      \end{tdblock}
      \vskip0.5em
      \begin{tdblock}[equal height group=src2]
      \end{tdblock}
    \end{column}
    \begin{column}{0.47\textwidth}
      \begin{tdblock}[equal height group=src]
      \end{tdblock}
      \vskip0.5em
      \begin{tdblock}[equal height group=src2]
      \end{tdblock}
    \end{column}
  \end{columns}
  \slidesources{Huyen2022}
\end{frame}

% ----------------------------------------------------------------------
\subsection{Data Formats}

% TEMPLATE: withmargin carrying a code panel. The heading is a serialization
% format, so the main column shows the book's own JSON object instead of
% describing it, captioned with the format the heading names. Margin takes the
% cost the book draws from it, the schema that is committed to the moment the
% data is written, and the pain of changing it afterwards.
\begin{frame}{JSON}
  \begin{withmargin}
    \begin{tdmain}
      \begin{tdcodet}{JSON}
      \end{tdcodet}
    \end{tdmain}
    \begin{tdmargin}
    \end{tdmargin}
  \end{withmargin}
  \slidesources{Huyen2022}
\end{frame}

% TEMPLATE: two block contrast. The heading is a versus, so each format takes
% one box rather than a bullet list, levelled by the equal height group, with
% the headers taken from the two sides the heading names. Each box holds what
% that layout makes cheap and what it makes expensive in the book's words, and
% the numpy against pandas timing example belongs in the column it favours.
\begin{frame}{Row-Major Versus Column-Major Format}
  \begin{columns}[T,onlytextwidth]
    \begin{column}{0.47\textwidth}
      \begin{tdblockt}[equal height group=rowcol]{Row-major}
      \end{tdblockt}
    \end{column}
    \begin{column}{0.47\textwidth}
      \begin{tdblockt}[equal height group=rowcol]{Column-major}
      \end{tdblockt}
    \end{column}
  \end{columns}
  \slidesources{Huyen2022}
\end{frame}

% TEMPLATE: two code panels side by side. A versus again, but the frame before
% it already spends the block contrast, and the point here is what the bytes
% look like, so each side takes a dark mono panel captioned with the side the
% heading names. Text panel takes the book's CSV row, binary panel the same row
% as it reads on disk, and the AWS saving for CSV against Parquet goes with it.
\begin{frame}{Text Versus Binary Format}
  \begin{columns}[T,onlytextwidth]
    \begin{column}{0.47\textwidth}
      \begin{tdcodet}{Text}
      \end{tdcodet}
    \end{column}
    \begin{column}{0.47\textwidth}
      \begin{tdcodet}{Binary}
      \end{tdcodet}
    \end{column}
  \end{columns}
  \slidesources{Huyen2022}
\end{frame}

% ----------------------------------------------------------------------
\subsection{Data Models}

% TEMPLATE: definition plus a code panel. The heading names a data model, so
% the term takes the titled definition box and the relational vocabulary,
% relation, tuple and normalization, lands inside it. Under it the mono panel
% takes the book's SQL query, which is on the slide to make one point: the
% query says what the result should be, not how to compute it.
\begin{frame}{Relational Model}
  \begin{tddef}{Relational model}
  \end{tddef}
  \vskip0.5em
  \begin{tdcode}
  \end{tdcode}
  \slidesources{Huyen2022}
\end{frame}

% TEMPLATE: two block contrast used as an opener. This heading owns the two
% deeper headings that follow, so its own frame previews them by name instead
% of going into either, levelled by the equal height group. Each box takes one
% child's single line claim; the book's reading of the acronym, Not Only SQL
% against nonrelational, belongs in the line above the row.
\begin{frame}{NoSQL}
  \begin{columns}[T,onlytextwidth]
    \begin{column}{0.47\textwidth}
      \begin{tdblockt}[equal height group=nosql]{Document model}
      \end{tdblockt}
    \end{column}
    \begin{column}{0.47\textwidth}
      \begin{tdblockt}[equal height group=nosql]{Graph model}
      \end{tdblockt}
    \end{column}
  \end{columns}
  \slidesources{Huyen2022}
\end{frame}

% TEMPLATE: withmargin, the chapter's workhorse, one claim with one aside. The
% main column takes the book's document and collection, and the consequence it
% draws, that the schema moves to the reader. Margin takes the price of that,
% the join the document model cannot do for you.
\begin{frame}{Document model}
  \begin{withmargin}
    \begin{tdmain}
    \end{tdmain}
    \begin{tdmargin}
    \end{tdmargin}
  \end{withmargin}
  \slidesources{Huyen2022}
\end{frame}

% TEMPLATE: full height side image. The book's content here is a figure, the
% graph of people and the countries they were born in, so the strip carries
% that diagram and the body keeps to nodes, edges and the query that walks an
% edge. First of this chapter's two image frames; the caption takes the
% figure's own one line reading, not a repeat of the frame title.
\begin{frame}{Graph model}
  \phimgside[0.33]{graph_model.png}{}
  \begin{sidebody}
  \end{sidebody}
  \slidesources{Huyen2022}
\end{frame}

% TEMPLATE: comparison table. The book already gives this contrast as a table,
% two columns of paired statements with no row labels, so it stays a table
% rather than a third block pair. Column headers are the two sides the heading
% names; each row takes one of the book's pairs, the defined schema against the
% schema on read, the warehouse against the lake.
\begin{frame}{Structured Versus Unstructured Data}
  \footnotesize
  \renewcommand{\arraystretch}{1.25}
  \begin{tabular}{@{}%
    >{\raggedright\arraybackslash}p{0.44\textwidth}%
    >{\raggedright\arraybackslash}p{0.44\textwidth}@{}}
    \toprule
    \textbf{Structured data} & \textbf{Unstructured data}\\
    \midrule
     & \\
     & \\
     & \\
     & \\
    \bottomrule
  \end{tabular}
  \slidesources{Huyen2022}
\end{frame}

% ----------------------------------------------------------------------
\subsection{Data Storage Engines and Processing}

% TEMPLATE: two block contrast. The heading names both processing types, so
% each takes one box, levelled by the equal height group. Transactional box
% holds the ACID properties and the latency each transaction is held to,
% analytical box the aggregated queries over columns. The line under the row is
% the book's own correction, that storage and processing have come apart.
\begin{frame}{Transactional and Analytical Processing}
  \begin{columns}[T,onlytextwidth]
    \begin{column}{0.47\textwidth}
      \begin{tdblockt}[equal height group=tap]{Transactional processing}
      \end{tdblockt}
    \end{column}
    \begin{column}{0.47\textwidth}
      \begin{tdblockt}[equal height group=tap]{Analytical processing}
      \end{tdblockt}
    \end{column}
  \end{columns}
  \slidesources{Huyen2022}
\end{frame}

% TEMPLATE: three block row. The heading itself names three stages in order, so
% the box headers are the heading, one stage per third of the width, levelled
% by the equal height group. Each box takes what that stage does in the book's
% words, with the transform box carrying the bulk of the work. ELT and the data
% lake belong in the line under the row, as the reversal of this order.
\begin{frame}{ETL: Extract, Transform, and Load}
  \begin{columns}[T,onlytextwidth]
    \begin{column}{0.31\textwidth}
      \begin{tdblockt}[equal height group=etl]{Extract}
      \end{tdblockt}
    \end{column}
    \begin{column}{0.31\textwidth}
      \begin{tdblockt}[equal height group=etl]{Transform}
      \end{tdblockt}
    \end{column}
    \begin{column}{0.31\textwidth}
      \begin{tdblockt}[equal height group=etl]{Load}
      \end{tdblockt}
    \end{column}
  \end{columns}
  \slidesources{Huyen2022}
\end{frame}

% ----------------------------------------------------------------------
\subsection{Modes of Dataflow}

% TEMPLATE: withmargin. One mode, one consequence, so the main column takes the
% mechanism, one process writes to a shared table and the other reads it, then
% the two reasons the book rules it out. Margin takes the access condition, both
% processes on the same database, and the read latency that makes it unusable
% for anything a user waits on.
\begin{frame}{Data Passing Through Databases}
  \begin{withmargin}
    \begin{tdmain}
    \end{tdmain}
    \begin{tdmargin}
    \end{tdmargin}
  \end{withmargin}
  \slidesources{Huyen2022}
\end{frame}

% TEMPLATE: full height side image. The book draws this mode as an
% architecture, two services sending requests to each other and the price
% service in the middle, so the strip carries that diagram and the body keeps
% to the request driven style and its two implementations, REST and RPC.
% Second and last image frame of the chapter; the rest goes into words.
\begin{frame}{Data Passing Through Services}
  \phimgside[0.33]{data_passing_through_services.png}{}
  \begin{sidebody}
  \end{sidebody}
  \slidesources{Huyen2022}
\end{frame}

% TEMPLATE: withmargin plus a takeaway. The image budget is spent, so the
% broker goes into words: main column takes the event driven mechanism,
% producers and consumers over a shared topic, with Kafka and Kinesis named as
% the log based brokers. Margin holds what the broker buys, the connections
% between services that fall away. The takeaway is the book's own verdict.
\begin{frame}{Data Passing Through Real-Time Transport}
  \begin{withmargin}
    \begin{tdmain}
    \end{tdmain}
    \begin{tdmargin}
    \end{tdmargin}
  \end{withmargin}
  \begin{takeaway}
  \end{takeaway}
  \slidesources{Huyen2022}
\end{frame}

% ----------------------------------------------------------------------
\subsection{Batch Processing Versus Stream Processing}

% TEMPLATE: two block contrast plus a takeaway. A versus, and the chapter turns
% on this one, so each side keeps its own box, levelled by the equal height
% group, headers from the two sides the heading names. Batch box takes bounded
% data and the static features, stream box unbounded data and the features that
% only exist in the window. Takeaway is the book's claim about the special case.
\begin{frame}{Batch Processing Versus Stream Processing}
  \begin{columns}[T,onlytextwidth]
    \begin{column}{0.47\textwidth}
      \begin{tdblockt}[equal height group=bstr]{Batch processing}
      \end{tdblockt}
    \end{column}
    \begin{column}{0.47\textwidth}
      \begin{tdblockt}[equal height group=bstr]{Stream processing}
      \end{tdblockt}
    \end{column}
  \end{columns}
  \begin{takeaway}
  \end{takeaway}
  \slidesources{Huyen2022}
\end{frame}

% ----------------------------------------------------------------------
\subsection{Summary}

% TEMPLATE: a bare takeaway, nothing else on the frame. The heading is the
% chapter summary, so it gets one centred line: the decision the whole chapter
% has been circling, how the data is stored and how it is passed between
% processes, and not a recap of everything said.
\begin{frame}{Summary}
  \begin{takeaway}
  \end{takeaway}
  \slidesources{Huyen2022}
\end{frame}

% ----------------------------------------------------------------------
% This chapter's own references. The refsection opened above scopes the
% citation numbering to this chapter, so chapter_pdfs/<this>.pdf resolves
% its own [n] markers instead of pointing at a list it does not contain.
\begin{frame}{References}
  \printbibliography[heading=none]
\end{frame}
\end{refsection}

%%% ch04_training_data.tex -- Intelligence Engineering, chapter 04 "Training
%% Data".
%% Spine transferred from Huyen, Designing Machine Learning Systems, chapter
%% 4. The subchapters and the frame titles under them are the book's own
%% section and subsection headings, in the book's order.
%%
%% STATE: titles and TEMPLATE layout only. No body text. Every frame carries
%% the layout it is meant to be filled into, and a % TEMPLATE comment saying
%% why that layout and what belongs in each slot. Filling them is the next pass.
%%
%% Fragment of intelligence_engineering.tex. One chapter is one lecture and
%% gets exactly ONE \lecturepage, at the top. Inside it every \subsection
%% opens on the subchapter divider, which the master emits automatically via
%% \AtBeginSubsection; do not write those divider frames by hand.

\begin{refsection}

\lecturepage{IV}{Training Data}{}
\section[Training Data]{Training Data}

% ----------------------------------------------------------------------
\subsection{Sampling}

% TEMPLATE: withmargin carrying a definition box. The heading is the term the
% first half of this subchapter turns on, so it takes the slide's one titled
% box and the book's kinds of nonprobability sampling are listed inside it.
% Margin takes the aside the book uses to justify opening here, the convenience
% of these samples against how unrepresentative they are.
\begin{frame}{Nonprobability Sampling}
  \begin{withmargin}
    \begin{tdmain}
      \begin{tddef}{Nonprobability sampling}
      \end{tddef}
    \end{tdmain}
    \begin{tdmargin}
    \end{tdmargin}
  \end{withmargin}
  \slidesources{Huyen2022}
\end{frame}

% TEMPLATE: two levelled boxes, untitled on purpose. The book gives this method
% one selection rule and one failure mode, which read as two equal halves
% rather than as a claim with an aside. The left box takes the rule and the
% probability it assigns each sample, the right box what it does to a rare
% class; the later pass opens each box with its own bold label.
\begin{frame}{Simple Random Sampling}
  \begin{columns}[T,onlytextwidth]
    \begin{column}{0.47\textwidth}
      \begin{tdblock}[equal height group=srs]
      \end{tdblock}
    \end{column}
    \begin{column}{0.47\textwidth}
      \begin{tdblock}[equal height group=srs]
      \end{tdblock}
    \end{column}
  \end{columns}
  \slidesources{Huyen2022}
\end{frame}

% TEMPLATE: withmargin. One mechanism plus one caveat, which is exactly the
% shape this layout is for. Main column takes how the population is divided
% into groups and sampled inside each of them, with the book's own class
% proportions. Margin takes the limit the book immediately raises, that the
% groups cannot always be drawn.
\begin{frame}{Stratified Sampling}
  \begin{withmargin}
    \begin{tdmain}
    \end{tdmain}
    \begin{tdmargin}
    \end{tdmargin}
  \end{withmargin}
  \slidesources{Huyen2022}
\end{frame}

% TEMPLATE: code panel captioned with the heading. The book makes this method
% concrete with a weights vector and a single call rather than with prose, and
% a weights vector only reads as itself in mono, so the panel is the whole
% slide. It takes that snippet, the weights it passes and the distribution they
% encode.
\begin{frame}{Weighted Sampling}
  \begin{tdcodet}{Weighted sampling}
  \end{tdcodet}
  \slidesources{Huyen2022}
\end{frame}

% TEMPLATE: three levelled boxes. The book states this one as an algorithm of
% three moves over a stream, and the moves only make sense in sequence, so one
% box each rather than a paragraph. The equal height group keeps the row level;
% the boxes stay untitled until the later pass numbers them and adds the
% probability that makes the reservoir uniform.
\begin{frame}{Reservoir Sampling}
  \begin{columns}[T,onlytextwidth]
    \begin{column}{0.31\textwidth}
      \begin{tdblock}[equal height group=resv]
      \end{tdblock}
    \end{column}
    \begin{column}{0.31\textwidth}
      \begin{tdblock}[equal height group=resv]
      \end{tdblock}
    \end{column}
    \begin{column}{0.31\textwidth}
      \begin{tdblock}[equal height group=resv]
      \end{tdblock}
    \end{column}
  \end{columns}
  \slidesources{Huyen2022}
\end{frame}

% TEMPLATE: withmargin. The book's case rests on one ratio between two
% distributions, so the main column carries that expression and the condition
% on the proposal distribution. Margin takes the use the book reaches for,
% where the samples are expensive and the distribution you want is not the one
% you can draw from.
\begin{frame}{Importance Sampling}
  \begin{withmargin}
    \begin{tdmain}
    \end{tdmain}
    \begin{tdmargin}
    \end{tdmargin}
  \end{withmargin}
  \slidesources{Huyen2022}
\end{frame}

% ----------------------------------------------------------------------
\subsection{Labeling}

% TEMPLATE: three levelled boxes, and the opener for the two frames under this
% heading. The book's case against hand labels is three costs of the same kind,
% peers of one another, so one box each rather than a bullet list. The two
% named problems below get their own frames and are deliberately not repeated
% here; the boxes stay untitled until the later pass names each cost.
\begin{frame}{Hand Labels}
  \begin{columns}[T,onlytextwidth]
    \begin{column}{0.31\textwidth}
      \begin{tdblock}[equal height group=hand]
      \end{tdblock}
    \end{column}
    \begin{column}{0.31\textwidth}
      \begin{tdblock}[equal height group=hand]
      \end{tdblock}
    \end{column}
    \begin{column}{0.31\textwidth}
      \begin{tdblock}[equal height group=hand]
      \end{tdblock}
    \end{column}
  \end{columns}
  \slidesources{Huyen2022}
\end{frame}

% TEMPLATE: withmargin. The book argues this one from a single annotated
% sentence and the disagreement it produces, so the main column takes that
% example and the differing counts it yields. Margin takes the book's remedy,
% the problem definition and the instructions that have to exist before the
% annotators start.
\begin{frame}{Label multiplicity}
  \begin{withmargin}
    \begin{tdmain}
    \end{tdmain}
    \begin{tdmargin}
    \end{tdmargin}
  \end{withmargin}
  \slidesources{Huyen2022}
\end{frame}

% TEMPLATE: definition box alone, full width. The heading is the book's own
% term and there is nothing here to contrast it against, so the definition
% takes the whole frame instead of sharing it. The box takes what has to be
% tracked per sample and the failure the book reports when it is not, a model
% that got worse after new data arrived.
\begin{frame}{Data lineage}
  \begin{tddef}{Data lineage}
  \end{tddef}
  \slidesources{Huyen2022}
\end{frame}

% TEMPLATE: booktabs table, pushed there on purpose rather than left as a
% withmargin. The book names its canonical cases in prose, but every one of
% them is the same two fields, the task and the label the system gets for free,
% so a table is the tighter transfer and makes the pattern visible. One row per
% case, column headers taken from the heading.
\begin{frame}{Natural Labels}
  \footnotesize
  \renewcommand{\arraystretch}{1.25}
  \begin{tabular}{@{}>{\raggedright\arraybackslash}p{0.38\textwidth}>{\raggedright\arraybackslash}p{0.50\textwidth}@{}}
    \toprule
    \textbf{Task} & \textbf{Natural label}\\
    \midrule
     & \\
     & \\
     & \\
     & \\
    \bottomrule
  \end{tabular}
  \slidesources{Huyen2022}
\end{frame}

% TEMPLATE: two levelled boxes. The heading is about a length, and a length
% only means something as a short end read against a long end, so each end
% takes one box headed by it. Each box takes the book's own examples and the
% window the label closes in, which is the number the deployment decision
% later hangs on.
\begin{frame}{Feedback loop length}
  \begin{columns}[T,onlytextwidth]
    \begin{column}{0.47\textwidth}
      \begin{tdblockt}[equal height group=fbl]{Short feedback loops}
      \end{tdblockt}
    \end{column}
    \begin{column}{0.47\textwidth}
      \begin{tdblockt}[equal height group=fbl]{Long feedback loops}
      \end{tdblockt}
    \end{column}
  \end{columns}
  \slidesources{Huyen2022}
\end{frame}

% TEMPLATE: two by two block grid, and the opener for the four frames under
% this heading. The heading introduces exactly four peer methods, none
% subordinate to another, so each takes one box headed with its own name rather
% than one long list. Each box carries only the one line the book's comparison
% uses to tell that method apart, how it gets its labels and whether it needs
% ground truths; the detail waits for the frames below.
\begin{frame}{Handling the Lack of Labels}
  \begin{columns}[T,onlytextwidth]
    \begin{column}{0.47\textwidth}
      \begin{tdblockt}[equal height group=lack]{Weak supervision}
      \end{tdblockt}
      \vskip0.5em
      \begin{tdblockt}[equal height group=lack2]{Transfer learning}
      \end{tdblockt}
    \end{column}
    \begin{column}{0.47\textwidth}
      \begin{tdblockt}[equal height group=lack]{Semi-supervision}
      \end{tdblockt}
      \vskip0.5em
      \begin{tdblockt}[equal height group=lack2]{Active learning}
      \end{tdblockt}
    \end{column}
  \end{columns}
  \slidesources{Huyen2022}
\end{frame}

% TEMPLATE: side image cover, the first of this chapter's two image slots. The
% book's own content here is a pipeline, heuristics written as functions that
% vote and are combined into labels, and a pipeline is a figure before it is a
% sentence, so the strip holds it. The body holds the definition of a labeling
% function and the kinds of heuristic the book lists.
\begin{frame}{Weak supervision}
  \phimgside[0.33]{weak_supervision.png}{Weak supervision}
  \begin{sidebody}
  \end{sidebody}
  \slidesources{Huyen2022}
\end{frame}

% TEMPLATE: three levelled boxes. The book names three ways to grow labels from
% a small seed set, peers of one another, so one box each rather than a bullet
% list, levelled by the equal height group. The boxes stay untitled until the
% later pass opens each with its method name and the assumption it rests on.
\begin{frame}{Semi-supervision}
  \begin{columns}[T,onlytextwidth]
    \begin{column}{0.31\textwidth}
      \begin{tdblock}[equal height group=semi]
      \end{tdblock}
    \end{column}
    \begin{column}{0.31\textwidth}
      \begin{tdblock}[equal height group=semi]
      \end{tdblock}
    \end{column}
    \begin{column}{0.31\textwidth}
      \begin{tdblock}[equal height group=semi]
      \end{tdblock}
    \end{column}
  \end{columns}
  \slidesources{Huyen2022}
\end{frame}

% TEMPLATE: withmargin. One mechanism plus one aside. Main column takes the
% book's order, the base task and its data, then the fine tuning step, then the
% zero shot case where no fine tuning happens at all. Margin takes what the
% book says this did to the cost of entry for tasks with little data.
\begin{frame}{Transfer learning}
  \begin{withmargin}
    \begin{tdmain}
    \end{tdmain}
    \begin{tdmargin}
    \end{tdmargin}
  \end{withmargin}
  \slidesources{Huyen2022}
\end{frame}

% TEMPLATE: two levelled boxes. The book's heuristics for choosing which sample
% to label next fall into two families, so they read side by side rather than
% as one list. Each box takes one family, the quantity it ranks candidates by
% and the book's query strategy for it; the later pass heads each box with the
% family name.
\begin{frame}{Active learning}
  \begin{columns}[T,onlytextwidth]
    \begin{column}{0.47\textwidth}
      \begin{tdblock}[equal height group=actv]
      \end{tdblock}
    \end{column}
    \begin{column}{0.47\textwidth}
      \begin{tdblock}[equal height group=actv]
      \end{tdblock}
    \end{column}
  \end{columns}
  \slidesources{Huyen2022}
\end{frame}

% ----------------------------------------------------------------------
\subsection{Class Imbalance}

% TEMPLATE: withmargin plus a takeaway. The book argues from one ratio and then
% from three consequences of it, so the main column runs them in that order and
% the margin carries the worked case it uses, the scan set and its counts. The
% takeaway keeps the book's punch in front of the claim, the model that scores
% well and never predicts the rare class.
\begin{frame}{Challenges of Class Imbalance}
  \begin{withmargin}
    \begin{tdmain}
    \end{tdmain}
    \begin{tdmargin}
    \end{tdmargin}
  \end{withmargin}
  \begin{takeaway}
  \end{takeaway}
  \slidesources{Huyen2022}
\end{frame}

% TEMPLATE: three levelled boxes, and the opener for the three frames under
% this heading. The heading introduces three peer families of fix and the book
% keeps them strictly apart, so one box each headed with its own name. Each box
% holds only the one line that tells that family from the other two, the level
% it acts on; the detail waits for the frames below.
\begin{frame}{Handling Class Imbalance}
  \begin{columns}[T,onlytextwidth]
    \begin{column}{0.31\textwidth}
      \begin{tdblockt}[equal height group=hci]{Evaluation metrics}
      \end{tdblockt}
    \end{column}
    \begin{column}{0.31\textwidth}
      \begin{tdblockt}[equal height group=hci]{Data-level methods}
      \end{tdblockt}
    \end{column}
    \begin{column}{0.31\textwidth}
      \begin{tdblockt}[equal height group=hci]{Algorithm-level methods}
      \end{tdblockt}
    \end{column}
  \end{columns}
  \slidesources{Huyen2022}
\end{frame}

% TEMPLATE: booktabs table with the header row left open. The book makes this
% point with a model by metric grid in which the same two models swap ranks the
% moment the metric changes, so the grid is the argument and a paragraph would
% lose it. The header row takes the book's metric names, the rows its two
% models, and the cells the numbers that reverse.
\begin{frame}{Using the right evaluation metrics}
  \footnotesize
  \renewcommand{\arraystretch}{1.25}
  \begin{tabular}{@{}llll@{}}
    \toprule
     & & & \\
    \midrule
     & & & \\
     & & & \\
    \bottomrule
  \end{tabular}
  \slidesources{Huyen2022}
\end{frame}

% TEMPLATE: two levelled boxes. Resampling runs in two directions and the
% book's named methods split cleanly along them, so one box each rather than
% one list that hides the split. Each box takes its direction, the methods the
% book names for it and what that direction risks, losing data on one side and
% overfitting on the other.
\begin{frame}{Data-level methods: Resampling}
  \begin{columns}[T,onlytextwidth]
    \begin{column}{0.47\textwidth}
      \begin{tdblock}[equal height group=rsmp]
      \end{tdblock}
    \end{column}
    \begin{column}{0.47\textwidth}
      \begin{tdblock}[equal height group=rsmp]
      \end{tdblock}
    \end{column}
  \end{columns}
  \slidesources{Huyen2022}
\end{frame}

% TEMPLATE: withmargin. The book's three modifications are all the same move on
% the loss function, so the main column carries them together with their
% expressions rather than splitting them into boxes that would hide how alike
% they are. Margin carries the shared idea, that a mistake on the rare class
% has to cost more than a mistake on the common one.
\begin{frame}{Algorithm-level methods}
  \begin{withmargin}
    \begin{tdmain}
    \end{tdmain}
    \begin{tdmargin}
    \end{tdmargin}
  \end{withmargin}
  \slidesources{Huyen2022}
\end{frame}

% ----------------------------------------------------------------------
\subsection{Data Augmentation}

% TEMPLATE: two levelled boxes. The book gives this one twice, once for images
% and once for text, the same move in two modalities, so the two modalities sit
% side by side rather than in one list. Each box takes that modality's
% transformations in the book's words and the reason the label survives them.
\begin{frame}{Simple Label-Preserving Transformations}
  \begin{columns}[T,onlytextwidth]
    \begin{column}{0.47\textwidth}
      \begin{tdblock}[equal height group=slpt]
      \end{tdblock}
    \end{column}
    \begin{column}{0.47\textwidth}
      \begin{tdblock}[equal height group=slpt]
      \end{tdblock}
    \end{column}
  \end{columns}
  \slidesources{Huyen2022}
\end{frame}

% TEMPLATE: side image cover, this chapter's second and last image slot. The
% book makes this point with a picture, a sample a human still reads one way
% while the model reads it another after a change too small to see, and that
% only lands as an image. The strip holds it; the body holds the definition and
% the defensive use the book draws from it.
\begin{frame}{Perturbation}
  \phimgside[0.33]{perturbation.png}{Perturbation}
  \begin{sidebody}
  \end{sidebody}
  \slidesources{Huyen2022}
\end{frame}

% TEMPLATE: withmargin carrying a code panel. The book's example here is a
% template with slots that get filled to generate samples, and a template only
% reads as itself in mono, so the panel takes it in the main column. Margin
% takes the boundary the book draws, what synthesis can stand in for and what
% it cannot.
\begin{frame}{Data Synthesis}
  \begin{withmargin}
    \begin{tdmain}
      \begin{tdcode}
      \end{tdcode}
    \end{tdmain}
    \begin{tdmargin}
    \end{tdmargin}
  \end{withmargin}
  \slidesources{Huyen2022}
\end{frame}

% ----------------------------------------------------------------------
\subsection{Summary}

% TEMPLATE: bare takeaway. A summary slide that repeats the chapter as bullets
% is dead weight, and the book's own summary is one claim rather than a recap,
% so this frame carries a single centred line and nothing else, the book's
% closing statement on training data and the work it still demands.
\begin{frame}{Summary}
  \begin{takeaway}
  \end{takeaway}
  \slidesources{Huyen2022}
\end{frame}

% ----------------------------------------------------------------------
% This chapter's own references. The refsection opened above scopes the
% citation numbering to this chapter, so chapter_pdfs/<this>.pdf resolves
% its own [n] markers instead of pointing at a list it does not contain.
\begin{frame}{References}
  \printbibliography[heading=none]
\end{frame}
\end{refsection}

%%% ch05_feature_engineering.tex -- Intelligence Engineering, chapter 05
%% "Feature Engineering".
%% Spine transferred from Huyen, Designing Machine Learning Systems, chapter
%% 5. The subchapters and the frame titles under them are the book's own
%% section and subsection headings, in the book's order.
%%
%% STATE: titles and TEMPLATE layout only. No body text. Every frame carries
%% the layout it is meant to be filled into, and a % TEMPLATE comment saying
%% why that layout and what belongs in each slot. Filling them is the next pass.
%%
%% Fragment of intelligence_engineering.tex. One chapter is one lecture and
%% gets exactly ONE \lecturepage, at the top. Inside it every \subsection
%% opens on the subchapter divider, which the master emits automatically via
%% \AtBeginSubsection; do not write those divider frames by hand.

\begin{refsection}

\lecturepage{V}{Feature Engineering}{}
\section[Feature Engineering]{Feature Engineering}

% ----------------------------------------------------------------------
\subsection{Learned Features Versus Engineered Features}

% TEMPLATE: two block contrast. The heading is a versus, so each side takes
% one box rather than a bullet list, levelled by the equal height group. Left
% box takes what the model learns by itself in the book's words, right box the
% work that stays manual. The line under them is for the consequence the book
% draws, that feature engineering remains the bulk of the effort.
\begin{frame}{Learned Features Versus Engineered Features}
  \begin{columns}[T,onlytextwidth]
    \begin{column}{0.47\textwidth}
      \begin{tdblockt}[equal height group=lfef]{Learned features}
      \end{tdblockt}
    \end{column}
    \begin{column}{0.47\textwidth}
      \begin{tdblockt}[equal height group=lfef]{Engineered features}
      \end{tdblockt}
    \end{column}
  \end{columns}
  \slidesources{Huyen2022}
\end{frame}

% ----------------------------------------------------------------------
\subsection{Common Feature Engineering Operations}

% TEMPLATE: booktabs overview table, the opener for Deletion and Imputation.
% The book introduces missing values with a sample dataset whose holes are the
% whole point, so the opener is that table and the blank cells stay blank on
% purpose. Five columns for the book's own fields, header row and row labels
% filled from the book. The two children below carry the two remedies.
\begin{frame}{Handling Missing Values}
  \footnotesize
  \renewcommand{\arraystretch}{1.25}
  \begin{tabular}{@{}lllll@{}}
    \toprule
     & & & & \\
    \midrule
     & & & & \\
     & & & & \\
     & & & & \\
     & & & & \\
    \bottomrule
  \end{tabular}
  \slidesources{Huyen2022}
\end{frame}

% TEMPLATE: tddef. The heading is a named technique, so it takes the titled
% definition box and nothing else competes with it on the frame. The box takes
% the book's own wording for what gets dropped and when dropping is safe; the
% two kinds of deletion the book distinguishes go in as the box's body.
\begin{frame}{Deletion}
  \begin{tddef}{Deletion}
  \end{tddef}
  \slidesources{Huyen2022}
\end{frame}

% TEMPLATE: tdcodet. Imputation is an operation the book shows as a fill
% applied to a column, so the dark mono panel carries it as the snippet the
% author would type in class rather than as prose. Caption is the heading. The
% panel takes the fill values the book names and the column it fills.
\begin{frame}{Imputation}
  \begin{tdcodet}{Imputation}
  \end{tdcodet}
  \slidesources{Huyen2022}
\end{frame}

% TEMPLATE: withmargin. Breaks the code run on purpose: the scaling point is
% one claim with one caveat, not a snippet. Main column takes the book's
% transformations and the range each maps into; margin takes the warning that
% the statistics come from the training split only, which the Data Leakage
% subchapter later turns into its own frame.
\begin{frame}{Scaling}
  \begin{withmargin}
    \begin{tdmain}
    \end{tdmain}
    \begin{tdmargin}
    \end{tdmargin}
  \end{withmargin}
  \slidesources{Huyen2022}
\end{frame}

% TEMPLATE: tdcodet. Discretization is bucketing, and the book makes it
% concrete with the bucket boundaries, which read as a listing and not as a
% sentence. Caption is the heading. The panel takes the book's own brackets;
% the arbitrariness of the boundaries belongs in the line under the panel.
\begin{frame}{Discretization}
  \begin{tdcodet}{Discretization}
  \end{tdcodet}
  \slidesources{Huyen2022}
\end{frame}

% TEMPLATE: withmargin, alternating with the code panels around it so the
% operations subchapter reads as claim, snippet, claim rather than six
% identical frames. Main column takes the book's problem, that the set of
% categories is not static in production; margin takes the hashing trick and
% the collision it trades for.
\begin{frame}{Encoding Categorical Features}
  \begin{withmargin}
    \begin{tdmain}
    \end{tdmain}
    \begin{tdmargin}
    \end{tdmargin}
  \end{withmargin}
  \slidesources{Huyen2022}
\end{frame}

% TEMPLATE: tdcodet. Crossing is a constructed column, so the panel shows the
% two source features and the combined one the book builds from them. Caption
% is the heading. The combinatorial blowup the book warns about is a number and
% goes in the line under the panel, not into the panel.
\begin{frame}{Feature Crossing}
  \begin{tdcodet}{Feature crossing}
  \end{tdcodet}
  \slidesources{Huyen2022}
\end{frame}

% TEMPLATE: two block contrast. The heading names both sides, so each takes one
% box levelled by the equal height group rather than a single slide of prose.
% Left box takes the lookup table the book describes, right box the function
% the book substitutes for it. The line under them is for what the two share.
\begin{frame}{Discrete and Continuous Positional Embeddings}
  \begin{columns}[T,onlytextwidth]
    \begin{column}{0.47\textwidth}
      \begin{tdblockt}[equal height group=pos]{Discrete}
      \end{tdblockt}
    \end{column}
    \begin{column}{0.47\textwidth}
      \begin{tdblockt}[equal height group=pos]{Continuous}
      \end{tdblockt}
    \end{column}
  \end{columns}
  \slidesources{Huyen2022}
\end{frame}

% ----------------------------------------------------------------------
\subsection{Data Leakage}

% TEMPLATE: three plus three block grid, the opener for the six causes below.
% The book names exactly six peers, none subordinate to another, so one box
% each sized to the count; the headers are this chapter's own six subheadings,
% shortened to fit the grid, and each box takes one line naming the mistake.
% Every box then gets its own frame, in the same order, underneath.
\begin{frame}{Common Causes for Data Leakage}
  \footnotesize
  \begin{columns}[T,onlytextwidth]
    \begin{column}{0.31\textwidth}
      \begin{tdblockt}[equal height group=cause1]{Time-correlated splits}
      \end{tdblockt}
    \end{column}
    \begin{column}{0.31\textwidth}
      \begin{tdblockt}[equal height group=cause1]{Scaling before splitting}
      \end{tdblockt}
    \end{column}
    \begin{column}{0.31\textwidth}
      \begin{tdblockt}[equal height group=cause1]{Test split statistics}
      \end{tdblockt}
    \end{column}
  \end{columns}
  \vskip0.5em
  \begin{columns}[T,onlytextwidth]
    \begin{column}{0.31\textwidth}
      \begin{tdblockt}[equal height group=cause2]{Data duplication}
      \end{tdblockt}
    \end{column}
    \begin{column}{0.31\textwidth}
      \begin{tdblockt}[equal height group=cause2]{Group leakage}
      \end{tdblockt}
    \end{column}
    \begin{column}{0.31\textwidth}
      \begin{tdblockt}[equal height group=cause2]{Data generation process}
      \end{tdblockt}
    \end{column}
  \end{columns}
  \slidesources{Huyen2022}
\end{frame}

% TEMPLATE: phimgside hero. The book argues this cause with a figure, the same
% data cut at random against cut by time, and the two cuts only make their
% point side by side, so the strip holds that figure and the caption names it.
% First of this chapter's two image frames. The body takes the book's claim and
% its stock price example, the reason tomorrow must not train today's model.
\begin{frame}{Splitting time-correlated data randomly instead of by time}
  \phimgside[0.33]{splitting_time_correlated_data_randomly_instead_of_by_time.png}{}
  \begin{sidebody}
  \end{sidebody}
  \slidesources{Huyen2022}
\end{frame}

% TEMPLATE: tdcodet. This cause is an ordering mistake in code, and the book
% shows it as one: the panel takes the wrong order and the right order as two
% short blocks so the difference is a line of code and not an argument.
% Caption is the heading. The global statistics that leak go in the panel.
\begin{frame}{Scaling before splitting}
  \begin{tdcodet}{Scaling before splitting}
  \end{tdcodet}
  \slidesources{Huyen2022}
\end{frame}

% TEMPLATE: withmargin. One mistake with one consequence, which is the layout
% this chapter uses when a claim needs a single aside. Main column takes the
% book's statistics that must never be computed across the split; margin takes
% the rule the book states to avoid it, split first and then impute.
\begin{frame}{Filling in missing data with statistics from the test split}
  \begin{withmargin}
    \begin{tdmain}
    \end{tdmain}
    \begin{tdmargin}
    \end{tdmargin}
  \end{withmargin}
  \slidesources{Huyen2022}
\end{frame}

% TEMPLATE: tddef. The heading names duplication as the thing to be understood,
% so the titled definition box carries it and the frame stays on one term. The
% box takes the book's sources of duplicates, the collection and processing
% step and the oversampling step, and what a duplicate does across the split.
\begin{frame}{Poor handling of data duplication before splitting}
  \begin{tddef}{Data duplication}
  \end{tddef}
  \slidesources{Huyen2022}
\end{frame}

% TEMPLATE: withmargin. The book makes this one land with a single example, so
% the main column takes the claim about examples that are strongly correlated
% and split apart anyway, and the margin takes that example, the repeated scans
% of the same patient, which is the proof and not a decoration.
\begin{frame}{Group leakage}
  \begin{withmargin}
    \begin{tdmain}
    \end{tdmain}
    \begin{tdmargin}
    \end{tdmargin}
  \end{withmargin}
  \slidesources{Huyen2022}
\end{frame}

% TEMPLATE: phimgside hero, the second and last image frame of the chapter.
% The heading is about a process, how the data came to exist, so the strip
% takes the provenance diagram from collection through the instrument to the
% label. The body takes the book's example of a label that rides along on the
% machine that produced the sample, and the advice to track data lineage.
\begin{frame}{Leakage from data generation process}
  \phimgside[0.33]{leakage_from_data_generation_process.png}{}
  \begin{sidebody}
  \end{sidebody}
  \slidesources{Huyen2022}
\end{frame}

% TEMPLATE: withmargin. The book answers this heading with measurements to run
% rather than with a list of peers, so the main column takes those measurements
% in the book's order and the margin takes the discipline around them, when in
% the lifecycle to run them and who on the team owns the test split.
\begin{frame}{Detecting Data Leakage}
  \begin{withmargin}
    \begin{tdmain}
    \end{tdmain}
    \begin{tdmargin}
    \end{tdmargin}
  \end{withmargin}
  \slidesources{Huyen2022}
\end{frame}

% ----------------------------------------------------------------------
\subsection{Engineering Good Features}

% TEMPLATE: booktabs table. The heading names both columns, feature and
% importance, and the book's point here is an ordering of numbers, which only
% reads as a ranking in a table. Rows take the book's own measured features,
% most important first. No vertical rules, booktabs only. The methods that
% produce the numbers belong in the line under the table.
\begin{frame}{Feature Importance}
  \footnotesize
  \renewcommand{\arraystretch}{1.25}
  \begin{tabular}{@{}ll@{}}
    \toprule
    \textbf{Feature} & \textbf{Importance}\\
    \midrule
     & \\
     & \\
     & \\
     & \\
     & \\
    \bottomrule
  \end{tabular}
  \slidesources{Huyen2022}
\end{frame}

% TEMPLATE: withmargin. The book measures generalization along two axes, and
% neither is a peer box because the second qualifies the first, so the main
% column runs both in the book's order. Margin takes the worked counterexample,
% the feature that scores well and still does not transfer to unseen data.
\begin{frame}{Feature Generalization}
  \begin{withmargin}
    \begin{tdmain}
    \end{tdmain}
    \begin{tdmargin}
    \end{tdmargin}
  \end{withmargin}
  \slidesources{Huyen2022}
\end{frame}

% ----------------------------------------------------------------------
\subsection{Summary}

% TEMPLATE: bare takeaway. The chapter's closing heading gets one centred line
% and nothing around it, so the sentence the lecture should leave behind is the
% only thing on the frame. Keep it to the book's own summary claim about
% features and model performance; no bullets, no recap of the subchapters.
\begin{frame}{Summary}
  \begin{takeaway}
  \end{takeaway}
  \slidesources{Huyen2022}
\end{frame}

% ----------------------------------------------------------------------
% This chapter's own references. The refsection opened above scopes the
% citation numbering to this chapter, so chapter_pdfs/<this>.pdf resolves
% its own [n] markers instead of pointing at a list it does not contain.
\begin{frame}{References}
  \printbibliography[heading=none]
\end{frame}
\end{refsection}

%%% ch06_model_development.tex -- Intelligence Engineering, chapter 06 "Model
%% Development and Offline Evaluation".
%% Spine transferred from Huyen, Designing Machine Learning Systems, chapter
%% 6. The subchapters and the frame titles under them are the book's own
%% section and subsection headings, in the book's order.
%%
%% STATE: titles and TEMPLATE layout only. No body text. Every frame carries
%% the layout it is meant to be filled into, and a % TEMPLATE comment saying
%% why that layout and what belongs in each slot. Filling them is the next pass.
%%
%% Fragment of intelligence_engineering.tex. One chapter is one lecture and
%% gets exactly ONE \lecturepage, at the top. Inside it every \subsection
%% opens on the subchapter divider, which the master emits automatically via
%% \AtBeginSubsection; do not write those divider frames by hand.

\begin{refsection}

\lecturepage{VI}{Model Development and Offline Evaluation}{}
\section[Model Development]{Model Development and Offline Evaluation}

% ----------------------------------------------------------------------
\subsection{Model Development and Training}

% TEMPLATE: withmargin, the chapter's opening claim plus one aside. This frame
% is the opener for the six tips below, so the main column takes the book's own
% framing of how a model is chosen for a problem, and the margin takes the
% caveat that qualifies it. Keep the six tips off this slide; each one gets its
% own box on the next frame.
\begin{frame}{Evaluating ML Models}
  \begin{withmargin}
    \begin{tdmain}
    \end{tdmain}
    \begin{tdmargin}
    \end{tdmargin}
  \end{withmargin}
  \slidesources{Huyen2022}
\end{frame}

% TEMPLATE: three by two block grid. The heading names a countable set of six
% peers, none subordinate to another, so one box each rather than a six item
% bullet list, levelled row by row with the equal height groups. Boxes are
% untitled on purpose: the heading counts the tips but does not name them, so
% each box takes one tip's name as its header and the book's reasoning as its
% body when the slide is filled.
\begin{frame}{Six tips for model selection}
  \begin{columns}[T,onlytextwidth]
    \begin{column}{0.31\textwidth}
      \begin{tdblock}[equal height group=tips1]
      \end{tdblock}
    \end{column}
    \begin{column}{0.31\textwidth}
      \begin{tdblock}[equal height group=tips1]
      \end{tdblock}
    \end{column}
    \begin{column}{0.31\textwidth}
      \begin{tdblock}[equal height group=tips1]
      \end{tdblock}
    \end{column}
  \end{columns}
  \vskip0.5em
  \begin{columns}[T,onlytextwidth]
    \begin{column}{0.31\textwidth}
      \begin{tdblock}[equal height group=tips2]
      \end{tdblock}
    \end{column}
    \begin{column}{0.31\textwidth}
      \begin{tdblock}[equal height group=tips2]
      \end{tdblock}
    \end{column}
    \begin{column}{0.31\textwidth}
      \begin{tdblock}[equal height group=tips2]
      \end{tdblock}
    \end{column}
  \end{columns}
  \slidesources{Huyen2022}
\end{frame}

% TEMPLATE: three block row, the overview layout for an opener. The heading
% introduces exactly three peers, its own subheadings, so the three take one
% levelled box each and the slide works as the map of the three frames that
% follow. Each box takes that method's one line definition in the book's words;
% the detail stays on the method's own frame.
\begin{frame}{Ensembles}
  \begin{columns}[T,onlytextwidth]
    \begin{column}{0.31\textwidth}
      \begin{tdblockt}[equal height group=ens]{Bagging}
      \end{tdblockt}
    \end{column}
    \begin{column}{0.31\textwidth}
      \begin{tdblockt}[equal height group=ens]{Boosting}
      \end{tdblockt}
    \end{column}
    \begin{column}{0.31\textwidth}
      \begin{tdblockt}[equal height group=ens]{Stacking}
      \end{tdblockt}
    \end{column}
  \end{columns}
  \slidesources{Huyen2022}
\end{frame}

% TEMPLATE: tddef. The heading is a term the book defines outright, bootstrap
% aggregating, so it takes the titled definition box rather than a claim with an
% aside. The box holds the sampling with replacement procedure and the vote or
% average that closes it; the variance argument follows inside the same box.
\begin{frame}{Bagging}
  \begin{tddef}{Bagging}
  \end{tddef}
  \slidesources{Huyen2022}
\end{frame}

% TEMPLATE: side image hero, the first of this chapter's two image slots. The
% book carries this one as a diagram, a sequence of weak learners each reweighting
% what the previous one got wrong, and a sequence is a picture before it is a
% sentence. The strip takes that diagram, the caption names it, and the sidebody
% takes the iteration in the book's steps.
\begin{frame}{Boosting}
  \phimgside[0.33]{boosting.png}{}
  \begin{sidebody}
  \end{sidebody}
  \slidesources{Huyen2022}
\end{frame}

% TEMPLATE: withmargin. Stacking has one mechanism, base learners feeding a
% meta learner, so the main column runs it in the book's order rather than
% splitting into boxes whose headers the heading does not license. Margin takes
% the aside on how the combining is done, majority vote or the learned
% combination.
\begin{frame}{Stacking}
  \begin{withmargin}
    \begin{tdmain}
    \end{tdmain}
    \begin{tdmargin}
    \end{tdmargin}
  \end{withmargin}
  \slidesources{Huyen2022}
\end{frame}

% TEMPLATE: two block contrast, levelled by the equal height group. The heading
% names two artefacts, and the book treats them as two jobs rather than one, so
% each takes a box headed with its own name from the heading. Each box takes what
% that job records and why it is kept; the two frames after this one open them
% up one at a time.
\begin{frame}{Experiment Tracking and Versioning}
  \begin{columns}[T,onlytextwidth]
    \begin{column}{0.47\textwidth}
      \begin{tdblockt}[equal height group=etv]{Experiment tracking}
      \end{tdblockt}
    \end{column}
    \begin{column}{0.47\textwidth}
      \begin{tdblockt}[equal height group=etv]{Versioning}
      \end{tdblockt}
    \end{column}
  \end{columns}
  \slidesources{Huyen2022}
\end{frame}

% TEMPLATE: withmargin. The book answers this heading with a list of what to
% watch during training, loss curves, performance metrics, speed and system
% metrics, so the main column takes that list in its order. Margin takes the
% warning the book attaches, that tracking everything quickly becomes
% overwhelming.
\begin{frame}{Experiment tracking}
  \begin{withmargin}
    \begin{tdmain}
    \end{tdmain}
    \begin{tdmargin}
    \end{tdmargin}
  \end{withmargin}
  \slidesources{Huyen2022}
\end{frame}

% TEMPLATE: code panel, caption taken from the heading. Versioning is the one
% heading here that is about commands and artefacts on disk, so the dark mono
% panel is the honest layout: it takes the book's code and data versioning
% snippet, the checksum or the tool invocation. The argument about why data
% versioning is the hard half goes under the panel when filled.
\begin{frame}{Versioning}
  \begin{tdcodet}{Versioning}
  \end{tdcodet}
  \slidesources{Huyen2022}
\end{frame}

% TEMPLATE: withmargin, the opener for the two parallelism frames. The book sets
% this one up as a condition rather than a comparison, the data or the model no
% longer fits on one machine, so the main column takes that condition and what
% it forces. Margin takes the practical aside, batch size and checkpointing.
\begin{frame}{Distributed Training}
  \begin{withmargin}
    \begin{tdmain}
    \end{tdmain}
    \begin{tdmargin}
    \end{tdmargin}
  \end{withmargin}
  \slidesources{Huyen2022}
\end{frame}

% TEMPLATE: tddef. The heading is a named technique with a one sentence
% definition in the book, so it takes the definition box. The box holds what is
% split across machines and how the gradients come back together; the
% synchronous and asynchronous variants follow inside the same box, since the
% heading does not license two boxes of their own.
\begin{frame}{Data parallelism}
  \begin{tddef}{Data parallelism}
  \end{tddef}
  \slidesources{Huyen2022}
\end{frame}

% TEMPLATE: side image hero, the second and last image slot of the chapter. The
% book draws this one, a model cut into components that sit on different machines
% and the pipeline that keeps them busy, and the picture carries the point that
% the parts are not independent. Strip takes the diagram, caption names it,
% sidebody takes the book's reading of it.
\begin{frame}{Model parallelism}
  \phimgside[0.33]{model_parallelism.png}{}
  \begin{sidebody}
  \end{sidebody}
  \slidesources{Huyen2022}
\end{frame}

% TEMPLATE: comparison table, the chapter's one table, and a deliberate push
% away from a third withmargin in this subchapter. The heading's two subheadings
% are soft and hard AutoML, so they take the two columns; the empty stub column
% takes one comparison axis per row, what is searched, what it costs and how far
% the practice has got. Four rows are provided, drop the ones that stay empty.
\begin{frame}{AutoML}
  \footnotesize
  \renewcommand{\arraystretch}{1.25}
  \begin{tabular}{@{}%
    >{\raggedright\arraybackslash}p{0.22\textwidth}%
    >{\raggedright\arraybackslash}p{0.35\textwidth}%
    >{\raggedright\arraybackslash}p{0.35\textwidth}@{}}
    \toprule
     & \textbf{Soft AutoML} & \textbf{Hard AutoML}\\
    \midrule
     & & \\
     & & \\
     & & \\
     & & \\
    \bottomrule
  \end{tabular}
  \slidesources{Huyen2022}
\end{frame}

% TEMPLATE: code panel, caption taken from the heading. Hyperparameter tuning is
% configuration before it is prose: the panel takes the book's search space and
% the search it is handed to, so the difference between the searches reads off
% the same snippet. The book's warning about tuning on the test split belongs
% under the panel, not inside it.
\begin{frame}{Soft AutoML: Hyperparameter tuning}
  \begin{tdcodet}{Hyperparameter tuning}
  \end{tdcodet}
  \slidesources{Huyen2022}
\end{frame}

% TEMPLATE: two block contrast, levelled by the equal height group. The heading
% itself names two things, so each takes a box headed with its own name and the
% two read off against one another rather than queueing in one column. Each box
% takes what that line of work hands over to the search and what it costs.
\begin{frame}{Hard AutoML: Architecture search and learned optimizer}
  \begin{columns}[T,onlytextwidth]
    \begin{column}{0.47\textwidth}
      \begin{tdblockt}[equal height group=hard]{Architecture search}
      \end{tdblockt}
    \end{column}
    \begin{column}{0.47\textwidth}
      \begin{tdblockt}[equal height group=hard]{Learned optimizer}
      \end{tdblockt}
    \end{column}
  \end{columns}
  \slidesources{Huyen2022}
\end{frame}

% ----------------------------------------------------------------------
\subsection{Model Offline Evaluation}

% TEMPLATE: withmargin plus a takeaway. The book answers this heading with a
% short list of baselines to compare against, so the main column takes them in
% the book's order and the margin takes the worked example that shows a number
% is meaningless without one. The takeaway keeps the book's own verdict on
% comparing against the right baseline in front of the list.
\begin{frame}{Baselines}
  \begin{withmargin}
    \begin{tdmain}
    \end{tdmain}
    \begin{tdmargin}
    \end{tdmargin}
  \end{withmargin}
  \begin{takeaway}
  \end{takeaway}
  \slidesources{Huyen2022}
\end{frame}

% TEMPLATE: three by two block grid, the overview layout for an opener with six
% children. The heading introduces exactly six methods, its own subheadings, so
% the grid is the map of the six frames that follow and each box is headed with
% that method's name. One line of what the method does per box, nothing more;
% the detail belongs on the method's own frame.
\begin{frame}{Evaluation Methods}
  \begin{columns}[T,onlytextwidth]
    \begin{column}{0.31\textwidth}
      \begin{tdblockt}[equal height group=em1]{Perturbation tests}
      \end{tdblockt}
    \end{column}
    \begin{column}{0.31\textwidth}
      \begin{tdblockt}[equal height group=em1]{Invariance tests}
      \end{tdblockt}
    \end{column}
    \begin{column}{0.31\textwidth}
      \begin{tdblockt}[equal height group=em1]{Directional expectation tests}
      \end{tdblockt}
    \end{column}
  \end{columns}
  \vskip0.5em
  \begin{columns}[T,onlytextwidth]
    \begin{column}{0.31\textwidth}
      \begin{tdblockt}[equal height group=em2]{Model calibration}
      \end{tdblockt}
    \end{column}
    \begin{column}{0.31\textwidth}
      \begin{tdblockt}[equal height group=em2]{Confidence measurement}
      \end{tdblockt}
    \end{column}
    \begin{column}{0.31\textwidth}
      \begin{tdblockt}[equal height group=em2]{Slice-based evaluation}
      \end{tdblockt}
    \end{column}
  \end{columns}
  \slidesources{Huyen2022}
\end{frame}

% TEMPLATE: withmargin. One test, one claim, so the main column takes what is
% perturbed and what the model's response to it is read as. Margin takes the
% book's worked example, the input that a user would never hand over as cleanly
% as the training set did.
\begin{frame}{Perturbation tests}
  \begin{withmargin}
    \begin{tdmain}
    \end{tdmain}
    \begin{tdmargin}
    \end{tdmargin}
  \end{withmargin}
  \slidesources{Huyen2022}
\end{frame}

% TEMPLATE: tddef. The heading is a test with a sharp definition, the input
% change that must leave the prediction alone, so the titled box carries it and
% the fairness consequence lands inside the same box. A definition box also
% breaks the withmargin run either side of it.
\begin{frame}{Invariance tests}
  \begin{tddef}{Invariance tests}
  \end{tddef}
  \slidesources{Huyen2022}
\end{frame}

% TEMPLATE: withmargin. The test is one rule about the sign of a change, so the
% main column states the rule and what a violation means. Margin takes the
% book's worked example, the feature whose increase may not move the prediction
% the wrong way.
\begin{frame}{Directional expectation tests}
  \begin{withmargin}
    \begin{tdmain}
    \end{tdmain}
    \begin{tdmargin}
    \end{tdmargin}
  \end{withmargin}
  \slidesources{Huyen2022}
\end{frame}

% TEMPLATE: code panel, caption taken from the heading. Calibration is checked
% rather than argued, and the book checks it with a curve and a library call, so
% the panel takes that call and the frequencies it is fed. The property being
% tested, predicted scores matching observed frequencies, goes under the panel
% when filled.
\begin{frame}{Model calibration}
  \begin{tdcodet}{Model calibration}
  \end{tdcodet}
  \slidesources{Huyen2022}
\end{frame}

% TEMPLATE: withmargin. The heading asks one question, whether a single
% prediction is worth showing, so the main column takes the threshold and what
% happens to the predictions below it. Margin takes the cost the book names for
% each choice, the useless prediction against the one that is never shown.
\begin{frame}{Confidence measurement}
  \begin{withmargin}
    \begin{tdmain}
    \end{tdmain}
    \begin{tdmargin}
    \end{tdmargin}
  \end{withmargin}
  \slidesources{Huyen2022}
\end{frame}

% TEMPLATE: withmargin plus a takeaway, closing the subchapter the way it
% opened. The book argues this one with numbers on slices of one dataset, so the
% main column takes the slice comparison and the margin takes why the slices
% exist at all, critical subgroups and the paradox behind the aggregate. The
% takeaway keeps the book's conclusion about the aggregate number.
\begin{frame}{Slice-based evaluation}
  \begin{withmargin}
    \begin{tdmain}
    \end{tdmain}
    \begin{tdmargin}
    \end{tdmargin}
  \end{withmargin}
  \begin{takeaway}
  \end{takeaway}
  \slidesources{Huyen2022}
\end{frame}

% ----------------------------------------------------------------------
\subsection{Summary}

% TEMPLATE: bare takeaway. The heading is the chapter's summary, so it gets one
% line and no scaffolding: the book's own closing claim about model development,
% big and centred, nothing competing with it on the slide.
\begin{frame}{Summary}
  \begin{takeaway}
  \end{takeaway}
  \slidesources{Huyen2022}
\end{frame}

% ----------------------------------------------------------------------
% This chapter's own references. The refsection opened above scopes the
% citation numbering to this chapter, so chapter_pdfs/<this>.pdf resolves
% its own [n] markers instead of pointing at a list it does not contain.
\begin{frame}{References}
  \printbibliography[heading=none]
\end{frame}
\end{refsection}

%%% ch07_deployment.tex -- Intelligence Engineering, chapter 07 "Model
%% Deployment and Prediction Service".
%% Spine transferred from Huyen, Designing Machine Learning Systems, chapter
%% 7. The subchapters and the frame titles under them are the book's own
%% section and subsection headings, in the book's order.
%%
%% STATE: titles and TEMPLATE layout only. No body text. Every frame carries
%% the layout it is meant to be filled into, and a % TEMPLATE comment saying
%% why that layout and what belongs in each slot. Filling them is the next pass.
%%
%% Fragment of intelligence_engineering.tex. One chapter is one lecture and
%% gets exactly ONE \lecturepage, at the top. Inside it every \subsection
%% opens on the subchapter divider, which the master emits automatically via
%% \AtBeginSubsection; do not write those divider frames by hand.

\begin{refsection}

\lecturepage{VII}{Model Deployment and Prediction Service}{}
\section[Deployment]{Model Deployment and Prediction Service}

% ----------------------------------------------------------------------
\subsection{Machine Learning Deployment Myths}

% TEMPLATE: withmargin, the chapter's workhorse and the right opener for a
% myth: one claim the book contradicts plus one aside that does the
% contradicting. Main column states the myth and then the correction in the
% book's words. Margin takes the counts, the models behind a single ride
% sharing app and the company totals the book quotes.
\begin{frame}{Myth 1: You Only Deploy One or Two ML Models at a Time}
  \begin{withmargin}
    \begin{tdmain}
    \end{tdmain}
    \begin{tdmargin}
    \end{tdmargin}
  \end{withmargin}
  \slidesources{Huyen2022}
\end{frame}

% TEMPLATE: one plain box plus a takeaway, and deliberately not a second
% withmargin in a row. The heading already is the myth as a full sentence, so
% the box needs no header over it: it holds the myth and the reasons a
% deployed model decays. The takeaway carries the book's verdict on software
% rot, which only lands as a line of its own.
\begin{frame}{Myth 2: If We Don't Do Anything, Model Performance Remains the
    Same}
  \begin{tdblock}
  \end{tdblock}
  \begin{takeaway}
  \end{takeaway}
  \slidesources{Huyen2022}
\end{frame}

% TEMPLATE: withmargin. The heading is a claim about frequency, so the main
% column takes what forces an update and how often the book says it happens.
% Margin holds the named deployment cadences the book cites company by
% company, which are the evidence and read better stacked than in prose.
\begin{frame}{Myth 3: You Won't Need to Update Your Models as Much}
  \begin{withmargin}
    \begin{tdmain}
    \end{tdmain}
    \begin{tdmargin}
    \end{tdmargin}
  \end{withmargin}
  \slidesources{Huyen2022}
\end{frame}

% TEMPLATE: tddef, kept bare. This myth turns on one word the heading names,
% scale, and the book settles it by saying what scale means before arguing
% who has to worry about it. The term box is the slide's only slot: it takes
% the book's own measure of scale and the share of engineers it covers.
\begin{frame}{Myth 4: Most ML Engineers Don't Need to Worry About Scale}
  \begin{tddef}{Scale}
  \end{tddef}
  \slidesources{Huyen2022}
\end{frame}

% ----------------------------------------------------------------------
\subsection{Batch Prediction Versus Online Prediction}

% TEMPLATE: booktabs comparison table. The subchapter heading is the versus
% and the book already carries this contrast as a table, so it stays a table
% here rather than two boxes. The two column headers are the two regimes the
% heading names; the row labels are the book's own comparison dimensions and
% get filled together with the cells in the next pass.
\begin{frame}{From Batch Prediction to Online Prediction}
  \footnotesize
  \renewcommand{\arraystretch}{1.35}
  \begin{tabular}{@{}%
    >{\raggedright\arraybackslash}p{0.22\textwidth}%
    >{\raggedright\arraybackslash}p{0.35\textwidth}%
    >{\raggedright\arraybackslash}p{0.35\textwidth}@{}}
    \toprule
     & \textbf{Batch prediction} & \textbf{Online prediction}\\
    \midrule
     & & \\
     & & \\
     & & \\
     & & \\
    \bottomrule
  \end{tabular}
  \slidesources{Huyen2022}
\end{frame}

% TEMPLATE: side image hero, the first of this chapter's two. The heading asks
% an architecture question and the book answers it with a drawing of the two
% pipelines, so the strip takes that figure and the sidebody takes the steps
% in the book's order. The caption names the figure; the gray placeholder
% stays until the asset is drawn.
\begin{frame}{Unifying Batch Pipeline and Streaming Pipeline}
  \phimgside[0.33]{unifying_batch_pipeline_and_streaming_pipeline.png}{Batch pipeline and streaming pipeline}
  \begin{sidebody}
  \end{sidebody}
  \slidesources{Huyen2022}
\end{frame}

% ----------------------------------------------------------------------
\subsection{Model Compression}

% TEMPLATE: withmargin. The heading names one technique, so the main column
% takes what it replaces and the parameter arithmetic that makes it a
% compression method at all. Margin holds the architectures the book names as
% the worked cases, with their counts.
\begin{frame}{Low-Rank Factorization}
  \begin{withmargin}
    \begin{tdmain}
    \end{tdmain}
    \begin{tdmargin}
    \end{tdmargin}
  \end{withmargin}
  \slidesources{Huyen2022}
\end{frame}

% TEMPLATE: tddef. The heading is a named method the book introduces by
% definition, so the term box is the whole slide rather than a claim with an
% aside. The box takes the book's definition and the tradeoff it states
% immediately after it.
\begin{frame}{Knowledge Distillation}
  \begin{tddef}{Knowledge distillation}
  \end{tddef}
  \slidesources{Huyen2022}
\end{frame}

% TEMPLATE: two untitled boxes, levelled by the equal height group. The book
% gives pruning two distinct meanings, so each takes a box of its own instead
% of two bullets in one. No box headers here on purpose: the heading names
% only pruning, and the two senses arrive with the text in the next pass.
\begin{frame}{Pruning}
  \begin{columns}[T,onlytextwidth]
    \begin{column}{0.47\textwidth}
      \begin{tdblock}[equal height group=prune]
      \end{tdblock}
    \end{column}
    \begin{column}{0.47\textwidth}
      \begin{tdblock}[equal height group=prune]
      \end{tdblock}
    \end{column}
  \end{columns}
  \slidesources{Huyen2022}
\end{frame}

% TEMPLATE: code panel plus a takeaway. Quantization is a statement about how
% a number is represented, so the panel is the honest slot: it takes the bit
% widths and the ranges the book lists, which only make their point written
% out. The takeaway keeps the book's caveat about rounding and overflow.
\begin{frame}{Quantization}
  \begin{tdcode}
  \end{tdcode}
  \begin{takeaway}
  \end{takeaway}
  \slidesources{Huyen2022}
\end{frame}

% ----------------------------------------------------------------------
\subsection{ML on the Cloud and on the Edge}

% TEMPLATE: two block contrast, levelled by the equal height group. This frame
% opens the two deep headings under it, so it stays an overview and not a
% detail layout. The heading names two operations and each takes one box,
% compiling a model for a target and optimizing it for that target; the boxes
% carry the book's account of each step and the children carry the detail.
\begin{frame}{Compiling and Optimizing Models for Edge Devices}
  \begin{columns}[T,onlytextwidth]
    \begin{column}{0.47\textwidth}
      \begin{tdblockt}[equal height group=edge]{Compiling}
      \end{tdblockt}
    \end{column}
    \begin{column}{0.47\textwidth}
      \begin{tdblockt}[equal height group=edge]{Optimizing}
      \end{tdblockt}
    \end{column}
  \end{columns}
  \slidesources{Huyen2022}
\end{frame}

% TEMPLATE: side image hero, the second and last of this chapter's two. The
% book makes this heading concrete with a drawing of what an optimization
% does to a computation graph, so the strip takes that figure and the sidebody
% takes the named techniques under it. Caption repeats the heading as the
% figure's name; placeholder stays until the asset exists.
\begin{frame}{Model optimization}
  \phimgside[0.33]{model_optimization.png}{Model optimization}
  \begin{sidebody}
  \end{sidebody}
  \slidesources{Huyen2022}
\end{frame}

% TEMPLATE: withmargin. The heading is a turn of the argument rather than a
% set of peers, so the main column runs it as one claim: the search space, and
% what a learned model does inside it. Margin takes the named system the book
% uses as its case and the cost the book attaches to the search.
\begin{frame}{Using ML to optimize ML models}
  \begin{withmargin}
    \begin{tdmain}
    \end{tdmain}
    \begin{tdmargin}
    \end{tdmargin}
  \end{withmargin}
  \slidesources{Huyen2022}
\end{frame}

% TEMPLATE: withmargin with a captioned code panel in the main column. The
% book's browser material is named runtimes and libraries, so the panel is
% what the slide should show rather than a sentence about them. Margin takes
% the promise the book attaches to running in a browser and the limit it
% names in the same breath.
\begin{frame}{ML in Browsers}
  \begin{withmargin}
    \begin{tdmain}
      \begin{tdcodet}{Browser}
      \end{tdcodet}
    \end{tdmain}
    \begin{tdmargin}
    \end{tdmargin}
  \end{withmargin}
  \slidesources{Huyen2022}
\end{frame}

% ----------------------------------------------------------------------
\subsection{Summary}

% TEMPLATE: statement, kept bare. A summary is a closing line, not a slide of
% bullets, so it takes the loud centred form and nothing else. The main line
% is the book's own summary of what deployment costs; the substatement is the
% consequence it carries into the next chapter.
\begin{frame}{Summary}
  \begin{statement}
    \substatement{}
  \end{statement}
  \slidesources{Huyen2022}
\end{frame}

% ----------------------------------------------------------------------
% This chapter's own references. The refsection opened above scopes the
% citation numbering to this chapter, so chapter_pdfs/<this>.pdf resolves
% its own [n] markers instead of pointing at a list it does not contain.
\begin{frame}{References}
  \printbibliography[heading=none]
\end{frame}
\end{refsection}

%%% ch08_distribution_shifts.tex -- Intelligence Engineering, chapter 08 "Data
%% Distribution Shifts and Monitoring".
%% Spine transferred from Huyen, Designing Machine Learning Systems, chapter
%% 8. The subchapters and the frame titles under them are the book's own
%% section and subsection headings, in the book's order.
%%
%% STATE: titles and TEMPLATE layout only. No body text. Every frame carries
%% the layout it is meant to be filled into, and a % TEMPLATE comment saying
%% why that layout and what belongs in each slot. Filling them is the next pass.
%%
%% Fragment of intelligence_engineering.tex. One chapter is one lecture and
%% gets exactly ONE \lecturepage, at the top. Inside it every \subsection
%% opens on the subchapter divider, which the master emits automatically via
%% \AtBeginSubsection; do not write those divider frames by hand.

\begin{refsection}

\lecturepage{VIII}{Data Distribution Shifts and Monitoring}{}
\section[Distribution Shifts]{Data Distribution Shifts and Monitoring}

% ----------------------------------------------------------------------
\subsection{Causes of ML System Failures}

% TEMPLATE: two by two block grid. The heading is plural and the book answers
% it with four peer causes of the same kind (dependency failure, deployment
% failure, hardware failure, downtime), so one box each rather than a bullet
% list, levelled by the equal height group. The box headers stay empty until
% the next pass names each box after the book's own cause.
\begin{frame}{Software System Failures}
  \begin{columns}[T,onlytextwidth]
    \begin{column}{0.47\textwidth}
      \begin{tdblock}[equal height group=ssf]
      \end{tdblock}
      \vskip0.5em
      \begin{tdblock}[equal height group=ssf]
      \end{tdblock}
    \end{column}
    \begin{column}{0.47\textwidth}
      \begin{tdblock}[equal height group=ssf]
      \end{tdblock}
      \vskip0.5em
      \begin{tdblock}[equal height group=ssf]
      \end{tdblock}
    \end{column}
  \end{columns}
  \slidesources{Huyen2022}
\end{frame}

% TEMPLATE: withmargin, the opener for the three child frames below. An opener
% needs room for an overview rather than a detail layout, so the main column
% names the three ML-specific failures in the book's order and leaves the
% detail to their own frames. Margin takes the aside on why these failures are
% harder to detect than the software ones on the frame before.
\begin{frame}{ML-Specific Failures}
  \begin{withmargin}
    \begin{tdmain}
    \end{tdmain}
    \begin{tdmargin}
    \end{tdmargin}
  \end{withmargin}
  \slidesources{Huyen2022}
\end{frame}

% TEMPLATE: booktabs comparison table. The heading names both sides of a
% difference and the book separates them along named dimensions, so it reads
% as a table rather than as two boxes, and it keeps this subchapter off a
% second box grid in three frames. The two column headers are the sides the
% heading names; the row labels are the book's own dimensions and get filled
% together with the cells in the next pass.
\begin{frame}{Production data differing from training data}
  \footnotesize
  \renewcommand{\arraystretch}{1.35}
  \begin{tabular}{@{}>{\raggedright\arraybackslash}p{0.24\textwidth} >{\raggedright\arraybackslash}p{0.32\textwidth} >{\raggedright\arraybackslash}p{0.32\textwidth}@{}}
    \toprule
     & \textbf{Training data} & \textbf{Production data}\\
    \midrule
     & & \\
     & & \\
     & & \\
     & & \\
    \bottomrule
  \end{tabular}
  \slidesources{Huyen2022}
\end{frame}

% TEMPLATE: withmargin with a definition box. The heading is a term the book
% defines and then separates from a neighbouring term, so the definition takes
% the main column's one titled box and the margin holds the distinction the
% book draws against outliers.
\begin{frame}{Edge cases}
  \begin{withmargin}
    \begin{tdmain}
      \begin{tddef}{Edge cases}
      \end{tddef}
    \end{tdmain}
    \begin{tdmargin}
    \end{tdmargin}
  \end{withmargin}
  \slidesources{Huyen2022}
\end{frame}

% TEMPLATE: side image hero on the grey placeholder, the first of this
% chapter's two image frames. A degenerate feedback loop is a dataflow, the
% predictions feeding the labels that feed the next model, so the book's own
% content here is a loop and it wants the figure, not a box. Sidebody takes
% the loop in the book's order, the caption names the loop.
\begin{frame}{Degenerate feedback loops}
  \phimgside[0.33]{degenerate_feedback_loops.png}{}
  \begin{sidebody}
  \end{sidebody}
  \slidesources{Huyen2022}
\end{frame}

% ----------------------------------------------------------------------
\subsection{Data Distribution Shifts}

% TEMPLATE: three block row, the opener for the three child frames below. The
% heading promises a countable set and its children name it, so each type
% takes one box headed with its own heading, levelled by the equal height
% group. Each box takes the one line that says which term of the joint
% distribution moves; the decomposition itself goes above the row.
\begin{frame}{Types of Data Distribution Shifts}
  \begin{columns}[T,onlytextwidth]
    \begin{column}{0.31\textwidth}
      \begin{tdblockt}[equal height group=types]{Covariate shift}
      \end{tdblockt}
    \end{column}
    \begin{column}{0.31\textwidth}
      \begin{tdblockt}[equal height group=types]{Label shift}
      \end{tdblockt}
    \end{column}
    \begin{column}{0.31\textwidth}
      \begin{tdblockt}[equal height group=types]{Concept drift}
      \end{tdblockt}
    \end{column}
  \end{columns}
  \slidesources{Huyen2022}
\end{frame}

% TEMPLATE: definition box on its own, full width. The heading is a term with
% a precise probabilistic meaning, so it takes the titled definition box and
% nothing competes with it on the frame. The causes the book lists and its
% worked example go inside that box.
\begin{frame}{Covariate shift}
  \begin{tddef}{Covariate shift}
  \end{tddef}
  \slidesources{Huyen2022}
\end{frame}

% TEMPLATE: withmargin. Label shift is not defined on its own but against the
% covariate shift on the frame before, so the main column carries the
% definition and the overlap between the two, and the margin takes the book's
% example where one shift brings the other with it.
\begin{frame}{Label shift}
  \begin{withmargin}
    \begin{tdmain}
    \end{tdmain}
    \begin{tdmargin}
    \end{tdmargin}
  \end{withmargin}
  \slidesources{Huyen2022}
\end{frame}

% TEMPLATE: withmargin with a definition box. Third of the three types and the
% only one with a time axis, so the definition keeps the main column's titled
% box for consistency with covariate shift, while the margin holds the cyclic
% and seasonal variation the book raises only here.
\begin{frame}{Concept drift}
  \begin{withmargin}
    \begin{tdmain}
      \begin{tddef}{Concept drift}
      \end{tddef}
    \end{tdmain}
    \begin{tdmargin}
    \end{tdmargin}
  \end{withmargin}
  \slidesources{Huyen2022}
\end{frame}

% TEMPLATE: three plain boxes in a row. Pushed to a grid on purpose, since a
% fourth withmargin in this subchapter would stall: the heading collects the
% shifts the three named types do not cover, and the book lists three of them
% (feature change, label schema change, non-stationarity), so one box each.
% Headers stay empty until the next pass names them.
\begin{frame}{General Data Distribution Shifts}
  \begin{columns}[T,onlytextwidth]
    \begin{column}{0.31\textwidth}
      \begin{tdblock}[equal height group=gen]
      \end{tdblock}
    \end{column}
    \begin{column}{0.31\textwidth}
      \begin{tdblock}[equal height group=gen]
      \end{tdblock}
    \end{column}
    \begin{column}{0.31\textwidth}
      \begin{tdblock}[equal height group=gen]
      \end{tdblock}
    \end{column}
  \end{columns}
  \slidesources{Huyen2022}
\end{frame}

% TEMPLATE: two block contrast, the opener for the two child frames below. The
% heading splits into the methods and the windows those methods run over, so
% each half takes one levelled box as an overview rather than a detail layout.
% Each box takes one line on what that half of detection does.
\begin{frame}{Detecting Data Distribution Shifts}
  \begin{columns}[T,onlytextwidth]
    \begin{column}{0.47\textwidth}
      \begin{tdblockt}[equal height group=det]{Statistical methods}
      \end{tdblockt}
    \end{column}
    \begin{column}{0.47\textwidth}
      \begin{tdblockt}[equal height group=det]{Time scale windows}
      \end{tdblockt}
    \end{column}
  \end{columns}
  \slidesources{Huyen2022}
\end{frame}

% TEMPLATE: captioned code panel. The book answers this heading with named
% tests and a library call, which read as a listing and not as prose, so the
% mono panel carries them. The caption is the heading itself; the panel takes
% the two sample test and the detector call in the book's form.
\begin{frame}{Statistical methods}
  \begin{tdcodet}{Statistical methods}
  \end{tdcodet}
  \slidesources{Huyen2022}
\end{frame}

% TEMPLATE: side image hero on the grey placeholder, the second and last image
% frame of the chapter, so the budget is spent here and on the feedback loop.
% The book makes this point with a plot, the same metric read as a cumulative
% against a sliding statistic, so the strip holds the figure and the sidebody
% the reading of it. Caption names the window the plot shows.
\begin{frame}{Time scale windows for detecting shifts}
  \phimgside[0.33]{time_scale_windows_for_detecting_shifts.png}{}
  \begin{sidebody}
  \end{sidebody}
  \slidesources{Huyen2022}
\end{frame}

% TEMPLATE: withmargin plus a takeaway. The book answers this heading with one
% decision and not a set of peers, train on a wider distribution or adapt the
% model, so the main column runs the options in the book's order and the
% margin holds the stateless against stateful retraining aside. The takeaway
% keeps the book's own caution about what retraining does not fix.
\begin{frame}{Addressing Data Distribution Shifts}
  \begin{withmargin}
    \begin{tdmain}
    \end{tdmain}
    \begin{tdmargin}
    \end{tdmargin}
  \end{withmargin}
  \begin{takeaway}
  \end{takeaway}
  \slidesources{Huyen2022}
\end{frame}

% ----------------------------------------------------------------------
\subsection{Monitoring and Observability}

% TEMPLATE: two by two block grid, the opener for the four child frames below.
% Four peers, none subordinate to another, so a levelled grid rather than a
% list, and each box is headed with its own child heading so the frame reads
% as the map of what follows. Each box takes one line on what that artifact
% tells you and how cheaply you get it.
\begin{frame}{ML-Specific Metrics}
  \begin{columns}[T,onlytextwidth]
    \begin{column}{0.47\textwidth}
      \begin{tdblockt}[equal height group=mlm]{Accuracy-related metrics}
      \end{tdblockt}
      \vskip0.5em
      \begin{tdblockt}[equal height group=mlm]{Predictions}
      \end{tdblockt}
    \end{column}
    \begin{column}{0.47\textwidth}
      \begin{tdblockt}[equal height group=mlm]{Features}
      \end{tdblockt}
      \vskip0.5em
      \begin{tdblockt}[equal height group=mlm]{Raw inputs}
      \end{tdblockt}
    \end{column}
  \end{columns}
  \slidesources{Huyen2022}
\end{frame}

% TEMPLATE: withmargin. Accuracy-related metrics are the most direct signal
% and the one that depends on a condition, user feedback arriving in
% production, so the main column carries the condition and what follows from
% it. Margin takes the book's own examples of that feedback.
\begin{frame}{Monitoring accuracy-related metrics}
  \begin{withmargin}
    \begin{tdmain}
    \end{tdmain}
    \begin{tdmargin}
    \end{tdmargin}
  \end{withmargin}
  \slidesources{Huyen2022}
\end{frame}

% TEMPLATE: single plain box plus a takeaway. Predictions are the one artifact
% low-dimensional enough to watch directly, which is one claim and not a
% comparison or a set, so it gets one untitled box and no margin. The takeaway
% keeps the book's own alert condition on a prediction distribution.
\begin{frame}{Monitoring predictions}
  \begin{tdblock}
  \end{tdblock}
  \begin{takeaway}
  \end{takeaway}
  \slidesources{Huyen2022}
\end{frame}

% TEMPLATE: withmargin with an uncaptioned code panel in the main column.
% Feature monitoring is schema and range checking, which the book shows as a
% validation snippet rather than a sentence, so the panel carries the checks.
% Margin takes the cost aside, what it means to check every feature of every
% request.
\begin{frame}{Monitoring features}
  \begin{withmargin}
    \begin{tdmain}
      \begin{tdcode}
      \end{tdcode}
    \end{tdmain}
    \begin{tdmargin}
    \end{tdmargin}
  \end{withmargin}
  \slidesources{Huyen2022}
\end{frame}

% TEMPLATE: withmargin. A short section carrying one claim, that raw inputs
% are the hardest artifact to reach because another team owns the data
% platform, so the main column takes the claim and the margin the aside on
% whose job this monitoring is.
\begin{frame}{Monitoring raw inputs}
  \begin{withmargin}
    \begin{tdmain}
    \end{tdmain}
    \begin{tdmargin}
    \end{tdmargin}
  \end{withmargin}
  \slidesources{Huyen2022}
\end{frame}

% TEMPLATE: three block row, the opener for the three child frames below. The
% toolbox heading names a set, so each tool takes one box headed with its own
% heading, levelled by the equal height group. Each box takes what that tool
% is for in the book's words, one line, with the detail left to its frame.
\begin{frame}{Monitoring Toolbox}
  \begin{columns}[T,onlytextwidth]
    \begin{column}{0.31\textwidth}
      \begin{tdblockt}[equal height group=tool]{Logs}
      \end{tdblockt}
    \end{column}
    \begin{column}{0.31\textwidth}
      \begin{tdblockt}[equal height group=tool]{Dashboards}
      \end{tdblockt}
    \end{column}
    \begin{column}{0.31\textwidth}
      \begin{tdblockt}[equal height group=tool]{Alerts}
      \end{tdblockt}
    \end{column}
  \end{columns}
  \slidesources{Huyen2022}
\end{frame}

% TEMPLATE: uncaptioned code panel, full width. A log is a record, so the
% slide shows one instead of describing one, and the frame title already
% names it. The panel takes a single log record with the metadata the book
% asks for; the distributed tracing point goes under it.
\begin{frame}{Logs}
  \begin{tdcode}
  \end{tdcode}
  \slidesources{Huyen2022}
\end{frame}

% TEMPLATE: withmargin. This one wanted a screenshot, but the chapter's two
% image frames are already spent on the feedback loop and the window plot, so
% the main column takes what a dashboard makes visible and to whom, and the
% margin the book's caveat about too many metrics on one board.
\begin{frame}{Dashboards}
  \begin{withmargin}
    \begin{tdmain}
    \end{tdmain}
    \begin{tdmargin}
    \end{tdmargin}
  \end{withmargin}
  \slidesources{Huyen2022}
\end{frame}

% TEMPLATE: three plain boxes in a row. The book gives an alert three parts of
% the same kind (the policy, the channels, the description), a countable set of
% peers, so one box each rather than three bullets. Headers stay empty until
% the next pass names each part after the book.
\begin{frame}{Alerts}
  \begin{columns}[T,onlytextwidth]
    \begin{column}{0.31\textwidth}
      \begin{tdblock}[equal height group=alert]
      \end{tdblock}
    \end{column}
    \begin{column}{0.31\textwidth}
      \begin{tdblock}[equal height group=alert]
      \end{tdblock}
    \end{column}
    \begin{column}{0.31\textwidth}
      \begin{tdblock}[equal height group=alert]
      \end{tdblock}
    \end{column}
  \end{columns}
  \slidesources{Huyen2022}
\end{frame}

% TEMPLATE: two block contrast plus a takeaway. The subchapter heading names
% both sides, so monitoring and observability take one levelled box each and
% the stronger notion reads off against the weaker instead of being asserted
% in prose. The takeaway is the book's one line on what observability buys
% that monitoring does not.
\begin{frame}{Observability}
  \begin{columns}[T,onlytextwidth]
    \begin{column}{0.47\textwidth}
      \begin{tdblockt}[equal height group=obs]{Monitoring}
      \end{tdblockt}
    \end{column}
    \begin{column}{0.47\textwidth}
      \begin{tdblockt}[equal height group=obs]{Observability}
      \end{tdblockt}
    \end{column}
  \end{columns}
  \begin{takeaway}
  \end{takeaway}
  \slidesources{Huyen2022}
\end{frame}

% ----------------------------------------------------------------------
\subsection{Summary}

% TEMPLATE: bare takeaway. A summary frame earns one line and nothing else, so
% no box, no margin and no grid; the takeaway holds the chapter's closing
% claim about shifts and monitoring in the book's own wording.
\begin{frame}{Summary}
  \begin{takeaway}
  \end{takeaway}
  \slidesources{Huyen2022}
\end{frame}

% ----------------------------------------------------------------------
% This chapter's own references. The refsection opened above scopes the
% citation numbering to this chapter, so chapter_pdfs/<this>.pdf resolves
% its own [n] markers instead of pointing at a list it does not contain.
\begin{frame}{References}
  \printbibliography[heading=none]
\end{frame}
\end{refsection}

%%% ch09_continual_learning.tex -- Intelligence Engineering, chapter 09
%% "Continual Learning and Test in Production".
%% Spine transferred from Huyen, Designing Machine Learning Systems, chapter
%% 9. The subchapters and the frame titles under them are the book's own
%% section and subsection headings, in the book's order.
%%
%% STATE: titles and TEMPLATE layout only. No body text. Every frame carries
%% the layout it is meant to be filled into, and a % TEMPLATE comment saying
%% why that layout and what belongs in each slot. Filling them is the next pass.
%%
%% Fragment of intelligence_engineering.tex. One chapter is one lecture and
%% gets exactly ONE \lecturepage, at the top. Inside it every \subsection
%% opens on the subchapter divider, which the master emits automatically via
%% \AtBeginSubsection; do not write those divider frames by hand.

\begin{refsection}

\lecturepage{IX}{Continual Learning and Test in Production}{}
\section[Continual Learning]{Continual Learning and Test in Production}

% ----------------------------------------------------------------------
\subsection{Continual Learning}

% TEMPLATE: two block contrast, levelled by the equal height group. The
% heading is a versus and the two regimes are defined against one another, so
% each takes one box rather than a shared bullet list. Each box holds how that
% regime initialises a training run and what it does with the previous model,
% in the book's words.
\begin{frame}{Stateless Retraining Versus Stateful Training}
  \begin{columns}[T,onlytextwidth]
    \begin{column}{0.47\textwidth}
      \begin{tdblockt}[equal height group=state]{Stateless retraining}
      \end{tdblockt}
    \end{column}
    \begin{column}{0.47\textwidth}
      \begin{tdblockt}[equal height group=state]{Stateful training}
      \end{tdblockt}
    \end{column}
  \end{columns}
  \slidesources{Huyen2022}
\end{frame}

% TEMPLATE: withmargin, the workhorse, and right for a heading that is a
% question the book answers with one argument. Main column takes the failure
% the question is really about, a model going stale against a shifting
% distribution. Margin takes the book's named cases, the ones where the
% distribution moves faster than any retraining schedule.
\begin{frame}{Why Continual Learning?}
  \begin{withmargin}
    \begin{tdmain}
    \end{tdmain}
    \begin{tdmargin}
    \end{tdmargin}
  \end{withmargin}
  \slidesources{Huyen2022}
\end{frame}

% TEMPLATE: three block row, 0.31 each, levelled by the equal height group.
% This frame opens the three deep headings under it, so it stays an overview:
% the three boxes are named by those children and each holds the challenge in
% one line. The detail stays on the three frames that follow.
\begin{frame}{Continual Learning Challenges}
  \begin{columns}[T,onlytextwidth]
    \begin{column}{0.31\textwidth}
      \begin{tdblockt}[equal height group=chal]{Fresh data access}
      \end{tdblockt}
    \end{column}
    \begin{column}{0.31\textwidth}
      \begin{tdblockt}[equal height group=chal]{Evaluation}
      \end{tdblockt}
    \end{column}
    \begin{column}{0.31\textwidth}
      \begin{tdblockt}[equal height group=chal]{Algorithm}
      \end{tdblockt}
    \end{column}
  \end{columns}
  \slidesources{Huyen2022}
\end{frame}

% TEMPLATE: side image hero, the first of this chapter's two. The book answers
% this heading with a drawing of where fresh data comes from, the warehouse
% route against the real time transport route, so the strip takes that figure
% and the sidebody takes the stations along it in the book's order. Caption
% names the figure; the gray placeholder stays until the asset is drawn.
\begin{frame}{Fresh data access challenge}
  \phimgside[0.33]{fresh_data_access_challenge.png}{Fresh data access}
  \begin{sidebody}
  \end{sidebody}
  \slidesources{Huyen2022}
\end{frame}

% TEMPLATE: one plain box plus a takeaway, and on purpose not a withmargin
% after the image frame. The heading names one challenge and the book argues
% it as one risk, so the box needs no header over it: it holds what a fresh
% model has to be tested against before it ever sees traffic. The takeaway
% carries the book's verdict on what a bad update costs.
\begin{frame}{Evaluation challenge}
  \begin{tdblock}
  \end{tdblock}
  \begin{takeaway}
  \end{takeaway}
  \slidesources{Huyen2022}
\end{frame}

% TEMPLATE: two untitled boxes, levelled by the equal height group. The book
% splits this challenge by model family, the ones that need the full dataset
% to be updated at all and the ones that can be updated from a data batch, so
% each family takes a box instead of two bullets in one. No box headers here:
% the heading names only the challenge, the families arrive with the text.
\begin{frame}{Algorithm challenge}
  \begin{columns}[T,onlytextwidth]
    \begin{column}{0.47\textwidth}
      \begin{tdblock}[equal height group=algo]
      \end{tdblock}
    \end{column}
    \begin{column}{0.47\textwidth}
      \begin{tdblock}[equal height group=algo]
      \end{tdblock}
    \end{column}
  \end{columns}
  \slidesources{Huyen2022}
\end{frame}

% TEMPLATE: two by two block grid, sized to the count the heading names. Four
% stages, each named by one of the deep headings under this frame, none
% subordinate to another, so one box each rather than a numbered list. This is
% the opener, so each box keeps to what the stage is; the four frames that
% follow carry the detail. Two equal height groups keep the rows level.
\begin{frame}{Four Stages of Continual Learning}
  \begin{columns}[T,onlytextwidth]
    \begin{column}{0.47\textwidth}
      \begin{tdblockt}[equal height group=stg1]{Stage 1}
      \end{tdblockt}
    \end{column}
    \begin{column}{0.47\textwidth}
      \begin{tdblockt}[equal height group=stg1]{Stage 2}
      \end{tdblockt}
    \end{column}
  \end{columns}
  \vskip0.5em
  \begin{columns}[T,onlytextwidth]
    \begin{column}{0.47\textwidth}
      \begin{tdblockt}[equal height group=stg2]{Stage 3}
      \end{tdblockt}
    \end{column}
    \begin{column}{0.47\textwidth}
      \begin{tdblockt}[equal height group=stg2]{Stage 4}
      \end{tdblockt}
    \end{column}
  \end{columns}
  \slidesources{Huyen2022}
\end{frame}

% TEMPLATE: withmargin. The stage is one state of practice, so the main column
% runs it as one claim: what triggers the retraining and who does each step by
% hand. Margin takes the book's reading of why a team sits here, the cadence
% it settles into and what it is not yet paying for.
\begin{frame}{Stage 1: Manual, stateless retraining}
  \begin{withmargin}
    \begin{tdmain}
    \end{tdmain}
    \begin{tdmargin}
    \end{tdmargin}
  \end{withmargin}
  \slidesources{Huyen2022}
\end{frame}

% TEMPLATE: captioned code panel plus a takeaway. This stage is reached by
% writing a script and a schedule, so the honest slot is the panel: it takes
% the steps the script has to chain and what the scheduler hands it, written
% out rather than described. The takeaway keeps the book's condition on the
% retraining frequency the schedule is allowed to claim.
\begin{frame}{Stage 2: Automated retraining}
  \begin{tdcodet}{Automated retraining}
  \end{tdcodet}
  \begin{takeaway}
  \end{takeaway}
  \slidesources{Huyen2022}
\end{frame}

% TEMPLATE: withmargin. The step from stage 2 to stage 3 is a reconfiguration
% rather than a new set of parts, so the main column takes what changes in the
% script and what the system now has to track per run. Margin takes the book's
% account of what stateful training buys, stated as the resource it stops
% spending.
\begin{frame}{Stage 3: Automated, stateful training}
  \begin{withmargin}
    \begin{tdmain}
    \end{tdmain}
    \begin{tdmargin}
    \end{tdmargin}
  \end{withmargin}
  \slidesources{Huyen2022}
\end{frame}

% TEMPLATE: tddef, kept bare. The last stage is where the chapter's title term
% finally earns its definition, so the term box is the whole slide: it takes
% what replaces the schedule as the trigger for an update, and the condition
% the book attaches to running that way.
\begin{frame}{Stage 4: Continual learning}
  \begin{tddef}{Continual learning}
  \end{tddef}
  \slidesources{Huyen2022}
\end{frame}

% TEMPLATE: withmargin, the overview form, since this frame opens the two deep
% headings under it and must not preempt them. Main column reframes the
% heading's question into the one the book says is answerable. Margin holds the
% aside that the answer is not a schedule but a measurement.
\begin{frame}{How Often to Update Your Models}
  \begin{withmargin}
    \begin{tdmain}
    \end{tdmain}
    \begin{tdmargin}
    \end{tdmargin}
  \end{withmargin}
  \slidesources{Huyen2022}
\end{frame}

% TEMPLATE: booktabs table, because the book settles this heading with an
% experiment and its numbers only make their point ranked against one another.
% Rows are the book's freshness buckets, the age of the training data; the
% second column takes the measured performance and gets its own header with
% the metric in the next pass. booktabs only, no vertical rules.
\begin{frame}{Value of data freshness}
  \footnotesize
  \renewcommand{\arraystretch}{1.35}
  \begin{tabular}{@{}%
    >{\raggedright\arraybackslash}p{0.34\textwidth}%
    >{\raggedright\arraybackslash}p{0.34\textwidth}@{}}
    \toprule
    \textbf{Data freshness} & \\
    \midrule
     & \\
     & \\
     & \\
    \bottomrule
  \end{tabular}
  \slidesources{Huyen2022}
\end{frame}

% TEMPLATE: two block contrast, levelled by the equal height group. The heading
% is a versus, so each kind of iteration takes one box rather than a bullet
% list: one box for what is changed in the model, one for what is changed in
% the data. Each holds the cost the book attaches to that move and the signal
% that says which one to spend on next.
\begin{frame}{Model iteration versus data iteration}
  \begin{columns}[T,onlytextwidth]
    \begin{column}{0.47\textwidth}
      \begin{tdblockt}[equal height group=iter]{Model iteration}
      \end{tdblockt}
    \end{column}
    \begin{column}{0.47\textwidth}
      \begin{tdblockt}[equal height group=iter]{Data iteration}
      \end{tdblockt}
    \end{column}
  \end{columns}
  \slidesources{Huyen2022}
\end{frame}

% ----------------------------------------------------------------------
\subsection{Test in Production}

% TEMPLATE: side image hero, the second and last of this chapter's two. The
% book makes this technique concrete with a drawing of traffic forked to two
% models with only one answering, so the strip takes that figure and the
% sidebody takes how the shadow traffic is logged and compared. Caption names
% the figure; placeholder stays until the asset exists.
\begin{frame}{Shadow Deployment}
  \phimgside[0.33]{shadow_deployment.png}{Shadow deployment}
  \begin{sidebody}
  \end{sidebody}
  \slidesources{Huyen2022}
\end{frame}

% TEMPLATE: withmargin. The technique is one procedure, so the main column runs
% it in the book's order: how traffic is split, what is measured, and when the
% result may be read. Margin takes the two conditions the book puts on the
% split, which are the part teams get wrong and read better stacked.
\begin{frame}{A/B Testing}
  \begin{withmargin}
    \begin{tdmain}
    \end{tdmain}
    \begin{tdmargin}
    \end{tdmargin}
  \end{withmargin}
  \slidesources{Huyen2022}
\end{frame}

% TEMPLATE: tddef, the form this deck uses for a named technique the book
% introduces by definition rather than by comparison. The term box is the whole
% slide: it takes the book's definition, the way traffic is shifted over, and
% the abort condition that makes it a release strategy and not a test.
\begin{frame}{Canary Release}
  \begin{tddef}{Canary release}
  \end{tddef}
  \slidesources{Huyen2022}
\end{frame}

% TEMPLATE: withmargin. The book argues this one against the split it replaces,
% so the main column takes what is interleaved and how a preference is read off
% a single user's session. Margin takes the position bias the method has to
% correct for and the sample size the book says it saves.
\begin{frame}{Interleaving Experiments}
  \begin{withmargin}
    \begin{tdmain}
    \end{tdmain}
    \begin{tdmargin}
    \end{tdmargin}
  \end{withmargin}
  \slidesources{Huyen2022}
\end{frame}

% TEMPLATE: three untitled boxes, 0.31 each, levelled by the equal height
% group. The book gates this technique on a short list of requirements the
% system has to meet before a bandit can be run at all, so one box per
% requirement rather than three bullets. No headers: the heading names only the
% method. This is the opener for the deep heading under it, so exploration
% against exploitation stays here and the contextual case follows.
\begin{frame}{Bandits}
  \begin{columns}[T,onlytextwidth]
    \begin{column}{0.31\textwidth}
      \begin{tdblock}[equal height group=band]
      \end{tdblock}
    \end{column}
    \begin{column}{0.31\textwidth}
      \begin{tdblock}[equal height group=band]
      \end{tdblock}
    \end{column}
    \begin{column}{0.31\textwidth}
      \begin{tdblock}[equal height group=band]
      \end{tdblock}
    \end{column}
  \end{columns}
  \slidesources{Huyen2022}
\end{frame}

% TEMPLATE: tddef. The heading names a term and then says what it is for, which
% is exactly a definition plus its use, so the term box is the slide's one
% slot. It takes what the context adds to a plain bandit and what is being
% explored once the context is in, with the book's data efficiency claim.
\begin{frame}{Contextual bandits as an exploration strategy}
  \begin{tddef}{Contextual bandits}
  \end{tddef}
  \slidesources{Huyen2022}
\end{frame}

% ----------------------------------------------------------------------
\subsection{Summary}

% TEMPLATE: statement, kept bare. A summary is a closing line, not a slide of
% bullets, so it takes the loud centred form and nothing else. The main line is
% the book's own summary of why a deployed model is never finished; the
% substatement is the consequence it carries, that evaluation moves into
% production.
\begin{frame}{Summary}
  \begin{statement}
    \substatement{}
  \end{statement}
  \slidesources{Huyen2022}
\end{frame}

% ----------------------------------------------------------------------
% This chapter's own references. The refsection opened above scopes the
% citation numbering to this chapter, so chapter_pdfs/<this>.pdf resolves
% its own [n] markers instead of pointing at a list it does not contain.
\begin{frame}{References}
  \printbibliography[heading=none]
\end{frame}
\end{refsection}

%%% ch10_infrastructure.tex -- Intelligence Engineering, chapter 10
%% "Infrastructure and Tooling for MLOps".
%% Spine transferred from Huyen, Designing Machine Learning Systems, chapter
%% 10. The subchapters and the frame titles under them are the book's own
%% section and subsection headings, in the book's order.
%%
%% STATE: titles and TEMPLATE layout only. No body text. Every frame carries
%% the layout it is meant to be filled into, and a % TEMPLATE comment saying
%% why that layout and what belongs in each slot. Filling them is the next pass.
%%
%% Fragment of intelligence_engineering.tex. One chapter is one lecture and
%% gets exactly ONE \lecturepage, at the top. Inside it every \subsection
%% opens on the subchapter divider, which the master emits automatically via
%% \AtBeginSubsection; do not write those divider frames by hand.

\begin{refsection}

\lecturepage{X}{Infrastructure and Tooling for MLOps}{}
\section[Infrastructure]{Infrastructure and Tooling for MLOps}

% ----------------------------------------------------------------------
\subsection{Storage and Compute}

% TEMPLATE: two block contrast, levelled by the equal height group. The heading
% is a versus, so each side takes one box rather than a bullet list and the box
% headers are read straight off the two sides the heading names. Public cloud
% box takes elasticity and the start without capital outlay; private box takes
% what owning the machines buys back, with the book's cloud bill numbers, the
% revenue share and the repatriation saving, kept as numbers.
\begin{frame}{Public Cloud Versus Private Data Centers}
  \begin{columns}[T,onlytextwidth]
    \begin{column}{0.47\textwidth}
      \begin{tdblockt}[equal height group=cloud]{Public cloud}
      \end{tdblockt}
    \end{column}
    \begin{column}{0.47\textwidth}
      \begin{tdblockt}[equal height group=cloud]{Private data centers}
      \end{tdblockt}
    \end{column}
  \end{columns}
  \slidesources{Huyen2022}
\end{frame}

% ----------------------------------------------------------------------
\subsection{Development Environment}

% TEMPLATE: captioned code panel, kept bare. The book's dev environment is a
% list of tools and the configuration that pins them, the IDE, versioning and
% CI/CD, so the honest slot is a panel and not prose about tooling. Caption is
% the heading's own subject; the panel takes the named tools and the versioning
% setup in the book's order.
\begin{frame}{Dev Environment Setup}
  \begin{tdcodet}{Dev environment}
  \end{tdcodet}
  \slidesources{Huyen2022}
\end{frame}

% TEMPLATE: withmargin, the chapter's workhorse and the right shape here: one
% claim plus one aside. Main column takes the book's recommendation, standardize
% the environment company wide and move development to the cloud, and what that
% move costs. Margin takes the resistance the book records and the remote IDE it
% names as the way around it.
\begin{frame}{Standardizing Dev Environments}
  \begin{withmargin}
    \begin{tdmain}
    \end{tdmain}
    \begin{tdmargin}
    \end{tdmargin}
  \end{withmargin}
  \slidesources{Huyen2022}
\end{frame}

% TEMPLATE: withmargin with a captioned code panel in the main column. The
% heading is about containers, which are configuration before they are an idea,
% so the panel takes the Dockerfile, the image against container distinction and
% the orchestration tools in the book's words. Margin takes why the same
% environment has to travel from dev to prod at all.
\begin{frame}{From Dev to Prod: Containers}
  \begin{withmargin}
    \begin{tdmain}
      \begin{tdcodet}{Containers}
      \end{tdcodet}
    \end{tdmain}
    \begin{tdmargin}
    \end{tdmargin}
  \end{withmargin}
  \slidesources{Huyen2022}
\end{frame}

% ----------------------------------------------------------------------
\subsection{Resource Management}

% TEMPLATE: three block row, headers taken from the three peers the heading
% names, levelled by the equal height group. The book's point is that these
% three are not the same thing, which a row of boxes shows and a bullet list
% hides. Each box takes what that tool schedules and what it cannot do, with
% the book's division of labour, jobs against resources, landing in the last.
\begin{frame}{Cron, Schedulers, and Orchestrators}
  \begin{columns}[T,onlytextwidth]
    \begin{column}{0.31\textwidth}
      \begin{tdblockt}[equal height group=rm]{Cron}
      \end{tdblockt}
    \end{column}
    \begin{column}{0.31\textwidth}
      \begin{tdblockt}[equal height group=rm]{Schedulers}
      \end{tdblockt}
    \end{column}
    \begin{column}{0.31\textwidth}
      \begin{tdblockt}[equal height group=rm]{Orchestrators}
      \end{tdblockt}
    \end{column}
  \end{columns}
  \slidesources{Huyen2022}
\end{frame}

% TEMPLATE: side image hero, the chapter's one and only figure. A workflow is a
% graph of steps and the book draws it as one, so the strip takes that figure
% and the sidebody takes the book's comparison of the named workflow managers,
% what each of them defines a step in and what it does about containers. The
% caption names the figure; the gray placeholder stays until the asset is drawn.
\begin{frame}{Data Science Workflow Management}
  \phimgside[0.33]{data_science_workflow_management.png}{Data science workflow}
  \begin{sidebody}
  \end{sidebody}
  \slidesources{Huyen2022}
\end{frame}

% ----------------------------------------------------------------------
\subsection{ML Platform}

% TEMPLATE: withmargin. The heading is one platform component and not a set of
% peers, so the main column runs it as a single claim: what the deployment
% component owes the rest of the platform once a model is pushed, online and
% batch prediction behind one service. Margin takes the named deployment tools
% the book lists and the caveat it attaches to buying them.
\begin{frame}{Model Deployment}
  \begin{withmargin}
    \begin{tdmain}
    \end{tdmain}
    \begin{tdmargin}
    \end{tdmargin}
  \end{withmargin}
  \slidesources{Huyen2022}
\end{frame}

% TEMPLATE: tddef, kept bare. Model store is a term the book has to define
% before it can show that most teams implement only a fraction of it, so the
% term box is the whole slide rather than a claim with an aside. The box takes
% the definition and then the artifacts the book says belong in a model store,
% which is the part the tools skip.
\begin{frame}{Model Store}
  \begin{tddef}{Model store}
  \end{tddef}
  \slidesources{Huyen2022}
\end{frame}

% TEMPLATE: three block row, untitled on purpose. The book settles this heading
% by naming three problems a feature store solves, so three levelled boxes
% rather than three bullets in one; the headers stay off because the heading
% names only the store itself and the three names arrive with the text in the
% next pass. Each box takes one problem and the tools the book credits with it.
\begin{frame}{Feature Store}
  \begin{columns}[T,onlytextwidth]
    \begin{column}{0.31\textwidth}
      \begin{tdblock}[equal height group=fs]
      \end{tdblock}
    \end{column}
    \begin{column}{0.31\textwidth}
      \begin{tdblock}[equal height group=fs]
      \end{tdblock}
    \end{column}
    \begin{column}{0.31\textwidth}
      \begin{tdblock}[equal height group=fs]
      \end{tdblock}
    \end{column}
  \end{columns}
  \slidesources{Huyen2022}
\end{frame}

% ----------------------------------------------------------------------
\subsection{Build Versus Buy}

% TEMPLATE: booktabs comparison table. This is the chapter's second versus and
% the two box contrast is already spent on the cloud frame, so the decision goes
% into a table: two columns headed by the two sides the heading names, the book's
% deciding questions as the rows. Row labels and cells are filled together in
% the next pass, which is why both are empty here.
\begin{frame}{Build Versus Buy}
  \footnotesize
  \renewcommand{\arraystretch}{1.35}
  \begin{tabular}{@{}%
    >{\raggedright\arraybackslash}p{0.22\textwidth}%
    >{\raggedright\arraybackslash}p{0.35\textwidth}%
    >{\raggedright\arraybackslash}p{0.35\textwidth}@{}}
    \toprule
     & \textbf{Build} & \textbf{Buy}\\
    \midrule
     & & \\
     & & \\
     & & \\
    \bottomrule
  \end{tabular}
  \slidesources{Huyen2022}
\end{frame}

% ----------------------------------------------------------------------
\subsection{Summary}

% TEMPLATE: takeaway, kept bare. A summary is a closing line and not a slide of
% bullets, so it takes the big centred form and nothing else; the chapter has no
% other takeaway, so this one lands. The line is the book's own closing claim,
% that what a team needs from infrastructure is set by the number and the scale
% of its use cases.
\begin{frame}{Summary}
  \begin{takeaway}
  \end{takeaway}
  \slidesources{Huyen2022}
\end{frame}

% ----------------------------------------------------------------------
% This chapter's own references. The refsection opened above scopes the
% citation numbering to this chapter, so chapter_pdfs/<this>.pdf resolves
% its own [n] markers instead of pointing at a list it does not contain.
\begin{frame}{References}
  \printbibliography[heading=none]
\end{frame}
\end{refsection}

%%% ch11_human_side.tex -- Intelligence Engineering, chapter 11 "The Human Side
%% of Machine Learning".
%% Spine transferred from Huyen, Designing Machine Learning Systems, chapter
%% 11. The subchapters and the frame titles under them are the book's own
%% section and subsection headings, in the book's order.
%%
%% STATE: titles and TEMPLATE layout only. No body text. Every frame carries
%% the layout it is meant to be filled into, and a % TEMPLATE comment saying
%% why that layout and what belongs in each slot. Filling them is the next pass.
%%
%% Fragment of intelligence_engineering.tex. One chapter is one lecture and
%% gets exactly ONE \lecturepage, at the top. Inside it every \subsection
%% opens on the subchapter divider, which the master emits automatically via
%% \AtBeginSubsection; do not write those divider frames by hand.

\begin{refsection}

\lecturepage{XI}{The Human Side of Machine Learning}{}
\section[Human Side]{The Human Side of Machine Learning}

% ----------------------------------------------------------------------
\subsection{User Experience}

% TEMPLATE: withmargin, the chapter's workhorse and the right opener for a
% requirement: one claim plus the aside that makes it concrete. Main column
% takes what consistency means to a user and the tension the book sets it
% against. Margin holds the book's worked case, the travel site whose model
% overrode the filters a user had set, and the verdict its designers reached.
\begin{frame}{Ensuring User Experience Consistency}
  \begin{withmargin}
    \begin{tdmain}
    \end{tdmain}
    \begin{tdmargin}
    \end{tdmargin}
  \end{withmargin}
  \slidesources{Huyen2022}
\end{frame}

% TEMPLATE: side image hero, the first of this chapter's two. The book makes
% this heading concrete with a figure of several predictions offered at once,
% so the strip takes that figure and the sidebody takes why a mostly correct
% prediction is worthless to a user who cannot check it, and what showing many
% predictions buys instead. The gray placeholder stays until the asset exists.
\begin{frame}{Combatting ``Mostly Correct'' Predictions}
  \phimgside[0.33]{combatting_mostly_correct_predictions.png}{Mostly correct predictions}
  \begin{sidebody}
  \end{sidebody}
  \slidesources{Huyen2022}
\end{frame}

% TEMPLATE: tddef plus a takeaway, and deliberately not a third layout of the
% same kind in a row. The heading is the book's own term for a backup route,
% so the term box is the slide: it takes the definition and the rule that
% decides when the fallback fires. The takeaway keeps the speed against
% accuracy tradeoff the book closes the heading with.
\begin{frame}{Smooth Failing}
  \begin{tddef}{Smooth failing}
  \end{tddef}
  \begin{takeaway}
  \end{takeaway}
  \slidesources{Huyen2022}
\end{frame}

% ----------------------------------------------------------------------
\subsection{Team Structure}

% TEMPLATE: withmargin. The heading is about who has to be in the room, so the
% main column takes the subject matter experts the book names and the stages it
% says they must reach beyond labeling. Margin holds the friction the book
% reports once those experts are asked to work inside engineers' tooling.
\begin{frame}{Cross-functional Teams Collaboration}
  \begin{withmargin}
    \begin{tdmain}
    \end{tdmain}
    \begin{tdmargin}
    \end{tdmargin}
  \end{withmargin}
  \slidesources{Huyen2022}
\end{frame}

% TEMPLATE: two block contrast, levelled by the equal height group. This frame
% opens the two deep headings under it, so it stays a comparison and not a
% detail layout. The headers are the two approaches the children name; each box
% takes that approach in a line, and the children carry the workflow itself.
\begin{frame}{End-to-End Data Scientists}
  \begin{columns}[T,onlytextwidth]
    \begin{column}{0.47\textwidth}
      \begin{tdblockt}[equal height group=e2e]{Approach 1}
      \end{tdblockt}
    \end{column}
    \begin{column}{0.47\textwidth}
      \begin{tdblockt}[equal height group=e2e]{Approach 2}
      \end{tdblockt}
    \end{column}
  \end{columns}
  \slidesources{Huyen2022}
\end{frame}

% TEMPLATE: side image hero, the second and last of this chapter's two. The
% book draws this approach as a handoff, the data scientists on one side and
% the team that owns production on the other, so the strip takes that diagram
% and the sidebody takes the communication overhead the book charges it with.
\begin{frame}{Approach 1: Have a separate team to manage production}
  \phimgside[0.33]{approach_1_have_a_separate_team_to_manage_production.png}{A separate team to manage production}
  \begin{sidebody}
  \end{sidebody}
  \slidesources{Huyen2022}
\end{frame}

% TEMPLATE: withmargin. The heading is one claim about ownership, so the main
% column takes what owning the entire process means step by step and the
% infrastructure it presupposes. Margin holds the named practice the book cites
% as its case and the objection it quotes from the data scientists themselves.
\begin{frame}{Approach 2: Data scientists own the entire process}
  \begin{withmargin}
    \begin{tdmain}
    \end{tdmain}
    \begin{tdmargin}
    \end{tdmargin}
  \end{withmargin}
  \slidesources{Huyen2022}
\end{frame}

% ----------------------------------------------------------------------
\subsection{Responsible AI}

% TEMPLATE: booktabs comparison table, and the opener for the two case studies
% under it. One column per case study named in the children, so the slide reads
% as a pair rather than as two sentences. Row labels are left empty on purpose:
% they are the dimensions the next pass fills, what the system did and who
% carried the cost.
\begin{frame}{Irresponsible AI: Case Studies}
  \footnotesize
  \renewcommand{\arraystretch}{1.35}
  \begin{tabular}{@{}%
    >{\raggedright\arraybackslash}p{0.22\textwidth}%
    >{\raggedright\arraybackslash}p{0.35\textwidth}%
    >{\raggedright\arraybackslash}p{0.35\textwidth}@{}}
    \toprule
     & \textbf{Case study I} & \textbf{Case study II}\\
    \midrule
     & & \\
     & & \\
     & & \\
     & & \\
    \bottomrule
  \end{tabular}
  \slidesources{Huyen2022}
\end{frame}

% TEMPLATE: withmargin. The case is an argument from numbers, so the main
% column runs what the grader was asked to do and what it graded on instead.
% Margin takes the book's figures on where the downgrades landed, which is the
% evidence the whole case rests on.
\begin{frame}{Case study I: Automated grader's biases}
  \begin{withmargin}
    \begin{tdmain}
    \end{tdmain}
    \begin{tdmargin}
    \end{tdmargin}
  \end{withmargin}
  \slidesources{Huyen2022}
\end{frame}

% TEMPLATE: one plain box plus a takeaway, and not a second withmargin in a
% row. The heading already carries the claim, so the box needs no header over
% it: it takes the released dataset, the fields left inside it and the
% reidentification that followed. The takeaway is the book's verdict on what
% the word anonymized actually buys.
\begin{frame}{Case study II: The danger of ``anonymized'' data}
  \begin{tdblock}
  \end{tdblock}
  \begin{takeaway}
  \end{takeaway}
  \slidesources{Huyen2022}
\end{frame}

% TEMPLATE: withmargin, opener for the seven deep headings under it. Seven
% peers do not grid, so this frame stays one claim about why a framework is
% needed at all and the children carry one recommendation each. Margin holds
% the book's own caveat, that this is a starting point and not a standard.
\begin{frame}{A Framework for Responsible AI}
  \begin{withmargin}
    \begin{tdmain}
    \end{tdmain}
    \begin{tdmargin}
    \end{tdmargin}
  \end{withmargin}
  \slidesources{Huyen2022}
\end{frame}

% TEMPLATE: five block grid, three across then two, levelled by two equal
% height groups. The book answers this heading with peer sources of bias spread
% along the pipeline, none subordinate to another, so one box each rather than
% a bullet list. Boxes stay untitled: the heading names no source, the names
% arrive with the text. Trim or add a box to the count the book gives.
\begin{frame}{Discover sources for model biases}
  \begin{columns}[T,onlytextwidth]
    \begin{column}{0.31\textwidth}
      \begin{tdblock}[equal height group=bias1]
      \end{tdblock}
    \end{column}
    \begin{column}{0.31\textwidth}
      \begin{tdblock}[equal height group=bias1]
      \end{tdblock}
    \end{column}
    \begin{column}{0.31\textwidth}
      \begin{tdblock}[equal height group=bias1]
      \end{tdblock}
    \end{column}
  \end{columns}
  \vskip0.4em
  \begin{columns}[T,onlytextwidth]
    \begin{column}{0.47\textwidth}
      \begin{tdblock}[equal height group=bias2]
      \end{tdblock}
    \end{column}
    \begin{column}{0.47\textwidth}
      \begin{tdblock}[equal height group=bias2]
      \end{tdblock}
    \end{column}
  \end{columns}
  \slidesources{Huyen2022}
\end{frame}

% TEMPLATE: withmargin. The heading names a limitation and not a set of peers,
% so the main column takes what the data on its own cannot settle and what the
% book asks for in its place. Margin holds the case it argues from and the
% people it says the data never reached.
\begin{frame}{Understand the limitations of the data-driven approach}
  \begin{withmargin}
    \begin{tdmain}
    \end{tdmain}
    \begin{tdmargin}
    \end{tdmargin}
  \end{withmargin}
  \slidesources{Huyen2022}
\end{frame}

% TEMPLATE: booktabs table, both column headers taken from the heading's own
% words. A tradeoff is a pair, so one row per desideratum against what it
% costs reads as a table and not as prose. Rows stay empty for the pairs the
% book names and the argument it makes for each.
\begin{frame}{Understand the trade-offs between different desiderata}
  \footnotesize
  \renewcommand{\arraystretch}{1.35}
  \begin{tabular}{@{}%
    >{\raggedright\arraybackslash}p{0.45\textwidth}%
    >{\raggedright\arraybackslash}p{0.45\textwidth}@{}}
    \toprule
    \textbf{Desideratum} & \textbf{Trade-off}\\
    \midrule
     & \\
     & \\
     & \\
    \bottomrule
  \end{tabular}
  \slidesources{Huyen2022}
\end{frame}

% TEMPLATE: one titled box plus a takeaway. The heading is a short imperative,
% so it becomes the box header itself and the box takes the reason behind it,
% what a bias costs once everything else is already built on top of it. The
% takeaway carries the book's line on the cheapest moment to act.
\begin{frame}{Act early}
  \begin{tdblockt}{Act early}
  \end{tdblockt}
  \begin{takeaway}
  \end{takeaway}
  \slidesources{Huyen2022}
\end{frame}

% TEMPLATE: captioned code panel. A model card is a document with fixed
% fields, so the honest slot is a panel showing those fields rather than a
% sentence about them. The caption is the artifact the heading names; the panel
% takes the book's field list and the question each field is there to answer.
\begin{frame}{Create model cards}
  \begin{tdcodet}{Model card}
  \end{tdcodet}
  \slidesources{Huyen2022}
\end{frame}

% TEMPLATE: withmargin. The heading asks for a process and not a list of peers,
% so the main column takes what that process has to cover and the failure mode
% the book names, fixes bolted on after the fact. Margin holds the tooling and
% the audits the book points at.
\begin{frame}{Establish processes for mitigating biases}
  \begin{withmargin}
    \begin{tdmain}
    \end{tdmain}
    \begin{tdmargin}
    \end{tdmargin}
  \end{withmargin}
  \slidesources{Huyen2022}
\end{frame}

% TEMPLATE: three block row, levelled by the equal height group, and a break
% from the column layout before it. The book answers this heading with named
% places to follow the field, so one box per kind of source rather than a run
% of names inside a sentence. Boxes stay untitled until those names are in.
\begin{frame}{Stay up-to-date on responsible AI}
  \begin{columns}[T,onlytextwidth]
    \begin{column}{0.31\textwidth}
      \begin{tdblock}[equal height group=rai]
      \end{tdblock}
    \end{column}
    \begin{column}{0.31\textwidth}
      \begin{tdblock}[equal height group=rai]
      \end{tdblock}
    \end{column}
    \begin{column}{0.31\textwidth}
      \begin{tdblock}[equal height group=rai]
      \end{tdblock}
    \end{column}
  \end{columns}
  \slidesources{Huyen2022}
\end{frame}

% ----------------------------------------------------------------------
\subsection{Summary}

% TEMPLATE: statement, kept bare. A summary is a closing line and not a slide
% of bullets, so it takes the loud centred form and nothing else. The main line
% is the book's own summing up of the human side; the substatement is the
% consequence it carries out of the chapter.
\begin{frame}{Summary}
  \begin{statement}
    \substatement{}
  \end{statement}
  \slidesources{Huyen2022}
\end{frame}

% ----------------------------------------------------------------------
% This chapter's own references. The refsection opened above scopes the
% citation numbering to this chapter, so chapter_pdfs/<this>.pdf resolves
% its own [n] markers instead of pointing at a list it does not contain.
\begin{frame}{References}
  \printbibliography[heading=none]
\end{frame}
\end{refsection}


% ---------------------------------------------------------- bibliography
% There is none at deck level, on purpose. Every chapter wraps itself in a
% \begin{refsection} and closes on its own References frame, which is what
% lets a standalone chapter PDF resolve its own [n] markers. A closing frame
% here would reprint the same bib once more, after twelve chapters have each
% already printed the part of it they use. If you ever want one anyway, it
% needs \nocite{*}: refsection 0, the material outside every chapter, cites
% nothing, so a bare \printbibliography typesets an empty frame.

% ===================================================================== close

%\begin{frame}{What you learned}
%\end{frame}

% The closing statement slot.
\begin{frame}[plain]
  This is the last slide.\\
  Look how far you've come.

\end{frame}

\end{document}




\iffalse





\fi